\documentclass[aps,preprint]{revtex4}%
\usepackage{amssymb}
\usepackage{amsfonts}
\usepackage{amsmath}
\usepackage{graphicx}%
\setcounter{MaxMatrixCols}{30}
%TCIDATA{OutputFilter=latex2.dll}
%TCIDATA{Version=5.50.0.2960}
%TCIDATA{CSTFile=revtex4.cst}
%TCIDATA{Created=Saturday, June 11, 2011 09:15:54}
%TCIDATA{LastRevised=Tuesday, August 09, 2011 10:37:48}
%TCIDATA{<META NAME="GraphicsSave" CONTENT="32">}
%TCIDATA{<META NAME="SaveForMode" CONTENT="1">}
%TCIDATA{BibliographyScheme=Manual}
%TCIDATA{<META NAME="DocumentShell" CONTENT="Articles\SW\REVTeX 4">}
%BeginMSIPreambleData
\providecommand{\U}[1]{\protect\rule{.1in}{.1in}}
%EndMSIPreambleData
\newtheorem{theorem}{Theorem}
\newtheorem{acknowledgement}[theorem]{Acknowledgement}
\newtheorem{algorithm}[theorem]{Algorithm}
\newtheorem{axiom}[theorem]{Axiom}
\newtheorem{claim}[theorem]{Claim}
\newtheorem{conclusion}[theorem]{Conclusion}
\newtheorem{condition}[theorem]{Condition}
\newtheorem{conjecture}[theorem]{Conjecture}
\newtheorem{corollary}[theorem]{Corollary}
\newtheorem{criterion}[theorem]{Criterion}
\newtheorem{definition}[theorem]{Definition}
\newtheorem{example}[theorem]{Example}
\newtheorem{exercise}[theorem]{Exercise}
\newtheorem{lemma}[theorem]{Lemma}
\newtheorem{notation}[theorem]{Notation}
\newtheorem{problem}[theorem]{Problem}
\newtheorem{proposition}[theorem]{Proposition}
\newtheorem{remark}[theorem]{Remark}
\newtheorem{solution}[theorem]{Solution}
\newtheorem{summary}[theorem]{Summary}
\newenvironment{proof}[1][Proof]{\noindent\textbf{#1.} }{\ \rule{0.5em}{0.5em}}
\begin{document}
\preprint{ }
\title[N-Representable Approximations]{Systematic N-Representable Approximations using Thouless Coherent States}
\author{Brian Weiner}
\affiliation{Department of Physics, Pennsylvania State University, DuBois PA 15801}
\keywords{N-Representability, Fock Space, Positive Extensions, Approximations, Coherent
States, Group Manifolds}
\pacs{31.15.A-,31.15.-p}

\begin{abstract}
The N-Representability Problem entails characterizing the set of second order
reduced states that are contractions of N-electron states of the Fermion Fock
algebra. This problem is formulated in the form of finding the conditions that
a positive linear functional defined on a subspace of this algebra must
satisfy in order to be extended to the whole algebra. As this algebra is a
w*-algebra one can utilize a theorem by Kadison that shows it is sufficient to
consider the values of linear functionals on projectors contained in the
subspace in order to determine whether they have positive extensions. Thus we
find the form of projectors belonging to the subspace of one and two particle
operators and subsequently show that the extension conditions needed in the
N-Representability Problem correspond to generalized P, Q and G conditions
plus the additional constraints that the functionals be dispersion free on the
number operator and their values on one particle operators determined by their
values on two particle operators.

It is shown that this set of projectors form a homogeneous space under the
action of the unitary group of one electron Hilbert space. Using the Thouless
parameterization of this group generated by two electron Thouless Coherent
States the preceeding N-Representability conditions are expressed in terms of
conditions on Thouless Q symbols (also known as covariant symbols) of positive
operators acting in two electron Hilbert space. It is observed that the
Q-symbols form a Banach space of bounded real vlaued functions defined on
complex euclidean vectors. As the N-electron energy can be expressed as a
linear functional of Q-symbols that are constrained to N-Representable one can
determine the ground state of an N-electron system by performing a constrained
linear variation over functions from this Banach space. N-Representable
approximations can then be generated by utilizing various approximation
schemes for these functions.

\end{abstract}
\date{07/13/2011}
\startpage{0}
\maketitle
\tableofcontents


\section{Introduction}

In order to develop a sufficiently accurate quantum mechanical description of
the electronic portion of atomic and molecular systems containing
$N$-electrons it has been found that one only needs to consider a subset,
$\mathfrak{H}_{N2},$ of Hamiltonian operators that describe the individual
kinetic energies, interactions with external and internal electric and
magnetic fields and only the \emph{simultaneous pair wise }interactions of the
electrons. This subset acts in the $N$-fold antisymmetric tensor product
Hilbert space $\mathcal{H}^{N}$ and correspond to operators acting in the
antisymmetric tensor product space $\mathcal{H}^{2}$ that describe pairs of
electrons. The stationary $N$-electron states, $\left\{  \left.  D_{j}%
^{N}\right\vert 0\leq j\leq\varkappa\right\}  ,$ associated with these
Hamiltonians belong to the convex set, $S_{N},$ of positive, normalized, trace
class operators acting in $\mathcal{H}^{N}$ that is defined by
\begin{subequations}
\label{NDensityOP}%
\begin{equation}
S_{N}=\left\{  D\in\mathcal{B}_{1}\left(  \mathcal{H}^{N}\right)
\subseteq\mathcal{B}\left(  \mathcal{H}^{N}\right)  \left\vert
\operatorname*{Tr}\nolimits_{N}\left\{  D^{N}\right\}  =1,D^{N}\geq O\right.
\right\}  ,
\end{equation}
where $\mathcal{B}\left(  \mathcal{H}^{N}\right)  $ and $\mathcal{B}%
_{1}\left(  \mathcal{H}^{N}\right)  $ are the Banach spaces of bounded and
trace class operators \cite{BanachSpace} respectively acting in $\mathcal{H}%
^{N}$ and $\operatorname*{Tr}\nolimits_{N}\left\{  {}\right\}  $ is the trace
over $\mathcal{H}^{N}.$ The energies, $\left\{  \left.  E_{Nj}\right\vert
0\leq j\leq\varkappa\right\}  ,$ of these stationary states are the stationary
values of the restriction of the linear functional
\end{subequations}
\begin{equation}
E_{H}\left(  D\right)  =\operatorname*{Tr}\nolimits_{N}\left\{  HD\right\}
\label{LinFun}%
\end{equation}
to the convex set $S_{N}$ at the points $\left\{  D_{j}^{N}\right\}  .$ In
particular the ground state energy, $E_{N0},$ is the solution to the
constrained minimization of the linear functional $E_{H}$ over this convex set
i.e.
\begin{equation}
E_{N0}=E_{H}\left(  D_{0}^{N}\right)  =\min_{D\in S_{N}}\operatorname*{Tr}%
\nolimits_{N}\left\{  HD\right\}  . \label{Minval}%
\end{equation}
As $E_{H}$ is linear its minimum value occurs at extreme points,
$\operatorname*{ext}\left\{  S_{N}\right\}  ,$ of $S_{N},$ thus one only has
to the check the values of $E_{H}$ at these points$,$ which \emph{are known}
to be one dimensional projectors onto subspaces of $\mathcal{H}^{N}$ and thus
correspond to, up to a phase factor, unit vectors, $\left\vert \Psi
\right\rangle ,$ in $\mathcal{H}^{N}.$ The values of $E_{H}$ at these points
is given by
\begin{equation}
E_{H}\left(  D\right)  =\operatorname*{Tr}\nolimits_{N}\left\{  HD\right\}
=\left\langle \left.  \Psi\right\vert H\Psi\right\rangle , \label{Enval}%
\end{equation}
which can be computed to any prescribed degree of accuracy by introducing
finite dimensional approximations to $\mathcal{H}^{1}$ and hence
$\mathcal{H}^{N}$. There is a natural duality between $\mathcal{B}\left(
\mathcal{H}^{N}\right)  $ and $\mathcal{B}_{1}\left(  \mathcal{H}^{N}\right)
$ given by
\begin{equation}
\left(  X\left\vert Y\right.  \right)  _{\mathcal{H}^{N}}=\operatorname*{Tr}%
\nolimits_{N}\left\{  X^{\dagger}Y\right\}  ;X\in\mathcal{B}\left(
\mathcal{H}^{N}\right)  ,Y\in\mathcal{B}_{1}\left(  \mathcal{H}^{N}\right)  ,
\end{equation}
where $X^{\dagger}$ is the adjoint of $X,$ which can be extended to unbounded
linear functionals corresponding to unbounded operators acting in
$\mathcal{H}^{N}$ and allows the expression eq. $\left(  \ref{LinFun}\right)
$ to be written as
\begin{equation}
E_{H}\left(  D\right)  =\left(  H\left\vert D\right.  \right)  _{\mathcal{H}%
^{N}}.
\end{equation}
As the Hamiltonian $H\in\mathfrak{H}_{N2}$ the linear problem in eq. $\left(
\ref{MinVal}\right)  $ can be equivalently expressed as
\begin{align}
E_{N0}  &  =\min_{D\in S_{N}}\left(  H\left\vert D\right.  \right)
_{\mathcal{H}^{N}}=\min_{D\in S_{N}}\operatorname*{Tr}\left\{  HD\right\}
\label{MinVal2}\\
&  =\min_{D\in\operatorname*{ext}\left\{  S_{2N}\right\}  }\left(
K_{2N}\left\vert D\right.  \right)  _{\mathcal{H}^{2}}=\min_{D\in
\operatorname*{ext}\left\{  S_{2N}\right\}  }\operatorname*{Tr}\left\{
K_{2N}D\right\}  ,
\end{align}
(If the ground state is not in the continuum then $\min$ is correct and not
$\inf)$ where $K_{2N}$ is the reduced Hamiltonian operator that acts in
$\mathcal{H}^{2}$ which corresponds to $H$ through the injective expansion map
$\Gamma_{2}^{N}$ and $S_{2N}$ is the convex set of Second Order Reduced
Density Operators (SORDO's) contained in $S_{2}$ that corresponds to $S_{2N}$
through the surjective contraction map $L_{N}^{2}$ that are adjoints of each
other i.e.
\begin{equation}
\left(  H\left\vert L_{N}^{2}\left(  D\right)  \right.  \right)
_{\mathcal{H}^{2}}=\left(  \Gamma_{2}^{N}\left(  K_{2N}\right)  \left\vert
D\right.  \right)  _{\mathcal{H}^{N}}.
\end{equation}


Finite dimensional approximations of the functional eq. $\left(
\ref{LinFun}\right)  $ determining the energy of $N$-electron states are
obtained by the introduction of subsets of orthonormal bases of $\mathcal{H}%
^{1}$ and thus $\mathcal{H}^{N}.$ The matrices $\mathbb{D}^{N}$ representing
the operators $D^{N}$ satisfy the conditions
\begin{align}
\operatorname*{Tr}\left\{  \mathbb{D}^{N}\right\}   &  =1\nonumber\\
\mathbb{D}^{N}  &  \geq\mathbb{O} \label{NDensity}%
\end{align}
and are determined by
\begin{equation}
\eta_{N}=\frac{1}{2}\binom{r}{N}\left(  \binom{r}{N}+1\right)  \sim r^{2N}
\label{NDenSize}%
\end{equation}
complex parameters, where $r$ is the rank of the orthonormal basis set that
spans a finite dimensional approximation to $\mathcal{H}^{1}.$ SORDO's are
represented by Second Order Reduced Density Matrices (SORDM's), $\mathbb{D}%
_{2}^{N},$ which are determined by
\begin{equation}
\eta_{2}=\frac{1}{2}\binom{r}{2}\left(  \binom{r}{2}+1\right)  \sim r^{4}
\label{2DenSize}%
\end{equation}
complex parameters. A crucial difference is that $\eta_{2}$ does not depend on
the number of electrons while $\eta_{N}$ does, hence
\begin{equation}
\frac{\eta_{N}}{\eta_{2}}\sim r^{2N-4}, \label{RatioSize}%
\end{equation}
which increases rapidly as the number of electrons increases. Just
concentrating on the ratio of the number of parameters is however misleading
as their values are constrained and the complexity of the constraining
equations must be taken into account. In the case of $N$-electron density
matrices the constraints are eq. $\left(  \ref{NDensity}\right)  ,$ while in
the case of SORDM's they are
\begin{equation}
\mathbb{D}_{2}^{N}\in S_{2N}. \label{2Density}%
\end{equation}
The convex set described in eq. $\left(  \ref{NDensity}\right)  $ is the
positive portion of the unit hypersphere in the space of trace class
operators, $\mathcal{B}\left(  \mathcal{H}^{N}\right)  ,$ acting in the
Hilbert space $\mathcal{H}^{N},$ whose extreme points are well characterized,
in contrast the extreme points of the convex set described in eq. $\left(
\ref{2Density}\right)  $ have not all been described and the search for a
description is called the $N$-representability problem, first explicitly
introduced by Husimi \cite{Husimi} then Coulson \cite{Coulson}.

The formulation of the problem can be extended to distinguish between pure
states, ensemble (also known as mixed) states, states that describe coherent
and incoherent mixtures of varying number of particles as well as between
finite and infinite dimensional versions of the problem. In this paper
"representability" is to be interpretted as "ensemble representability" in
contrast to "pure state representability" and we consider the finite
dimensional version of the problem. Solutions to the infinite dimensional
problem would then be expressed as limits of the constraints that we derive
and their validity depend on the continuity of the constraining relationships
to the topology the limits are taken in.

Even though the significant physical interactions that determine the nature of
the state are most pair wise the fact that electrons are fermions implies a
basic underlying exchange symmetry [thus a conservation principle and an
invariance group (the Symmetric Group)] that can be thought of as an
$N$-electron interaction. This basic symmetry interaction determines a class
of possible electron states from which a physical state is determined by
coulombic interaction between electrons and electrons and nuclei, the kinetic
energy of the electrons and the presence of electric and magnetic fields.
States determined by the $\eta_{2}$ parameters are called \emph{reduced
states} in contrast to those determined by those determined by the $\eta_{N}$
parameters. Every state determines a reduced state however the set of
conditions that intrinsically characterize reduced states are not known and
constitute the $N$-representability problem.

We formulate the problem of determining whether an element of the set of
positive trace class operators, $\mathcal{B}_{1+}\left(  \mathcal{H}%
^{2}\right)  ,$ acting in $\mathcal{H}^{2}$ actually is a reduced state by
considering whether the bounded positive linear functional that it determines
on the set of bounded operators, $\mathcal{B}\left(  \mathcal{H}^{2}\right)
,$ acting in $\mathcal{H}^{2}$ can be extended to a positive bounded linear
functional acting on the set of bounded operators acting in $\mathcal{H}^{N}$
i.e. we formulate the conditions for a positive trace class operator acting in
$\mathcal{H}^{2}$ to be $N$-representable in terms of whether it can be
extended to a bounded positive linear functional on $\mathcal{B}\left(
\mathcal{H}^{N}\right)  .$

In order to do this we observe that operators in $\mathcal{B}\left(
\mathcal{H}^{N}\right)  $ for any $N$ are determined by the values of linear
functionals on the Fermion Fock algebra, which we review in section $\left(
{}\right)  .$ The Fermion Fock algebra is a $w^{\ast}$-algebra and thus one
can apply general results from this context to the extension problem in the
Fock algebra and hence the representability and $N$-representability problems.
Emch has discussed the conditions that a positive functional defined on a
subspace of an algebra must satisfy in order to be extend to a state thus be a
partial state, while Segal has extended consideration to subspaces not
containing the identity and Kadison has shown how these conditions can be
expressed in terms of values of the linear functionals on idempotents, i.e.
projectors, in these subspaces.

The subspaces associated with SORDO's are the spaces formed by one and two
particle second quantized operators. We show that the idempotents within these
subspaces are decomposable i.e products of one particle or one hole operators
and can be all generated from a reference one particle basis by the action of
the unitary group $\mathcal{U}\left(  \mathcal{H}^{1}\right)  $ of one
particle space $\mathcal{H}^{1}.$ We show that Kadisons theorem leads to
generalizations of the P, Q and G conditions which are \emph{necessary and
sufficient }for\emph{ }representability.

If a functional is $N$-representable then its values on one particle operators
are defined by its values on two particle operators and one need only check
the \emph{necessary and sufficient} conditions that the functional must
satisfy in order to be a reduced $N$-particle state on decomposable two
particle idempotents and sums of one and two particle idempotents as discussed
in section $\left(  {}\right)  $.

The values of functionals on every decomposable two particle idempotent can be
expressed in terms of the action of the one particle unitary group
$\mathcal{U}\left(  \mathcal{H}^{1}\right)  $ on a reference decomposable
two-particle state i.e. a two particle Independent Particle State (IPS).

The group $\mathcal{U}\left(  \mathcal{H}^{1}\right)  $ can be expressed as a
coset factorization with respect to the isotropy group of a two particle
reference IPS and that generates a complex parameterization of the group by
considering elements of the coset to be products of elements of the general
linear group $GL\left(  \mathcal{H}^{1}\right)  $. The action of
$\mathcal{U}\left(  \mathcal{H}^{1}\right)  $ on a two particle IP\ reference
state produces a complex manifold of two particle Thouless Coherent States
(TCS's). The values of the functionals on the idempotents then correspond to
values of the TCS Q-symbols on this complex manifold. The \emph{necessary and
sufficient }conditions for a $N$-representable extension to exist can then be
totally recast in terms of \emph{known }norm, positivity and linear
constraints on Q-symbols and by using the P-symbol of the reduced $N$-particle
Hamiltonian (which is dual to its Q-symbol) the $N$-particle energy
functional, $E_{N}\left(  Q\right)  $ can be formulated as a linear functional
of the Q-symbols. We explicitly give the transformation from the Q-symbol of
the reduced Hamiltonian to its P-symbol.

Finding the ground electronic state can then be expressed as Linear Programing
(LP) problem or a Semi Definite Programing (SDP) one and in the infinite
dimensional case is equivalent to solving the Time Independent Schr\H{o}dinger
Equation (TISE). In the finite dimensional case, produced by finite
dimensional approximations to the space of Q-symbols, approximations of the
TISE which converge to the exact solution.

\section{Background \ref{background}}

\subsection{Fermion Fock Algebra \ref{FockAlgebra}}

A fermionic state is a normalized, bounded positive linear function, $f$,
defined on the Fermion Fock Algebra $\mathcal{F}$ $\left(  \text{Appendix
\ref{Fock} }\right)  $
\begin{equation}
f:\mathcal{F\rightarrow}\mathbb{R},
\end{equation}
where $\mathcal{F}$ is generated by polynomials of second quantized operators
that satisfy the Canonical Anticommutation Relationships (CAR) and form the
algebra of bounded operators $\mathcal{B}\left(  \mathcal{H}_{\mathcal{F}%
}\right)  $ defined on the Fermion Fock space $\mathcal{H}_{\mathcal{F}}.$
(One often needs the values of states on unbounded operators as physical
Hamiltonians are often unbounded and many methods have been developed to
obtain such values from values on bounded operators). A state is called a
normal state if it can be identified with a trace class operator, otherwise it
is an abnormal state. In this paper only normal states are considered. Normal
states defined on the Fermion Fock algebra are in one to one correspondence
with normalized density operators i.e.
\begin{equation}
f\left(  X\right)  =\operatorname*{Tr}\left\{  XD_{f}\right\}  ;\,D_{f}\in
S_{\mathcal{F}},
\end{equation}
where the convex set, $S_{\mathcal{F}},$ of density operators in $\mathcal{F}$
is given by
\begin{equation}
S_{\mathcal{F}}=\left\{  D\left\vert D\in\mathcal{B}_{1}\left(  \mathcal{H}%
_{\mathcal{F}}\right)  ;D\geq0;\operatorname{Tr}D=1\right.  \right\}
\end{equation}
and $\mathcal{B}_{1}\left(  \mathcal{H}_{\mathcal{F}}\right)  $ is the space
of trace class operators. The scalar product%
\begin{equation}
\left(  \left.  \cdot\right\vert \left.  \cdot\right.  \right)  :\mathcal{B}%
\left(  \mathcal{H}_{\mathcal{F}}\right)  \times\mathcal{B}_{1}\left(
\mathcal{H}_{\mathcal{F}}\right)  \rightarrow\mathbb{C}%
\end{equation}
can be expressed in terms of the trace operation in $\mathcal{B}\left(
\mathcal{H}_{\mathcal{F}}\right)  $ as
\begin{equation}
\left(  \left.  A\right\vert B\right)  =\operatorname*{Tr}\left\{  A^{\dagger
}B\right\}  ;\quad A\in\mathcal{B}\left(  \mathcal{H}_{\mathcal{F}}\right)
,B\in\mathcal{B}_{1}\left(  \mathcal{H}_{\mathcal{F}}\right)  .
\end{equation}


\subsection{Contraction and Expansion Maps \ref{ContExpan}}

In order to analyse representability conditions one needs to consider two
forms of the decomposition of the Fock algebra $\mathcal{F},$ one in terms of
operators on and between subspaces associated with a definite number of
particles i.e.
\begin{equation}
\mathcal{F}=%
%TCIMACRO{\dbigoplus \limits_{0\leq N,M\leq r}}%
%BeginExpansion
{\displaystyle\bigoplus\limits_{0\leq N,M\leq r}}
%EndExpansion
\mathcal{B}\left(  \mathcal{H}^{N},\mathcal{H}^{M}\right)  ,
\end{equation}
(note that $\mathcal{B}\left(  \mathcal{H}^{N}\right)  =\mathcal{B}\left(
\mathcal{H}^{N},\mathcal{H}^{N}\right)  $) and the other in terms of
polynomials of second quantized operators i.e.
\begin{equation}
\mathcal{F}=%
%TCIMACRO{\dbigoplus \limits_{0\leq N,M\leq r}}%
%BeginExpansion
{\displaystyle\bigoplus\limits_{0\leq N,M\leq r}}
%EndExpansion
\mathcal{B}\left(  \mathcal{H}_{\mathcal{F}}\right)  _{NM},
\end{equation}
where
\begin{equation}
\mathcal{B}\left(  \mathcal{H}_{\mathcal{F}}\right)  _{NM}=\left\{
\sum_{\substack{1\leq j_{1}<\cdots<j_{N}\leq r\\1\leq k_{1}<\cdots<k_{M}\leq
r}}X_{j_{1}\cdots j_{N}k_{1}\cdots k_{N}}a_{j_{1}}^{\dagger}\cdots a_{j_{N}%
}^{\dagger}a_{k_{M}}\cdots a_{k_{1}}\right\}
\end{equation}
and $\mathcal{B}\left(  \mathcal{H}_{\mathcal{F}}\right)  _{00}$ is the one
dimensional space generated by the identitiy operator $I_{\mathcal{H}%
_{\mathcal{F}}}.$

\subsection{Expansion Map \label{Expansion copy(1)}}

The non singular expansion map \cite{weiner2011}\cite{Kummar} $\left(
\text{Appendix \ref{Expan}}\right)  $
\begin{align}
\Gamma &  :\mathcal{F\rightarrow F}\nonumber\\
\Gamma &  :%
%TCIMACRO{\dbigoplus \limits_{0\leq N,M\leq r}}%
%BeginExpansion
{\displaystyle\bigoplus\limits_{0\leq N,M\leq r}}
%EndExpansion
\mathcal{B}\left(  \mathcal{H}^{N},\mathcal{H}^{M}\right)
\mathcal{\rightarrow}%
%TCIMACRO{\dbigoplus \limits_{0\leq N,M\leq r}}%
%BeginExpansion
{\displaystyle\bigoplus\limits_{0\leq N,M\leq r}}
%EndExpansion
\mathcal{B}\left(  \mathcal{H}_{\mathcal{F}}\right)  _{NM}%
\end{align}
has components
\begin{equation}
\Gamma_{MN}\left(  X\right)  =\sum_{0\leq P\leq M,0\leq Q\leq N}\sqrt
{\binom{M}{P}\binom{N}{Q}}\Gamma_{PQ}^{MN}\left(  X_{PQ}\right)
\delta_{M\left(  N+P-Q\right)  }.
\end{equation}
that can be expressed in terms of elementary the $\left(  P,Q\right)
\rightarrow\left(  M,N\right)  $ expansion maps, $\Gamma_{PQ}^{MN},$ as
\begin{equation}
\Gamma_{PQ}^{MN}\left(  X_{PQ}\right)  =\sum_{1\leq j_{1}<\cdots<j_{N}%
\leq\infty}\left\vert \varphi_{l_{1}}\cdots\varphi_{l_{P}}\varphi_{j_{1}%
}\cdots\varphi_{j_{N}}\right\rangle \left\langle \varphi_{m_{1}}\cdots
\varphi_{m_{Q}}\varphi_{j_{1}}\cdots\varphi_{j_{N}}\right\vert .
\end{equation}
The expansion map is actually the second quantization map (Appendix
\ref{Second}) i.e.
\begin{equation}
\Gamma\left(  \left\vert \varphi_{j_{1}}\cdots\varphi_{j_{N}}\right\rangle
\left\langle \varphi_{k_{1}}\cdots\varphi_{k_{M}}\right\vert \right)
=a_{j_{1}}^{\dagger}\cdots a_{j_{N}}^{\dagger}a_{k_{M}}\cdots a_{k_{1}};0\leq
M,N\leq\infty,1\leq j_{1}<\cdots<j_{N}\leq\infty,1\leq k_{1}<\cdots<k_{N}%
\leq\infty,
\end{equation}
where $M=0$ or $N=0$ corresponds to the vacuum state vector $\left\vert
\phi\right\rangle $ and
\begin{equation}
\Gamma\left(  \left\vert \phi\right\rangle \left\langle \phi\right\vert
\right)  =I_{\mathcal{H}_{\mathcal{F}}}.
\end{equation}


In the study of the $N$-representability problems one is only interested in
the , $\Gamma_{e}=$ $\left.  \Gamma\right\vert _{\mathcal{F}_{e}},$ to
particle conserving operators i.e.
\begin{equation}
\Gamma_{e}:%
%TCIMACRO{\dbigoplus \limits_{0\leq N\leq r}}%
%BeginExpansion
{\displaystyle\bigoplus\limits_{0\leq N\leq r}}
%EndExpansion
\mathcal{B}\left(  \mathcal{H}^{N}\right)  \rightarrow%
%TCIMACRO{\dbigoplus \limits_{0\leq N\leq r}}%
%BeginExpansion
{\displaystyle\bigoplus\limits_{0\leq N\leq r}}
%EndExpansion
\mathcal{B}\left(  \mathcal{H}_{\mathcal{F}}\right)  _{N}%
\end{equation}
where $\mathcal{B}\left(  \mathcal{H}^{N}\right)  \equiv$ $\mathcal{B}\left(
\mathcal{H}^{N,N}\right)  $, and $\mathcal{B}\left(  \mathcal{H}_{\mathcal{F}%
}\right)  _{N}\equiv\mathcal{B}\left(  \mathcal{H}_{\mathcal{F}}\right)
_{N,N}.$ The restriction, which is also non singular, can be expanded in
component form as
\begin{equation}
\Gamma\left(  X\right)  =\sum_{0\leq N\leq\infty}\Gamma_{N}\left(  X\right)  ,
\end{equation}
where the components are defined in terms of the elementary $P\rightarrow N$
expansion maps \cite{Kummar} (Appendix \ref{Expansion}) as
\begin{equation}
\Gamma_{N}\left(  X\right)  =\sum_{0\leq P\leq N}\binom{N}{P}\Gamma_{P}%
^{N}\left(  X_{P}\right)  .
\end{equation}
The inverse of the restricted expansion map is given by \cite{Kummar70}
\begin{equation}
\Gamma_{e}^{-1}\left(  X\right)  =\sum_{0\leq N\leq\infty}\Gamma_{N}%
^{-1}\left(  X\right)  ,
\end{equation}
where%
\begin{equation}
\Gamma_{N}^{-1}\left(  X\right)  =\sum_{0\leq P\leq N}\left(  -1\right)
^{N-P}\binom{N}{P}\Gamma_{P}^{N}\left(  X_{P}\right)  .
\end{equation}


\subsection{Contraction Map \label{Contraction copy(1)}}

The expansion map $\Gamma$ is the adjoint of the contraction map, $L,$
\cite{Kummar70}
\begin{equation}
L:\mathcal{B}_{1}\left(  \mathcal{H}_{\mathcal{F}}\right)  \rightarrow
\mathcal{B}_{1}\left(  \mathcal{H}_{\mathcal{F}}\right)  ,
\end{equation}
i.e. $\Gamma=L^{\dagger},$ with respect to inner product $\left(
\cdot\left\vert \cdot\right.  \right)  $
\begin{equation}
\left(  \Gamma\left(  X\right)  \left\vert Y\right.  \right)  =\left(
X\left\vert L\left(  Y\right)  \right.  \right)  \quad\forall Y\in
\mathcal{B}_{1}\left(  \mathcal{H}_{\mathcal{F}}\right)  .
\end{equation}
The contraction map transforms between the second quantized and first
quantized decompositions of $\mathcal{F}$
\begin{equation}
L:%
%TCIMACRO{\dbigoplus \limits_{0\leq N,M\leq r}}%
%BeginExpansion
{\displaystyle\bigoplus\limits_{0\leq N,M\leq r}}
%EndExpansion
\mathcal{B}_{1}\left(  \mathcal{H}_{\mathcal{F}}\right)  _{N,M}\rightarrow%
%TCIMACRO{\dbigoplus \limits_{0\leq N,M\leq r}}%
%BeginExpansion
{\displaystyle\bigoplus\limits_{0\leq N,M\leq r}}
%EndExpansion
\mathcal{B}_{1}\left(  \mathcal{H}^{N},\mathcal{H}^{M}\right)  .
\end{equation}
and is defined by its action on basis elements
\begin{equation}
L\left(  Y\right)  =\sum_{0\leq N,M\leq r}\sum_{\substack{1\leq j_{1}%
<\cdots<j_{M}\leq r\\1\leq k_{1}<\cdots<k_{N}\leq r}}\left(  \left.  a_{k_{1}%
}^{\dagger}\cdots a_{k_{N}}^{\dagger}a_{j_{M}}\cdots a_{j_{1}}\right\vert
Y\right)  \left\vert \varphi_{j_{1}}\cdots\varphi_{j_{M}}\right\rangle
\left\langle \varphi_{k_{1}}\cdots\varphi_{k_{N}}\right\vert ,
\end{equation}
where $M=0$ or $N=0$ corresponds to the vacuum state vector $\left\vert
\phi\right\rangle $ and
\begin{equation}
L\left(  Y_{00}\right)  =\operatorname*{Tr}\left\{  Y\right\}  \left\vert
\phi\right\rangle \left\langle \phi\right\vert .
\end{equation}


The restriction, $L_{e},$ to $\mathcal{F}_{e}$, which is also non singular,
can be expanded in component form as
\begin{equation}
L_{e}\left(  Y\right)  =\sum_{0\leq P\leq\infty}L_{eP}\left(  Y\right)  ,
\end{equation}
where the components are defined in terms of the standard $N\rightarrow P$
contraction maps \cite{Kummar} (Appendix \ref{Expansion}) as
\begin{equation}
L_{eP}\left(  Y\right)  =\sum_{P\leq N\leq\infty}\binom{N}{P}L_{N}^{P}\left(
Y_{N}\right)  .
\end{equation}
The inverse of the restricted contraction map is given by \cite{Kummar70}
\begin{equation}
L_{e}^{-1}\left(  Y\right)  =\sum_{0\leq P\leq\infty}L_{eP}^{-1}\left(
Y\right)  ,
\end{equation}
where%
\begin{equation}
L_{eP}^{-1}\left(  Y\right)  =\sum_{P\leq N\leq r}\left(  -1\right)
^{N-P}\binom{N}{P}L_{N}^{P}\left(  Y_{N}\right)  .
\end{equation}


\subsection{Unitary Group of $\mathcal{H}^{1}$}

The unitary group, $U\left(  \mathcal{H}^{1}\right)  $ of one particle space
$\mathcal{H}^{1}$ is defined by
\begin{equation}
U\left(  \mathcal{H}^{1}\right)  =\left\{  \left.  u\right\vert \left\langle
\left.  u\varphi\right\vert u\psi\right\rangle _{1}=\left\langle \left.
\varphi\right\vert \psi\right\rangle _{1},\forall\left\vert \varphi
\right\rangle ,\left\vert \psi\right\rangle \in\mathcal{H}^{1}\right\}  ,
\end{equation}
where $\left\langle \left.  {}\right\vert \right\rangle _{1}$ is the inner
product in $\mathcal{H}^{1}.$ In particular both first and second quantized
forms of $\mathcal{F}_{e}$ are composed of invariant subspaces of $U\left(
\mathcal{H}^{1}\right)  $ i.e.
\begin{equation}
\mathcal{R}_{\mathcal{F}_{e}}\left(  U\left(  \mathcal{H}^{1}\right)  \right)
:\mathcal{B}\left(  \mathcal{H}^{N}\right)  \rightarrow\mathcal{B}\left(
\mathcal{H}^{N}\right)  ;\quad0\leq N\leq r
\end{equation}
and
\begin{equation}
\mathcal{R}_{\mathcal{F}_{e}}\left(  U\left(  \mathcal{H}^{1}\right)  \right)
:\mathcal{B}\left(  \mathcal{H}\right)  _{N}\rightarrow\mathcal{B}\left(
\mathcal{H}\right)  _{N};\quad0\leq N\leq r,
\end{equation}
that produce irreducible representations, $\mathcal{R}_{\mathcal{F}_{e}NF}$
and $\mathcal{R}_{\mathcal{F}_{e}NS}$ respectively, where
\begin{align}
\mathcal{R}_{\mathcal{F}_{e}NF}\left(  U\left(  \mathcal{H}^{1}\right)
\right)   &  =\left.  \mathcal{R}_{\mathcal{F}_{e}}\left(  U\left(
\mathcal{H}^{1}\right)  \right)  \right\vert _{\mathcal{B}\left(
\mathcal{H}^{N}\right)  }\\
\mathcal{R}_{\mathcal{F}_{e}NS}\left(  U\left(  \mathcal{H}^{1}\right)
\right)   &  =\left.  \mathcal{R}_{\mathcal{F}_{e}}\left(  U\left(
\mathcal{H}^{1}\right)  \right)  \right\vert _{\mathcal{B}\left(
\mathcal{H}\right)  _{N}}.
\end{align}
The representation $\mathcal{R}_{\mathcal{F}_{e}}$ is defined as
\begin{equation}
\left[  \mathcal{R}_{\mathcal{F}_{e}}\left(  u\right)  \right]  \left(
X\right)  =\widehat{\mathcal{R}_{\mathcal{H}_{\mathcal{F}}}\left(  u\right)
}\left(  X\right)  =\mathcal{R}_{\mathcal{H}_{\mathcal{F}}}\left(  u\right)
X\mathcal{R}_{\mathcal{H}_{\mathcal{F}}}\left(  u\right)  ^{\dagger},
\label{AdjointRep}%
\end{equation}
where
\begin{equation}
\mathcal{R}_{\mathcal{H}_{\mathcal{F}}}\left(  u\right)  =\exp\left\{
i\sum_{1\leq j,k\leq r}\lambda_{\varphi jk}\left(  u\right)  a_{j}^{\dagger
}a_{k}\right\}  ,\lambda_{\varphi jk}\left(  u\right)  =\overline
{\lambda_{\varphi kj}\left(  u\right)  }\in\mathbb{C}. \label{lamda}%
\end{equation}
The complex variables $\left\{  \lambda_{\varphi jk}\right\}  $ are a set of
coordinates for the manifold formed by the group $U\left(  \mathcal{H}%
^{1}\right)  $ i.e form a parameterization of $U\left(  \mathcal{H}%
^{1}\right)  .$

The maps $\Gamma,L$ both commute with $\mathcal{R}_{\mathcal{F}_{e}}\left(
U\left(  \mathcal{H}^{1}\right)  \right)  $ i.e.
\begin{equation}
\left[  \Gamma\circ\mathcal{R}_{\mathcal{F}_{e}}\left(  U\left(
\mathcal{H}^{1}\right)  \right)  \right]  \left(  X\right)  =\left[
\mathcal{R}_{\mathcal{F}_{e}}\left(  U\left(  \mathcal{H}^{1}\right)  \right)
\circ\Gamma\right]  \left(  X\right)
\end{equation}
and
\begin{equation}
\left[  L\circ\mathcal{R}_{\mathcal{F}_{e}}\left(  U\left(  \mathcal{H}%
^{1}\right)  \right)  \right]  \left(  X\right)  =\left[  \mathcal{R}%
_{\mathcal{F}_{e}}\left(  U\left(  \mathcal{H}^{1}\right)  \right)  \circ
L\right]  \left(  X\right)
\end{equation}
and the components of these maps intertwine with the irreducible
representations $\left\{  \mathcal{R}_{\mathcal{F}_{e}NF}\right\}  $ i.e.
\begin{equation}
\Gamma_{P}^{N}\circ\mathcal{R}_{\mathcal{F}_{e}PF}=\mathcal{R}_{\mathcal{F}%
_{e}NF}\circ\Gamma_{P}^{N}%
\end{equation}%
\begin{equation}
L_{N}^{P}\circ\mathcal{R}_{\mathcal{F}_{e}PF}=\mathcal{R}_{\mathcal{F}_{e}%
NF}\circ L_{N}^{P}.
\end{equation}


The representation $\mathcal{R}_{\mathcal{F}_{e}}$ is actually identical to
the representation $\mathcal{R}_{\mathcal{F}}$ which is also given by the
right hand side of eq. $\left(  \ref{AdjointRep}\right)  .$

\section{Reduced Density Operators \label{Reduced Density Operators}}

In order to obtain representability conditions for SORDO's the appropriate
subspace of $\mathcal{B}\left(  \mathcal{H}_{\mathcal{F}}\right)  $ to
consider is the direct sum formed by one electron operators $\mathcal{B}%
\left(  \mathcal{H}_{\mathcal{F}}\right)  _{1}$ and two electron
$\mathcal{B}\left(  \mathcal{H}_{\mathcal{F}}\right)  _{2}$ operators i.e.
\begin{equation}
\Lambda_{2}=%
%TCIMACRO{\dbigoplus \limits_{1\leq P\leq2}}%
%BeginExpansion
{\displaystyle\bigoplus\limits_{1\leq P\leq2}}
%EndExpansion
\mathcal{B}\left(  \mathcal{H}_{\mathcal{F}}\right)  _{P},
\end{equation}
as reduced states defined on $\Lambda_{2}$ correspond to SORDO's through%
\begin{equation}
D_{2}=L_{2}\left(  D\right)  =\sum_{\substack{1\leq j_{1}<j_{2}\leq r\\1\leq
k_{1}<k_{2}\leq r}}f\left(  a_{j_{1}}^{\dagger}a_{j_{2}}^{\dagger}a_{k_{2}%
}a_{k_{1}}\right)  \left\vert \varphi_{j_{1}}\varphi_{j_{2}}\right\rangle
\left\langle \varphi_{k_{1}}\varphi_{k_{2}}\right\vert ,
\end{equation}
where%
\begin{equation}
f\left(  a_{j_{1}}^{\dagger}a_{j_{2}}^{\dagger}a_{k_{2}}a_{k_{1}}\right)
=\left(  \left.  a_{k_{1}}^{\dagger}a_{k_{2}}^{\dagger}a_{j_{2}}a_{j_{1}%
}\right\vert D\right)
\end{equation}
and to FORDO's through
\begin{equation}
D_{1}=L_{1}\left(  D\right)  =\sum_{\substack{1\leq j_{1}\leq r\\1\leq
k_{1}\leq r}}f\left(  a_{j_{1}}^{\dagger}a_{k_{1}}\right)  \left\vert
\varphi_{j_{1}}\right\rangle \left\langle \varphi_{k_{1}}\right\vert ,
\end{equation}
where%

\begin{equation}
f\left(  a_{j_{1}}^{\dagger}a_{k_{1}}\right)  =\left(  \left.  a_{k_{1}%
}^{\dagger}a_{j_{1}}\right\vert D\right)
\end{equation}
when $D\in S_{\mathcal{F}}.$

The action of the group $U\left(  \mathcal{H}^{1}\right)  $ on states and
first and second order reduced states is given by%
\begin{equation}
D\left(  u\right)  =\hat{u}\left(  D\right)  =uDu^{\dagger},\quad u\in
U\left(  \mathcal{H}^{1}\right)
\end{equation}
leading to
\begin{align}
D_{f2}\left(  u\right)   &  =L_{2}\left(  D_{f}\left(  u\right)  \right)
=\sum_{\substack{1\leq j_{1}<j_{2}\leq r\\1\leq k_{1}<k_{2}\leq r}}f\left(
\hat{u}^{\dagger}\left(  a_{j_{1}}^{\dagger}a_{j_{2}}^{\dagger}a_{k_{2}%
}a_{k_{1}}\right)  \right)  \left\vert \varphi_{j_{1}}\varphi_{j_{2}%
}\right\rangle \left\langle \varphi_{k_{1}}\varphi_{k_{2}}\right\vert
\nonumber\\
&  =\sum_{\substack{1\leq j_{1}<j_{2}\leq r\\1\leq k_{1}<k_{2}\leq r}}\left(
\left.  \hat{u}^{\dagger}\left(  a_{j_{1}}^{\dagger}a_{j_{2}}^{\dagger
}a_{k_{2}}a_{k_{1}}\right)  \right\vert D_{f}\right)  =\sum_{\substack{1\leq
j_{1}<j_{2}\leq r\\1\leq k_{1}<k_{2}\leq r}}\left(  \left.  a_{k_{1}}%
^{\dagger}a_{k_{2}}^{\dagger}a_{j_{2}}a_{j_{1}}\right\vert \hat{u}\left(
D_{f}\right)  \right)
\end{align}
for the SORDO. Noting that by the intertwining relationship
\begin{equation}
\hat{u}_{2}\left(  L_{2}\left(  D\right)  \right)  =L_{2}\left(  D\left(
u\right)  \right)
\end{equation}
one obtains
\begin{align}
D_{f2}\left(  u\right)   &  =\sum_{\substack{1\leq j_{1}<j_{2}\leq r\\1\leq
k_{1}<k_{2}\leq r}}f\left(  a_{j_{1}}^{\dagger}a_{j_{2}}^{\dagger}a_{k_{2}%
}a_{k_{1}}\right)  u_{2}\left\vert \varphi_{j_{1}}\varphi_{j_{2}}\right\rangle
\left\langle \varphi_{k_{1}}\varphi_{k_{2}}\right\vert u_{2}^{\dagger}%
=\sum_{\substack{1\leq j_{1}<j_{2}\leq r\\1\leq k_{1}<k_{2}\leq r}}f\left(
a_{j_{1}}^{\dagger}a_{j_{2}}^{\dagger}a_{k_{2}}a_{k_{1}}\right)  \left\vert
\varphi_{j_{1}}\left(  u\right)  \varphi_{j_{2}}\left(  u\right)
\right\rangle \left\langle \varphi_{k_{1}}\left(  u\right)  \varphi_{k_{2}%
}\left(  u\right)  \right\vert \nonumber\\
&  =\sum_{\substack{1\leq j_{1}<j_{2}\leq r\\1\leq k_{1}<k_{2}\leq r}}f\left(
a_{j_{1}}^{\dagger}\left(  u^{\dagger}\right)  a_{j_{2}}^{\dagger}\left(
u^{\dagger}\right)  a_{k_{2}}\left(  u^{\dagger}\right)  a_{k_{1}}\left(
u^{\dagger}\right)  \right)  \left\vert \varphi_{j_{1}}\varphi_{j_{2}%
}\right\rangle \left\langle \varphi_{k_{1}}\varphi_{k_{2}}\right\vert .
\label{SORDO}%
\end{align}


One can observe that one can always find unitary operators $u_{j_{1}j_{2}}$
and $u_{k_{1}k_{2}}$ such that
\begin{equation}
a_{j_{1}}^{\dagger}\left(  u^{\dagger}\right)  a_{j_{2}}^{\dagger}\left(
u^{\dagger}\right)  a_{k_{2}}\left(  u^{\dagger}\right)  a_{k_{1}}\left(
u^{\dagger}\right)  =u^{\dagger}a_{j_{1}}^{\dagger}a_{j_{2}}^{\dagger}%
a_{k_{2}}a_{k_{1}}u=u_{j_{1}j_{2}}^{\dagger}a_{1}^{\dagger}a_{2}^{\dagger
}a_{2}a_{1}u_{k_{1}k_{2}}%
\end{equation}
and in particular
\begin{equation}
a_{j_{1}}^{\dagger}\left(  u\right)  a_{j_{2}}^{\dagger}\left(  u\right)
a_{j_{2}}\left(  u\right)  a_{j_{1}}\left(  u\right)  =u_{j_{1}j_{2}}%
^{\dagger}a_{1}^{\dagger}a_{2}^{\dagger}a_{2}a_{1}u_{j_{1}j_{2}}=\hat
{u}_{j_{1}j_{2}}^{\dagger}\left(  a_{1}^{\dagger}a_{2}^{\dagger}a_{2}%
a_{1}\right)
\end{equation}
hence the values of $f$ on the basis elements of $\mathcal{B}\left(
\mathcal{H}_{\mathcal{F}}\right)  _{2}$ can be generated from a reference
element $a_{1}^{\dagger}a_{2}^{\dagger}a_{2}a_{1}$ and $D_{f2}\left(
u\right)  $ expressed by
\begin{equation}
D_{f2}\left(  u\right)  =D_{f2}\left(  \mathbf{u}\left(  u\right)  \right)
=\sum_{\substack{1\leq j_{1}<j_{2}\leq r\\1\leq k_{1}<k_{2}\leq r}}f\left(
u_{j_{1}j_{2}}^{\dagger}a_{1}^{\dagger}a_{2}^{\dagger}a_{2}a_{1}u_{k_{1}k_{2}%
}\right)  \left\vert \varphi_{j_{1}}\varphi_{j_{2}}\right\rangle \left\langle
\varphi_{k_{1}}\varphi_{k_{2}}\right\vert ,
\end{equation}
where
\begin{equation}
\mathbf{u}\left(  u\right)  =\left(  u_{11}^{\dagger},\cdots,u_{rr}^{\dagger
}\right)  .
\end{equation}


The manifold of projectors $\mathcal{P}_{2}$ $\subset\mathcal{B}\left(
\mathcal{H}_{\mathcal{F}}\right)  _{2}$ is a homogeneous manifold of $U\left(
\mathcal{H}^{1}\right)  $ formed by the orbit generated by $a_{1}^{\dagger
}a_{2}^{\dagger}a_{2}a_{1}$ under the group $U\left(  \mathcal{H}^{1}\right)
$ i.e.
\begin{equation}
\mathcal{P}_{2}=\left\{  \left.  \hat{u}^{\dagger}\left(  a_{1}^{\dagger}%
a_{2}^{\dagger}a_{2}a_{1}\right)  \right\vert u\in U\left(  \mathcal{H}%
^{1}\right)  \right\}  ,
\end{equation}
one should note that this includes all elements formed by \emph{all
}decomposable\emph{ }bases including the other elements of the basis that
includes $\left\{  \left.  a_{j}^{\dagger},a_{j}\right\vert 1\leq
j\leq2\right\}  $ formed by $\left\{  \left.  a_{j}^{\dagger},a_{j}\right\vert
3\leq j\leq r\right\}  $.

\section{N-Representability Conditions \label{NREP_Con}}

It was shown \cite{Weiner 2011} that by finding the conditions that a linear
functional defined on a subspace generated by one and two particle operators
must satisfy in order to be extended to a state of the Fermion Fock algebra,
and be dispersion free on the number operator, one obtains ensemble
\emph{Necessary and Sufficient }$N$-representability conditions. These
conditions can be expressed in terms of diagonal matrix elements of
representations of operators belonging $\mathcal{B}_{1}\left(  \mathcal{H}%
^{2}\right)  $ produced by decomposable bases of $\mathcal{H}^{2},$ however in
order to be sufficient conditions the matrix elements must satisfy them in
\emph{all} \emph{decomposable} bases, i.e. the conditions on the matrices
produced by these elements must be conserved under the under the action of the
representations, $\mathcal{R}_{\varphi\mathcal{H}^{2}}\left(  U\left(
\mathcal{H}^{1}\right)  \right)  ,$ of the one electron unitary group
$U\left(  \mathcal{H}^{1}\right)  $ on $\mathcal{H}^{2}$ generated by the
orthonormal basis $\left\{  \left.  \left\vert \varphi_{j}\right\rangle
\right\vert 1\leq j\leq r\right\}  $ of $\mathcal{H}^{1}.$ These conditions
are given by the following list :

\begin{enumerate}
\item Normalization%

\begin{equation}
\operatorname*{Tr}\left\{  D_{2}\left(  u\right)  \right\}  =\sum_{1\leq
j<k\leq r}D_{2jkjk}\left(  u\right)  =\operatorname*{Tr}\left\{
D_{2}\right\}  =\binom{N}{2}\quad\forall u\in\mathcal{U}\left(  \mathcal{H}%
^{1}\right)  , \label{NormCon}%
\end{equation}
as the trace is invariant under unitary transformation.

\item Generalized P-Condition%
\begin{equation}
1\geq D_{2abab}\left(  u\right)  \geq0\,\forall u\in\mathcal{U}\left(
\mathcal{H}^{1}\right)  \label{PCond}%
\end{equation}


\item Generalized Q-Condition%
\begin{equation}
1\geq\left(  N-1\right)  ^{-1}\sum_{1\leq l\leq r}D_{2alal}\left(  u\right)
-D_{2abab}\left(  u\right)  \geq0\,\forall u\in\mathcal{U}\left(
\mathcal{H}^{1}\right)  , \label{Qcond}%
\end{equation}
it should be observed that this condition is equivalent to
\begin{equation}
1\geq\left(  N-1\right)  ^{-1}\sum_{1\leq l\leq r}D_{2blbl}\left(  u\right)
-D_{2abab}\left(  u\right)  \geq0\,\forall u\in\mathcal{U}\left(
\mathcal{H}^{1}\right)  ,
\end{equation}
as one can always make the variable transformation
\begin{equation}
u\rightarrow wu,
\end{equation}
where
\begin{equation}
D_{2blbl}\left(  u\right)  =D_{2alal}\left(  wu\right)  .
\end{equation}


\item Generalized G-Condition%
\begin{equation}
1\geq\left(  N-1\right)  ^{-1}\sum_{1\leq l\leq r}\left\{  D_{2alal}\left(
u\right)  +D_{2blbl}\left(  u\right)  \right\}  -D_{2abab}\left(  u\right)
\geq0\,\,\forall u\in\mathcal{U}\left(  \mathcal{H}^{1}\right)  ,
\label{GCond}%
\end{equation}

\end{enumerate}

where the matrix elements $\left\{  D_{2jkjk}\left(  u\right)  \right\}  $ can
be considered as the values of a linear functional $f$ on the subspace
$\mathcal{B}\left(  \mathcal{H}_{\mathcal{F}}\right)  _{2}$ given by
\begin{equation}
D_{2jkjk}\left(  u\right)  =f\left(  \widehat{u}^{\dagger}\left(
a_{j}^{\dagger}a_{k}^{\dagger}a_{k}a_{j}\right)  \right)  =f\left(
a_{j}^{\dagger}\left(  u\right)  a_{k}^{\dagger}\left(  u\right)  a_{k}\left(
u\right)  a_{j}\left(  u\right)  \right)  =f\left(  u_{jk}^{\dagger}%
a_{1}^{\dagger}a_{2}^{\dagger}a_{2}a_{1}u_{jk}\right)  . \label{matel}%
\end{equation}


The set
\begin{equation}
\left\{  \left.  D_{21212}\left(  u\right)  =f\left(  u^{\dagger}%
a_{1}^{\dagger}a_{2}^{\dagger}a_{2}a_{1}u\right)  \right\vert u\in
\mathcal{U}\left(  \mathcal{H}^{1}\right)  \right\}  \label{allmat}%
\end{equation}
clearly contains the subsets
\begin{equation}
\left\{  \left.  D_{2jkjk}\left(  u\right)  =f\left(  u_{jk}^{\dagger}%
a_{1}^{\dagger}a_{2}^{\dagger}a_{2}a_{1}u_{jk}\right)  \right\vert u_{jk}%
\in\mathcal{U}\left(  \mathcal{H}^{1}\right)  ,1\leq j,k\leq r\right\}
\label{subset}%
\end{equation}
generated by any basis of $\mathcal{H}^{1}$ and thus contains all possible
diagonal values of $f$ on products of the form $a^{\dagger}b^{\dagger}ba.$

The conditions eqs. $\left(  \ref{NormCon},\ref{PCond},\ref{GCond}\right)  $
and the lower bound in eq. $\left(  \ref{Qcond}\right)  $ are known from the
study of the Diagonal Representability Problem
\cite{DavidsonI,DavidsonII,DavidsonIII,WeinholdWilson,Erdahl78}. However,
these conditions were not caste in terms of existance of positive extensions,
which allows one to utilize general theorems from functional analysis that
lead to the consideration of bounds on the values of functionals on the
manifold of projectors contained in $\Lambda_{2}.$ In \cite{weiner2011} we
argued that all such projectors are of a particular form and they constitute a
homogeneous manifold of the unitary group $U\left(  \mathcal{H}^{1}\right)  .$
In appendix \ref{Projection} we give an improved proof of this assertion.

\section{Group Manifold and Coherent States}

The group manifold formed by the elements of the group $U\left(
\mathcal{H}^{1}\right)  $ can be parameterized by the $\left\{  \lambda
_{\varphi jk}\right\}  $ introduced in eq. $\left(  \ref{lamda}\right)  ,$
however a more useful parameterization for the study of SORDO's is the one
generated by employing two electron Thouless Coherent State (TCS) vectors,
which are parameterized by $\left\{  \left.  z_{ph}\in\mathbb{C}\right\vert
1\leq h\leq2,3\leq p\leq r\right\}  $ and defined by ( \cite{weiner}%
\cite{Orhn}\cite{Thouless}\cite{Peremilov} and appendix \ref{Thouless})
\begin{equation}
\left\vert \mathbf{z}\right\rangle =\left\vert \varphi_{1}\left(
\mathbf{z}\right)  \varphi_{2}\left(  \mathbf{z}\right)  \right\rangle
;\quad\mathbf{z\in}\mathbb{C}^{r\left(  r-2\right)  },
\end{equation}
where
\begin{equation}
\left\vert \mathbf{z}\right\rangle =\left\vert \varphi_{1}\left(
\mathbf{z}\right)  \varphi_{2}\left(  \mathbf{z}\right)  \right\rangle
=\exp\left\{  \sum\limits_{\substack{1\leq h\leq2\\3\leq p\leq r}}z_{ph}%
a_{p}^{\dagger}a_{h}\right\}  \left\vert \varphi_{1}\varphi_{2}\right\rangle
\label{CS}%
\end{equation}
and the transformed reference one particle basis elements $\left\{  \left.
\left\vert \varphi_{j}\left(  \mathbf{z}\right)  \right\rangle \right\vert
1\leq j\leq2\right\}  $ are given by
\begin{equation}
\left\vert \varphi_{j}\left(  \mathbf{z}\right)  \right\rangle =\exp\left\{
\sum\limits_{\substack{1\leq h\leq2\\3\leq p\leq r}}z_{ph}a_{p}^{\dagger}%
a_{h}\right\}  \left\vert \varphi_{j}\right\rangle ;1\leq j\leq2.
\end{equation}
The choice of the reference vector%
\begin{equation}
\left\vert \Phi_{2}\right\rangle =\left\vert \varphi_{1}\varphi_{2}%
\right\rangle
\end{equation}
is arbitrary. The manifold $\mathfrak{M}_{2}$ formed by the TCS's, which is
the set of two particle Independent Particle States (IPS's), is holomorpically
parameterized by the variables $\mathbf{z.}$

Coherent States have many significant properties \cite{Gilmore}
\cite{Peremilov} \cite{weiner} \cite{Orhn} \cite{Klauder} most deriving from
the fact that they form resolutions of projectors, in particular in this case
the projector $P_{\mathcal{H}^{2}}$ onto $\mathcal{H}^{2}$ i.e. (appendix
\ref{Thouless}\cite{Blaiz&Orland} )
\begin{align}
P_{\mathcal{H}^{2}}  &  =\int\exp\left\{  \sum\limits_{\substack{1\leq
h\leq2\\3\leq p\leq r}}z_{ph}a_{p}^{\dagger}a_{h}\right\}  \left\vert \Phi
_{2}\right\rangle \left\langle \Phi_{2}\right\vert \exp\left\{  \sum
\limits_{\substack{1\leq h\leq2\\3\leq p\leq r}}z_{ph}a_{p}^{\dagger}%
a_{h}\right\}  d\mu\left(  \mathbf{z}\right) \nonumber\\
&  =\int\left\vert \mathbf{z}\right\rangle \left\langle \mathbf{z}\right\vert
d\mu\left(  \mathbf{z}\right)  =\int\left\vert \mathbf{z}\right\rangle
\left\langle \mathbf{z}\right\vert \det\left(  \mathbb{I}_{2}+\mathbf{z}%
^{\dagger}\mathbf{z}\right)  ^{-\left(  r+1\right)  }\frac{\dim\left\{
\mathcal{H}^{2}\right\}  }{\operatorname*{vol}\left\{  \frac{U\left(
\mathcal{H}^{1}\right)  }{U\left(  \mathcal{H}_{p}^{1}\right)  }\right\}
}d\mathbf{\bar{z}}d\mathbf{z}, \label{CS Proj}%
\end{align}
i.e.
\begin{equation}
d\mu\left(  \mathbf{z}\right)  =\det\left(  \mathbb{I}_{2}+\mathbf{z}%
^{\dagger}\mathbf{z}\right)  ^{-\left(  r+1\right)  }\frac{\dim\left\{
\mathcal{H}^{2}\right\}  }{\operatorname*{vol}\left\{  \frac{U\left(
\mathcal{H}^{1}\right)  }{U\left(  \mathcal{H}_{p}^{1}\right)  }\right\}
}d\mathbf{\bar{z}}d\mathbf{z}%
\end{equation}
In terms of the measure $\tilde{\mu}$
\begin{align}
P_{\mathcal{H}^{2}}  &  =%
%TCIMACRO{\dint }%
%BeginExpansion
{\displaystyle\int}
%EndExpansion
\det\left(  \mathbb{I}_{2}+\mathbf{z}^{\dagger}\mathbf{z}\right)  ^{-1}%
\exp\left\{  \sum\limits_{\substack{1\leq h\leq2\\3\leq p\leq r}}z_{ph}%
a_{p}^{\dagger}a_{h}\right\}  \left\vert \Phi_{2}\right\rangle \left\langle
\Phi_{2}\right\vert \exp\left\{  \sum\limits_{\substack{1\leq h\leq2\\3\leq
p\leq r}}z_{ph}a_{p}^{\dagger}a_{h}\right\}  d\tilde{\mu}\left(
\mathbf{z}\right) \nonumber\\
&  =\int\left\vert \mathbf{\tilde{z}}\right\rangle \left\langle \mathbf{\tilde
{z}}\right\vert d\tilde{\mu}\left(  \mathbf{z}\right)  =\int\left\vert
\mathbf{\tilde{z}}\right\rangle \left\langle \mathbf{\tilde{z}}\right\vert
\det\left(  \mathbb{I}_{2}+\mathbf{z}^{\dagger}\mathbf{z}\right)  ^{-r}%
\frac{\dim\left\{  \mathcal{H}^{2}\right\}  }{\operatorname*{vol}\left\{
\frac{U\left(  \mathcal{H}^{1}\right)  }{U\left(  \mathcal{H}_{p}^{1}\right)
}\right\}  }d\mathbf{\bar{z}}d\mathbf{z},
\end{align}
i.e.
\begin{equation}
d\tilde{\mu}\left(  \mathbf{z}\right)  =\det\left(  \mathbb{I}_{2}%
+\mathbf{z}^{\dagger}\mathbf{z}\right)  ^{-r}\left(  2\pi i\right)
^{-2\left(  r-2\right)  }d\mathbf{\bar{z}}d\mathbf{z.}%
\end{equation}
\cite{Blaiz&Orland}. The measure $\mu$ is bounded \cite{Gilmore} and defined
by
\begin{equation}
\mu\left(  \Delta\right)  =\int_{\Delta\subseteq\mathbb{C}^{2\left(
r-2\right)  }}d\mu\left(  \mathbf{z}\right)  =\int_{\mathbb{C}^{2\left(
r-2\right)  }}\mathit{p}_{\Delta}\left(  \mathbf{z}\right)  d\mu\left(
\mathbf{z}\right)  =\int_{\mathbb{C}^{2\left(  r-2\right)  }}\mathit{p}%
_{\Delta}\left(  \mathbf{z}\right)  \det\left(  \mathbb{I}_{2}+\mathbf{z}%
^{\dagger}\mathbf{z}\right)  ^{-\left(  r+1\right)  }\left(  2\pi i\right)
^{-2\left(  r-2\right)  }d\mathbf{\bar{z}}d\mathbf{z,} \label{meas}%
\end{equation}
where
\begin{equation}
\mathit{p}_{\Delta}\left(  \mathbf{z}\right)  =\left\{
\begin{array}
[c]{c}%
1\text{ if }\mathbf{z}\text{ or }\mathbf{\bar{z}}\text{ }\in\Delta\\
0\text{ if }\mathbf{z}\text{ or }\mathbf{\bar{z}}\text{ }\notin\Delta
\end{array}
\right.  .
\end{equation}
One can reexpress eq. $\left(  \ref{meas}\right)  $ as
\begin{equation}
\mu\left(  \Delta\right)  =\int_{\Delta\subseteq\mathbb{C}^{2\left(
r-2\right)  }}d\mu\left(  \mathbf{z}\right)  =\int_{\mathbb{C}^{2\left(
r-2\right)  }}\mathit{p}_{\Delta}\left(  \mathbf{z}\right)  d\mu\left(
\mathbf{z}\right)  =\int_{\mathbb{C}^{2\left(  r-2\right)  }}\mathit{p}%
_{\Delta}\left(  \mathbf{z}\right)  e^{\eta}\frac{\dim\left\{  \mathcal{H}%
^{2}\right\}  }{\operatorname*{vol}\left\{  \frac{U\left(  \mathcal{H}%
^{1}\right)  }{U\left(  \mathcal{H}_{p}^{1}\right)  }\right\}  }%
d\mathbf{\bar{z}}d\mathbf{z,}%
\end{equation}
where
\begin{equation}
\eta=-\left(  r+1\right)  \operatorname*{Tr}\left\{  \mathbb{I}_{2}%
+\mathbf{z}^{\dagger}\mathbf{z}\right\}  u
\end{equation}
and%
\begin{equation}
\operatorname*{Tr}\left\{  \mathbb{I}_{2}+\mathbf{z}^{\dagger}\mathbf{z}%
\right\}  =\quad2+\mathbf{z}_{1}^{\dagger}\mathbf{z}_{1}+\mathbf{z}%
_{2}^{\dagger}\mathbf{z}_{2}.
\end{equation}
In detail
\begin{align}
\operatorname*{vol}\left\{  U\left(  \mathcal{H}^{1}\right)  \right\}   &
=2^{\frac{r}{2}}\pi^{\binom{r}{2}}\sqrt{r+1}%
%TCIMACRO{\dprod \limits_{k=1}^{r-1}}%
%BeginExpansion
{\displaystyle\prod\limits_{k=1}^{r-1}}
%EndExpansion
\left(  \frac{1}{k!}\right) \\
\operatorname*{vol}\left\{  U\left(  \mathcal{H}_{p}^{1}\right)  \right\}   &
=2^{\frac{r-2}{2}}\pi^{\binom{r-2}{2}}\sqrt{r-1}%
%TCIMACRO{\dprod \limits_{k=1}^{r-3}}%
%BeginExpansion
{\displaystyle\prod\limits_{k=1}^{r-3}}
%EndExpansion
\left(  \frac{1}{k!}\right) \\
\operatorname*{vol}\left\{  \frac{U\left(  \mathcal{H}^{1}\right)  }{U\left(
\mathcal{H}_{p}^{1}\right)  }\right\}   &  =\frac{2^{\frac{r}{2}}\pi
^{\frac{r\left(  r-1\right)  }{2}}\sqrt{r+1}%
%TCIMACRO{\dprod \limits_{k=1}^{r-1}}%
%BeginExpansion
{\displaystyle\prod\limits_{k=1}^{r-1}}
%EndExpansion
\left(  \frac{1}{k!}\right)  }{2^{\frac{r-2}{2}}\pi^{\frac{\left(  r-2\right)
\left(  r-3\right)  }{2}}\sqrt{r-1}%
%TCIMACRO{\dprod \limits_{k=1}^{r-3}}%
%BeginExpansion
{\displaystyle\prod\limits_{k=1}^{r-3}}
%EndExpansion
\left(  \frac{1}{k!}\right)  }=2\pi^{\frac{5r-6}{2}}\sqrt{\frac{r+1}{r-1}%
}\frac{1}{\left(  r-1\right)  \left(  r-2\right)  }\\
\frac{\dim\left\{  \mathcal{H}^{2}\right\}  }{\operatorname*{vol}\left\{
\frac{U\left(  \mathcal{H}^{1}\right)  }{U\left(  \mathcal{H}_{p}^{1}\right)
}\right\}  }  &  =\frac{r\left(  r-1\right)  \left(  r-1\right)  \left(
r-2\right)  }{4\pi^{\frac{5r-6}{2}}\sqrt{\frac{r+1}{r-1}}}\sim\frac{r^{3}}%
{\pi^{\frac{5r}{2}}}%
\end{align}
Clearly this measure is zero when $r=\infty$ on \emph{any }set $\Delta.$ If
$r=\infty$ then space $\mathbb{C}^{2\left(  r-2\right)  }$ becomes
$\mathbb{C}^{\infty}$ and has to be replaced by a function space, one can
observe this by first considering the coordinates to be values of a function
$z$ i.e.
\begin{equation}
\left(  z_{1},\cdots,z_{2\left(  r-2\right)  }\right)  \equiv\left(  z\left(
1\right)  ,\cdots,z\left(  2\left(  r-2\right)  \right)  \right)  ,
\end{equation}
where $z$ is complex valued map
\begin{equation}
z:\left\{  1,\cdots,2\left(  r-2\right)  \right\}  \rightarrow\mathbb{C}.
\label{finite z}%
\end{equation}
One can consider the values in eq. $\left(  \ref{finite z}\right)  $ as a
subset of the values of a function
\begin{equation}
z:\mathbb{N}\rightarrow\mathbb{C},
\end{equation}
which covers the case $r\rightarrow\infty.$ We can further extend this
reasoning to consider the countably infinite values $\left\{  \left.  z\left(
j\right)  \right\vert 1\leq j\leq\infty\right\}  $ to be a subset of values of
the function
\begin{equation}
z:\mathbb{R\rightarrow C},
\end{equation}
i.e. the restriction of $z$ to a countably infinte set points in $\mathbb{R}.$
(Wierstrauss theorem and $\mathbf{z}$ is replaced by a function $z)$. The
resolution of the identity then becomes
\begin{align}
P_{\mathcal{H}^{2}}  &  =\int\exp\left\{  \int\left\{  \Phi^{\dagger}\left(
f_{1}\right)  a_{1}+\Phi^{\dagger}\left(  f_{2}\right)  a_{2}\right\}
\right\}  \left\vert \Phi_{2}\right\rangle \left\langle \Phi_{2}\right\vert
\exp\left\{  \int\left\{  \Phi^{\dagger}\left(  f_{1}\right)  a_{1}%
+\Phi^{\dagger}\left(  f_{2}\right)  a_{2}\right\}  \right\}  d\sigma\left(
f\right) \nonumber\\
&  =\int\left\vert \mathbf{f}\right\rangle \left\langle \mathbf{f}\right\vert
d\sigma\left(  f\right)  =\int\left\vert z\right\rangle \left\langle
\mathbf{z}\right\vert \det\left(  \mathbb{I}_{2}+\mathbf{z}^{\dagger
}\mathbf{z}\right)  ^{-\left(  r+1\right)  }\left(  2\pi i\right)  ^{-2\left(
r-2\right)  }d\mathbf{\bar{z}}d\mathbf{z},
\end{align}


The coherent state definition eq. $\left(  \ref{CS}\right)  $ leads to
unnormalized CS's $\left\vert \mathbf{z}\right\rangle $ i.e.
\begin{equation}
\left\langle \mathbf{z\left\vert \mathbf{z}\right.  }\right\rangle
=\det\left(  \mathbb{I}_{2}+\mathbf{z}^{\dagger}\mathbf{z}\right)  ,
\end{equation}
which is taken into account in the measure $d\mu\left(  \mathbf{z}\right)  $
in eq. $\left(  \ref{Cs Proj}\right)  .$

The parameterization $\mathbf{z}$ is based on the coset decomposition of the
group $U\left(  r\right)  \equiv U\left(  \mathcal{H}^{1}\right)  $ given by
\begin{equation}
U\left(  r\right)  =\frac{U\left(  r\right)  }{U\left(  2\right)  \times
U\left(  r-2\right)  }\times\left(  U\left(  2\right)  \times U\left(
r-2\right)  \right)
\end{equation}
and the expression of the coset in terms of elements of $GL\left(
r,\mathbb{C}\right)  $ as
\begin{equation}
\frac{U\left(  r\right)  }{U\left(  2\right)  \times U\left(  r-2\right)
}=K_{+}K_{0}K_{-};\quad K_{+},K_{0},K_{-}\subset GL\left(  r,\mathbb{C}%
\right)  ,
\end{equation}
where we consider the Fermi-Fock realization of the unitary sub groups that
form the isotropy group of the reference state $\left\vert \varphi_{1}%
\varphi_{2}\right\rangle $
\begin{subequations}
\begin{align}
\mathcal{R}_{\mathcal{H}_{\mathcal{F}}}\left(  U\left(  2\right)  \right)   &
=\exp\left\{  i\sum\limits_{1\leq h,h^{\prime}\leq2}\beta_{hh^{\prime}}%
a_{h}^{\dagger}a_{h^{\prime}}\right\} \nonumber\\
&  =\exp\left\{  iN_{01}\left(  \beta\right)  \right\}  ;\beta_{hh^{\prime}%
}=\bar{\beta}_{h^{\prime}h}\in\mathbb{C},1\leq h\leq2\\
\mathcal{R}_{\mathcal{H}_{\mathcal{F}}}\left(  U\left(  r-2\right)  \right)
&  =\exp\left\{  i\sum\limits_{3\leq p,p^{\prime}\leq r}\beta_{pp^{\prime}%
}a_{p}^{\dagger}a_{p^{\prime}}\right\} \nonumber\\
&  =\exp\left\{  iN_{02}\right\}  ;\,\beta_{pp^{\prime}}=\bar{\beta
}_{p^{\prime}p}\in\mathbb{C},3\leq p\leq r
\end{align}
and the factors of the coset%
\end{subequations}
\begin{subequations}
\begin{align}
\mathcal{R}_{\mathcal{H}_{\mathcal{F}}}\left(  K_{+}\left(  \mathbf{z}\right)
\right)   &  =\exp\left\{  \sum\limits_{\substack{1\leq h\leq2\\3\leq p\leq
r}}z_{ph}a_{p}^{\dagger}a_{h}\right\}  =\exp\left\{  N_{+}\left(
\mathbf{z}\right)  \right\} \\
\mathcal{R}_{\mathcal{H}_{\mathcal{F}}}\left(  K_{0}\left(  \mathbf{z}\right)
\right)   &  =\exp\left\{  \sum\limits_{1\leq h,h^{\prime}\leq2}%
\alpha_{hh^{\prime}}\left(  \mathbf{z}\right)  a_{h}^{\dagger}a_{h^{\prime}%
}+\sum\limits_{3\leq p,p^{\prime}\leq r}\alpha_{pp^{\prime}}\left(
\mathbf{z}\right)  a_{p}^{\dagger}a_{p^{\prime}}\right\}  =\exp\left\{
N_{0}\left(  \mathbf{z}\right)  \right\}  ,\,\\
N_{0}\left(  \alpha\left(  \mathbf{z}\right)  \right)   &  =N_{01}\left(
\alpha\left(  \mathbf{z}\right)  \right)  +N_{02}\left(  \alpha\left(
\mathbf{z}\right)  \right)  ;\alpha_{hh^{\prime}}\left(  \mathbf{z}\right)
=\bar{\alpha}_{h^{\prime}h}\left(  \mathbf{z}\right)  ,\,\alpha_{pp^{\prime}%
}\left(  \mathbf{z}\right)  =\bar{\alpha}_{p^{\prime}p}\left(  \mathbf{z}%
\right)  \in\mathbb{C}\nonumber\\
\mathcal{R}_{\mathcal{H}_{\mathcal{F}}}\left(  K_{-}\left(  \mathbf{z}\right)
\right)   &  =\exp\left\{  -\sum\limits_{\substack{1\leq h\leq2\\3\leq p\leq
r}}\bar{z}_{ph}a_{h}^{\dagger}a_{p}\right\}  =\exp\left\{  N_{-}\left(
\mathbf{z}\right)  \right\}  .
\end{align}
but notice that the betas \emph{do not }depend on $\mathbf{z.}$ The group
elements belonging to the Fock representation $\mathcal{R}_{\mathcal{H}%
_{\mathcal{F}}}\left(  U\left(  r\right)  \right)  $ of $U\left(  r\right)  $
can be factored as
\end{subequations}
\begin{equation}
u\left(  \mathbf{\lambda}\left(  \mathbf{z,\alpha\left(  \mathbf{z}\right)
,\beta}\right)  \right)  =u\left(  \mathbf{z,\alpha\left(  \mathbf{z}\right)
,\beta}\right)  =c\left(  \mathbf{z}\right)  h\left(  \mathbf{\alpha}\left(
\mathbf{z}\right)  \right)  c\left(  -\mathbf{z}\right)  ^{\dagger}k\left(
\mathbf{\beta}\right)  ,
\end{equation}
where%
\begin{equation}
u\left(  \mathbf{\lambda}\left(  \mathbf{z,\alpha\left(  \mathbf{z}\right)
,\beta}\right)  \right)  =\exp\left\{  \sum\limits_{1\leq j,k\leq r}%
\lambda_{jk}a_{j}^{\dagger}a_{k}\right\}  ;\lambda_{jk}=-\bar{\lambda}_{kj}\,,
\end{equation}
\begin{equation}
u\left(  \mathbf{z,\alpha\left(  \mathbf{z}\right)  ,\beta}\right)  u\left(
\mathbf{z,\alpha\left(  \mathbf{z}\right)  ,\beta}\right)  ^{\dagger}=c\left(
\mathbf{z}\right)  h\left(  \mathbf{\alpha}\left(  \mathbf{z}\right)  \right)
c\left(  -\mathbf{z}\right)  ^{\dagger}k\left(  \mathbf{\beta}\right)
k\left(  \mathbf{\beta}\right)  ^{\dagger}c\left(  -\mathbf{z}\right)
h\left(  \mathbf{\alpha}\left(  \mathbf{z}\right)  \right)  c\left(
\mathbf{z}\right)  ^{\dagger}=I_{\mathcal{H}_{\mathcal{F}}}.
\end{equation}
and parameters $\mathbf{\alpha}\left(  \mathbf{z}\right)  $ are defined so
that $c\left(  \mathbf{z}\right)  h\left(  \mathbf{\alpha}\left(
\mathbf{z}\right)  \right)  c\left(  -\mathbf{z}\right)  ^{\dagger}$ is a
unitary operator \cite{Gilmore}.

The set of vectors
\begin{equation}
\left\{  \left.  u\left(  \mathbf{z,\alpha\left(  \mathbf{z}\right)  ,\beta
}\right)  \left\vert \varphi_{j}\right\rangle \right\vert 1\leq j\leq
r\right\}
\end{equation}
is always a conb as the factors $h\left(  \mathbf{\alpha}\left(
\mathbf{z}\right)  \right)  c\left(  -\mathbf{z}\right)  ^{\dagger}k\left(
\mathbf{\beta}\right)  $ \emph{do} effect $\left\{  \left.  \left\vert
\varphi_{j}\right\rangle \right\vert 3\leq j\leq r\right\}  ,$ in contrast to
the fact that $c\left(  -\mathbf{z}\right)  ^{\dagger}k\left(  \mathbf{\beta
}\right)  $ does not effect $\left\vert \varphi_{1}\varphi_{2}\right\rangle .$
The group $U\left(  2\right)  \times U\left(  r-2\right)  $ is the isotropy
group of the vector state $\left\vert \varphi_{1}\varphi_{2}\right\rangle ,$
the two particle state operator $\left\vert \varphi_{1}\varphi_{2}%
\right\rangle \left\langle \varphi_{1}\varphi_{2}\right\vert $ and its second
quantized image $a_{1}^{\dagger}a_{2}^{\dagger}a_{2}a_{1}$ i.e.
\begin{subequations}
\begin{align}
\exp\left\{  iN_{0}\left(  \beta\right)  \right\}  \left\vert \varphi
_{1}\varphi_{2}\right\rangle  &  =\left\vert \varphi_{1}\varphi_{2}%
\right\rangle \\
\exp\left\{  iN_{0}\left(  \beta\right)  \right\}  \left\vert \varphi
_{1}\varphi_{2}\right\rangle \left\langle \varphi_{1}\varphi_{2}\right\vert
\exp\left\{  -iN_{0}\left(  \beta\right)  \right\}   &  =\left\vert
\varphi_{1}\varphi_{2}\right\rangle \left\langle \varphi_{1}\varphi
_{2}\right\vert \\
\exp\left\{  iN_{0}\left(  \beta\right)  \right\}  a_{1}^{\dagger}%
a_{2}^{\dagger}a_{2}a_{1}\exp\left\{  -iN_{0}\left(  \beta\right)  \right\}
&  =a_{1}^{\dagger}a_{2}^{\dagger}a_{2}a_{1},
\end{align}
where
\end{subequations}
\begin{equation}
N_{0}\left(  \beta\right)  =N_{01}\left(  \beta\right)  +N_{02}\left(
\beta\right)  .
\end{equation}
\qquad\qquad\qquad

The action of the group $U\left(  \mathcal{H}^{1}\right)  $ on a reference
projector $a_{1}^{\dagger}a_{2}^{\dagger}a_{2}a_{1}$ generates values of $f$
on all projectors of the form $a^{\dagger}b^{\dagger}ba$ (see eqs. $\left(
\ref{matel},\ref{allmat},\ref{subset}\right)  $ through
\begin{equation}
f\left(  u^{\dagger}a_{1}^{\dagger}a_{2}^{\dagger}a_{2}a_{1}u\right)  ,
\end{equation}
i.e. all diagonal values, can be expressed by utilizing the TCS
parameterization as
\begin{equation}
f\left(  \mathbf{\bar{z},z}\right)  =f\left(  c\left(  \mathbf{z}\right)
^{\dagger}a_{1}^{\dagger}a_{2}^{\dagger}a_{2}a_{1}c\left(  \mathbf{z}\right)
\right)  ,
\end{equation}
showing that these values are generated by biholomorphic real valued functions
on $\mathbb{C}^{2\left(  r-2\right)  }.$ The N-Representability Conditions
listed in section \ref{NREP_Con} can all be expressed in terms such functions
and in the next section we study them in detail.

\section{Thouless P and Q symbols \label{P and Q Symbols}}

\subsection{Q Symbols}

The normalized Q \cite{Lieb} (also known as covariant the \cite{Berizin} or
Wick \cite{Peremelov}) symbol, $\tilde{Q}_{A},$ of an operator $A$ that acts
on $\mathcal{H}^{2}$ (and could be unbounded) is a biholomorphic function
defined as
\begin{equation}
\tilde{Q}_{A}\left(  \mathbf{\bar{z}},\mathbf{z}\right)  =\tilde{Q}_{A}\left(
\mathbf{z}\right)  =\frac{\left\langle \mathbf{z}\left\vert A\mathbf{z}%
\right.  \right\rangle }{\left\langle \mathbf{z}\left\vert \mathbf{z}\right.
\right\rangle }=\frac{Q_{A}\left(  \mathbf{z}\right)  }{\left\langle
\mathbf{z}\left\vert \mathbf{z}\right.  \right\rangle }%
\end{equation}
and the unnormalized symbol as
\begin{equation}
Q_{A}\left(  \mathbf{z}\right)  =\operatorname*{Tr}\left\{  A\left\vert
\mathbf{z}\right\rangle \left\langle \mathbf{z}\right\vert \right\}  ,
\end{equation}
giving
\begin{equation}
Q_{A}\left(  \mathbf{z}\right)  =\tilde{Q}_{A}\left(  \mathbf{z}\right)
\left\langle \mathbf{z}\left\vert \mathbf{z}\right.  \right\rangle ,
\end{equation}
where the normalization kernel $\left\langle \mathbf{z}\left\vert
\mathbf{z}\right.  \right\rangle $ is a special case of the overlap kernel
\begin{equation}
\left\langle \mathbf{z\left\vert \mathbf{z}^{\prime}\right.  }\right\rangle
=\det\left(  \mathbb{I}_{2}+\mathbf{z}^{\dagger}\mathbf{z}^{\prime}\right)  .
\end{equation}


If the operator $A$ is bounded and self adjoint then its Q-symbol, $\left(
\text{which is real valued}\right)  ,$ is also $\tilde{\mu}$ essentially
bounded $\left(  \text{appendix \ref{Banach}}\right)  $ as
\begin{equation}
\left\Vert A\right\Vert _{\mathcal{B}\left(  \mathcal{H}^{2}\right)  }%
=\sup_{\left\vert \Phi\right\rangle \in\mathcal{H}^{2}}\left\{  \left\vert
\frac{\left\langle \left.  \Phi\right\vert A\Phi\right\rangle }{\left\langle
\left.  \Phi\right\vert \Phi\right\rangle }\right\vert \right\}  \geq
\sup_{\left\vert \mathbf{z}\right\rangle \in\mathcal{H}^{2}}\left\{
\left\vert \frac{\left\langle \mathbf{z}A\mathbf{z}\right\rangle
}{\left\langle \left.  \mathbf{z}\right\vert \mathbf{z}\right\rangle
}\right\vert \right\}  =\sup_{\mathbf{z}\in\mathbb{C}^{2\left(  r-2\right)  }%
}\left\{  \left\vert \tilde{Q}_{A}\left(  \mathbf{\bar{z}},\mathbf{z}\right)
\right\vert \right\}  \geq\tilde{\mu}-\operatorname*{ess}\sup_{\mathbf{z}%
\in\mathbb{C}^{2\left(  r-2\right)  }}\left\{  \left\vert \tilde{Q}_{A}\left(
\mathbf{\bar{z}},\mathbf{z}\right)  \right\vert \right\}
\end{equation}
and the space of Q-symbols of such operators naturally form the real Banach
space $L_{\infty}\left(  \mathbb{C}^{2\left(  r-2\right)  },\tilde{\mu
}\right)  $ with the norm $\left\Vert {}\right\Vert _{L_{\infty}\left(
\mathbb{C}^{2\left(  r-2\right)  },\tilde{\mu}\right)  }$ defined as
\begin{equation}
\left\Vert \tilde{Q}\right\Vert _{L_{\infty}\left(  \mathbb{C}^{2\left(
r-2\right)  },\tilde{\mu}\right)  }=\tilde{\mu}-\operatorname*{ess}%
\sup_{\mathbf{z}\in\mathbb{C}^{2\left(  r-2\right)  }}\left\{  \left\vert
\tilde{Q}\left(  \mathbf{\bar{z}},\mathbf{z}\right)  \right\vert \right\}  .
\end{equation}
The measures $\mu$ and $\tilde{\mu}$ are equivalent $\left(  \text{appendix
\ref{Banach}}\right)  $ thus
\begin{equation}
\left\Vert \tilde{Q}\right\Vert _{L_{\infty}\left(  \mathbb{C}^{2\left(
r-2\right)  },\mu\right)  }=\mu-\operatorname*{ess}\sup_{\mathbf{z}%
\in\mathbb{C}^{2\left(  r-2\right)  }}\left\{  \left\vert \tilde{Q}\left(
\mathbf{\bar{z}},\mathbf{z}\right)  \right\vert \right\}
\end{equation}
and
\begin{equation}
\left\Vert Q\right\Vert _{L_{\infty}\left(  \mathbb{C}^{2\left(  r-2\right)
},\mu\right)  }=\mu-\operatorname*{ess}\sup_{\mathbf{z}\in\mathbb{C}^{2\left(
r-2\right)  }}\left\{  \left\vert Q\left(  \mathbf{\bar{z}},\mathbf{z}\right)
\right\vert \right\}  =\tilde{\mu}-\operatorname*{ess}\sup_{\mathbf{z}%
\in\mathbb{C}^{2\left(  r-2\right)  }}\left\{  \left\vert Q\left(
\mathbf{\bar{z}},\mathbf{z}\right)  \right\vert \right\}  =\left\Vert
Q\right\Vert _{L_{\infty}\left(  \mathbb{C}^{2\left(  r-2\right)  },\tilde
{\mu}\right)  }.
\end{equation}
However the preduals $L_{1}\left(  \mathbb{C}^{2\left(  r-2\right)  }%
,\tilde{\mu}\right)  $ and $L_{1}\left(  \mathbb{C}^{2\left(  r-2\right)
},\mu\right)  $ are not the same as it is to see that $L_{1}\left(
\mathbb{C}^{2\left(  r-2\right)  },\tilde{\mu}\right)  \subset L_{1}\left(
\mathbb{C}^{2\left(  r-2\right)  },\mu\right)  $ as $\left\Vert f\right\Vert
_{L_{1}\left(  \mathbb{C}^{2\left(  r-2\right)  },\tilde{\mu}\right)
}>\left\Vert f\right\Vert _{L_{1}\left(  \mathbb{C}^{2\left(  r-2\right)
},\mu\right)  }$, where
\begin{equation}
\left\Vert f\right\Vert _{L_{1}\left(  \mathbb{C}^{2\left(  r-2\right)
},\tilde{\mu}\right)  }=\int\left\vert f\left(  \mathbf{z}\right)  \right\vert
d\tilde{\mu}\left(  \mathbf{z}\right)  =\int\left\vert f\left(  \mathbf{z}%
\right)  \right\vert \det\left(  \mathbb{I}_{2}+\mathbf{z}^{\dagger}%
\mathbf{z}\right)  d\mu\left(  \mathbf{z}\right)
\end{equation}
and
\begin{equation}
\left\Vert f\right\Vert _{L_{1}\left(  \mathbb{C}^{2\left(  r-2\right)  }%
,\mu\right)  }=\int\left\vert f\left(  \mathbf{z}\right)  \right\vert
d\mu\left(  \mathbf{z}\right)  .
\end{equation}
The trace of an operator $A$ can be expressed as
\begin{equation}
\operatorname*{Tr}\left\{  A\right\}  =\int A\left\vert \mathbf{\tilde{z}%
}\right\rangle \left\langle \mathbf{\tilde{z}}\right\vert d\tilde{\mu}\left(
\mathbf{z}\right)  =\int\frac{\left\langle \mathbf{z}\left\vert A\mathbf{z}%
\right.  \right\rangle }{\left\langle \mathbf{z}\left\vert \mathbf{z}\right.
\right\rangle }d\tilde{\mu}\left(  \mathbf{z}\right)  =\int\frac{Q_{A}\left(
\mathbf{z}\right)  }{\left\langle \mathbf{z}\left\vert \mathbf{z}\right.
\right\rangle }d\tilde{\mu}\left(  \mathbf{z}\right)  =\int Q_{A}\left(
\mathbf{z}\right)  d\mu\left(  \mathbf{z}\right)
\end{equation}
and if $A\geq0$ then
\begin{equation}
\operatorname*{Tr}\left\{  A\right\}  =\left\Vert \tilde{Q}_{A}\right\Vert
_{L_{1}\left(  \mathbb{C}^{2\left(  r-2\right)  },\tilde{\mu}\right)
}=\left\Vert Q_{A}\right\Vert _{L_{1}\left(  \mathbb{C}^{2\left(  r-2\right)
},\mu\right)  }.
\end{equation}


\subsection{P Symbols}

\qquad The P (also known as the contravariant \cite{Berizin} or AntiWick
\cite{Peremelov}) symbol of an operator $A$ is formally defined as
\begin{equation}
A=\int P_{A}\left(  \mathbf{z}\right)  \left\vert \mathbf{z}\right\rangle
\left\langle \mathbf{z}\right\vert d\mu\left(  \mathbf{z}\right)  =\int
P_{A}\left(  \mathbf{z}\right)  \left\vert \mathbf{\tilde{z}}\right\rangle
\left\langle \mathbf{\tilde{z}}\right\vert d\tilde{\mu}\left(  \mathbf{z}%
\right)  ,
\end{equation}
i.e. the P-symbols with respect to normalized and unnormilized CS's can be
chosen to be the same, unlike the Q-symbols. There are informative discussions
of P-symbols in \cite{Tajourdhi} and \cite{Klauder}. The P-symbol of a bounded
operator might not be bounded i.e. not exist (even in finite dimensions) but
if $P_{A}\left(  \mathbf{z}\right)  $ does exists then
\begin{equation}
\operatorname*{Tr}\left\{  A\right\}  =\operatorname*{Tr}\left\{
AP_{\mathcal{H}^{2}}\right\}  =\operatorname*{Tr}\left\{  A\int\left\vert
\mathbf{z}\right\rangle \left\langle \mathbf{z}\right\vert d\mu\left(
\mathbf{z}\right)  \right\}  =\int P_{A}\left(  \mathbf{z}\right)
\left\langle \mathbf{z\left\vert \mathbf{z}\right.  }\right\rangle d\mu\left(
\mathbf{z}\right)  ,
\end{equation}
which can be expressed as
\begin{equation}
\operatorname*{Tr}\left\{  A\right\}  =\left\Vert P_{A}\right\Vert
_{L_{1}\left(  \mathbb{C}^{2\left(  r-2\right)  },\tilde{\mu}\right)  }%
\neq\left\Vert P_{A}\right\Vert _{L_{1}\left(  \mathbb{C}^{2\left(
r-2\right)  },\mu\right)  }.
\end{equation}
Thus $A$ is trace class if $P_{A}\in L_{1}\left(  \mathbb{C}^{2\left(
r-2\right)  },\tilde{\mu}\right)  $. The space of P-symbols is $L_{1}\left(
\mathbb{C}^{2\left(  r-2\right)  },\mu\right)  .$

\subsection{Berezin Transformation}

\bigskip\qquad This leads to the Berezin Transformation \cite{BerezinTrans}
between P-symbols and Q-symbols given by
\begin{align}
\left[  \mathfrak{Be}\left(  P_{A}\right)  \right]  \left(  \mathbf{z}%
^{\prime}\right)   &  =Q_{A}\left(  \mathbf{z}^{\prime}\right)  =\frac
{1}{\left\langle \mathbf{z}^{\prime}\mathbf{\left\vert \mathbf{z}^{\prime
}\right.  }\right\rangle }\int_{\mathbb{C}^{r\left(  r-2\right)  }}%
P_{A}\left(  \mathbf{z}\right)  \left\vert \left\langle \mathbf{z\left\vert
\mathbf{z}^{\prime}\right.  }\right\rangle \right\vert ^{2}d\mu\left(
\mathbf{z}\right) \nonumber\\
&  =\int\mathfrak{Be}\left(  \mathbf{z,z}^{\prime}\right)  P_{A}\left(
\mathbf{z}^{\prime}\right)  d\mu\left(  \mathbf{z}^{\prime}\right)  ,
\end{align}
where the Berezin integral kernel is
\begin{equation}
\mathfrak{Be}\left(  \mathbf{z,z}^{\prime}\right)  =\left\vert \left\langle
\mathbf{z\left\vert \mathbf{z}^{\prime}\right.  }\right\rangle \right\vert
^{2}.
\end{equation}
The inverse Berezin Transformation is formally given by
\begin{equation}
\mathfrak{Be}^{-1}\left(  Q_{A}\right)  \left(  \mathbf{z}\right)
=P_{A}\left(  \mathbf{z}\right)  =\int\mathfrak{Be}^{-1}\left(  \mathbf{z,z}%
^{\prime}\right)  Q_{A}\left(  \mathbf{z}^{\prime}\right)  d\mu\left(
\mathbf{z}^{\prime}\right)  .
\end{equation}


\subsection{Bases of P and Q Symbols}

One can define a basis of Q-symbols, that in the finite dimensional case spans
the space of symbols, by
\begin{align}
Q_{j_{1}j_{2}k_{1}k_{2}}\left(  \mathbf{z}\right)   &  =\operatorname*{Tr}%
\left\{  \left\vert \mathbf{z}\right\rangle \left\langle \mathbf{z}\right\vert
\left\vert \varphi_{j_{1}}\varphi_{j_{2}}\right\rangle \left\langle
\varphi_{k_{1}}\varphi_{k_{2}}\right\vert \right\} \nonumber\\
&  =\left\langle \varphi_{k_{1}}\varphi_{k_{2}}\left\vert \left\vert
\mathbf{z}\right\rangle \left\langle \mathbf{z}\right\vert \varphi_{j_{1}%
}\varphi_{j_{2}}\right.  \right\rangle \nonumber\\
&  =\gamma_{2}\left(  \mathbf{z}\right)  _{j_{1}j_{2}k_{1}k_{2}}\left\langle
\mathbf{z\left\vert \mathbf{z}\right.  }\right\rangle ,
\end{align}
where $\gamma_{2}\left(  \mathbf{z}\right)  $ is the normalized density matrix
of the two particle state operator $\frac{\left\vert \mathbf{z}\right\rangle
\left\langle \mathbf{z}\right\vert }{\left\langle \mathbf{z\left\vert
\mathbf{z}\right.  }\right\rangle }$ with respect to the orthonormal basis
$\left\{  \left.  \left\vert \varphi_{j_{1}}\varphi_{j_{2}}\right\rangle
\right\vert 1\leq j_{1}<j_{2}\leq r\right\}  $ of $\mathcal{H}^{2}.$
\begin{align}
Q_{j_{1}j_{2}k_{1}k_{2}}\left(  \mathbf{z}\right)   &  =\operatorname*{Tr}%
\left\{  \left\vert \mathbf{z}\right\rangle \left\langle \mathbf{z}\right\vert
\left\vert \varphi_{j_{1}}\varphi_{j_{2}}\right\rangle \left\langle
\varphi_{k_{1}}\varphi_{k_{2}}\right\vert \right\}  =\nonumber\\
&  =\operatorname*{Tr}\left\{  u\left(  \mathbf{z}\right)  \left\vert
\varphi_{1}\varphi_{2}\right\rangle \left\langle \varphi_{1}\varphi
_{2}\right\vert u\left(  \mathbf{z}\right)  ^{\dagger}\left\vert
\varphi_{j_{1}}\varphi_{j_{2}}\right\rangle \left\langle \varphi_{k_{1}%
}\varphi_{k_{2}}\right\vert \right\} \\
&  =\operatorname*{Tr}\left\{  \left\vert \varphi_{1}\varphi_{2}\right\rangle
\left\langle \varphi_{1}\varphi_{2}\right\vert u\left(  \mathbf{z}\right)
^{\dagger}\left\vert \varphi_{j_{1}}\varphi_{j_{2}}\right\rangle \left\langle
\varphi_{k_{1}}\varphi_{k_{2}}\right\vert u\left(  \mathbf{z}\right)  \right\}
\\
&  =\left\langle \varphi_{1}\varphi_{2}\right\vert u\left(  \mathbf{z}\right)
^{\dagger}\left\vert \varphi_{j_{1}}\varphi_{j_{2}}\right\rangle \left\langle
\varphi_{k_{1}}\varphi_{k_{2}}\left\vert u\left(  \mathbf{z}\right)
\varphi_{1}\varphi_{2}\right.  \right\rangle \\
&  =\left\langle \varphi_{1}\varphi_{2}\right\vert u\left(  \mathbf{z}\right)
^{\dagger}\left\vert \varphi_{j_{1}}\varphi_{j_{2}}\right\rangle \left\langle
\varphi_{k_{1}}\varphi_{k_{2}}\left\vert u\left(  \mathbf{z}\right)
\varphi_{1}\varphi_{2}\right.  \right\rangle \\
&  =\left\langle \varphi_{1}\varphi_{2}\left\vert \widehat{u\left(
\mathbf{z}\right)  }^{\dagger}\left(  \left\vert \varphi_{j_{1}}\varphi
_{j_{2}}\right\rangle \left\langle \varphi_{k_{1}}\varphi_{k_{2}}\right\vert
\right)  \varphi_{1}\varphi_{2}\right.  \right\rangle \\
&  =\left\langle \varphi_{k_{1}}\varphi_{k_{2}}\left\vert \left\vert
\mathbf{z}\right\rangle \left\langle \mathbf{z}\right\vert \varphi_{j_{1}%
}\varphi_{j_{2}}\right.  \right\rangle \nonumber\\
&  =\gamma_{2}\left(  \mathbf{z}\right)  _{j_{1}j_{2}k_{1}k_{2}}\left\langle
\mathbf{z\left\vert \mathbf{z}\right.  }\right\rangle ,
\end{align}
basis is linearly independent and any Q-symbol $Q_{A}\left(  \mathbf{z}%
\right)  $ can be expanded as
\begin{equation}
Q_{A}\left(  \mathbf{z}\right)  =Q_{A}\left(  \mathbf{\bar{z},z}\right)
=\sum_{\substack{1\leq l_{1}<l_{2}\leq r,\\1\leq m_{1}<m_{2}\leq r}%
}Q_{j_{1}j_{2}k_{1}k_{2}}\left(  \mathbf{z}\right)  A_{j_{1}j_{2}k_{1}k_{2}}%
\end{equation}
and
\begin{equation}
A_{j_{1}j_{2}k_{1}k_{2}}=\left\{  Q_{A}\left\vert P_{j_{1}j_{2}k_{1}k_{2}%
}\right.  \right\}  .
\end{equation}
A basis of P-symbols, that in the finite dimensional case also spans the space
of symbols, can be constructed from the Q-symbol basis by using the inverse
Berizin transformation i.e.
\begin{equation}
P_{j_{1}j_{2}k_{1}k_{2}}\left(  \mathbf{z}\right)  =\int\mathfrak{Be}%
^{-1}\left(  \mathbf{z,z}^{\prime}\right)  Q_{j_{1}j_{2}k_{1}k_{2}}\left(
\mathbf{z}^{\prime}\right)  d\mu\left(  \mathbf{z}^{\prime}\right)  ,
\end{equation}
and any P-symbol $P_{A}\left(  \mathbf{z}\right)  $ can be expanded as
\begin{equation}
P_{A}\left(  \mathbf{z}\right)  =P_{A}\left(  \mathbf{\bar{z},z}\right)
=\sum_{\substack{1\leq l_{1}<l_{2}\leq r,\\1\leq m_{1}<m_{2}\leq r}%
}P_{j_{1}j_{2}k_{1}k_{2}}\left(  \mathbf{z}\right)  A_{j_{1}j_{2}k_{1}k_{2}}%
\end{equation}
and
\begin{equation}
A_{j_{1}j_{2}k_{1}k_{2}}=\left\{  Q_{j_{1}j_{2}k_{1}k_{2}}\left\vert
P_{A}\right.  \right\}  .
\end{equation}
The integral kernel of the inverse Berezin transform is given by \cite{Weiner}%
\begin{align}
\mathfrak{Be}^{-1}\left(  \mathbf{z,z}^{\prime}\right)   &  =\sum
_{_{\substack{1\leq l_{1}<l_{2}\leq r,\\1\leq m_{1}<m_{2}\leq r,\\1\leq
j_{1}<j_{2}\leq r,\\1\leq k_{1}<k_{2}\leq r}}}Q_{l_{1}l_{2}m_{1}m_{2}}\left(
\mathbf{z}\right)  \left(  \Delta^{-2}\right)  _{l_{1}l_{2}m_{1}m_{2}%
j_{1}j_{2}k_{1}k_{2}}Q_{j_{1}j_{2}k_{1}k_{2}}\left(  \mathbf{z}^{\prime
}\right) \nonumber\\
&  =\left\langle \mathbf{z\left\vert \mathbf{z}\right.  }\right\rangle
^{2}\sum_{\substack{1\leq l_{1}<l_{2}\leq r,\\1\leq m_{1}<m_{2}\leq r,\\1\leq
j_{1}<j_{2}\leq r,\\1\leq k_{1}<k_{2}\leq r}}\gamma_{2}\left(  \mathbf{z}%
\right)  _{l_{1}l_{2}m_{1}m_{2}}\left(  \Delta^{-2}\right)  _{l_{1}l_{2}%
m_{1}m_{2}j_{1}j_{2}k_{1}k_{2}}\gamma_{2}\left(  \mathbf{z}^{\prime}\right)
_{j_{1}j_{2}k_{1}k_{2}}, \label{InvBer}%
\end{align}
where
\begin{align}
\Delta_{l_{1}l_{2}m_{1}m_{2}j_{1}j_{2}k_{1}k_{2}}  &  =\int\overline
{\operatorname*{Tr}\left\{  \left\vert \mathbf{z}\right\rangle \left\langle
\mathbf{z}\right\vert \left\vert \varphi_{l_{1}}\varphi_{l_{2}}\right\rangle
\left\langle \varphi_{m_{1}}\varphi_{m_{2}}\right\vert \right\}
}\operatorname*{Tr}\left\{  \left\vert \mathbf{z}\right\rangle \left\langle
\mathbf{z}\right\vert \left\vert \varphi_{j_{1}}\varphi_{j_{2}}\right\rangle
\left\langle \varphi_{k_{1}}\varphi_{k_{2}}\right\vert \right\}  d\mu\left(
\mathbf{z}\right) \nonumber\\
&  =\int\overline{\gamma_{2}\left(  \mathbf{z}\right)  }_{m_{1}m_{2}l_{1}%
l_{2}}\gamma_{2}\left(  \mathbf{z}\right)  _{j_{1}j_{2}k_{1}k_{2}}\left\langle
\mathbf{z\left\vert \mathbf{z}\right.  }\right\rangle ^{2}d\mu\left(
\mathbf{z}\right)  .
\end{align}
and can be expanded, if $\Delta^{-1}$ exists, pointwise as
\begin{align}
P_{j_{1}j_{2}k_{1}k_{2}}\left(  \mathbf{z}\right)   &  =\sum_{1\leq
p_{1}<p_{2}\leq r,1\leq q_{1}<q_{2}\leq r}Q_{p_{1}p_{2}q_{1}q_{2}}\left(
\mathbf{z}\right)  \Delta_{p_{1}p_{2}q_{1}q_{2}j_{1}j_{2}k_{1}k_{2}}%
^{-1}\nonumber\\
&  =\sum_{\substack{1\leq l_{1}<l_{2}\leq r,\\1\leq m_{1}<m_{2}\leq
r}}\operatorname*{Tr}\left\{  \left\vert \mathbf{z}\right\rangle \left\langle
\mathbf{z}\right\vert \left\vert \varphi_{l_{1}}\varphi_{l_{2}}\right\rangle
\left\langle \varphi_{m_{1}}\varphi_{m_{2}}\right\vert \right\}  \left(
\Delta^{-1}\right)  _{l_{1}l_{2}m_{1}m_{2}j_{1}j_{2}k_{1}k_{2}}\nonumber\\
&  =\left\langle \mathbf{z\left\vert \mathbf{z}\right.  }\right\rangle
^{2}\sum_{\substack{1\leq l_{1}<l_{2}\leq r,\\1\leq m_{1}<m_{2}\leq r}%
}\gamma_{2}\left(  \mathbf{z}\right)  _{m_{1}m_{2}l_{1}l_{2}}\left(
\Delta^{-1}\right)  _{l_{1}l_{2}m_{1}m_{2}j_{1}j_{2}k_{1}k_{2}}.
\end{align}
In finite dimensions one can introduce a scaler product by
\begin{equation}
\left\{  \left.  Q\right\vert P\right\}  =\int\overline{Q\left(
\mathbf{z}\right)  }P\left(  \mathbf{z}\right)  d\mu\left(  \mathbf{z}\right)
.
\end{equation}
The Q and P-symbols $\left\{  Q_{j_{1}j_{2}k_{1}k_{2}}\right\}  ,$ $\left\{
P_{j_{1}j_{2}k_{1}k_{2}}\right\}  $ form biorthogonal bases (Appendix ) i.e.
\begin{equation}
\left\{  \left.  Q_{j_{1}j_{2}k_{1}k_{2}}\right\vert P_{l_{1}l_{2}m_{1}m_{2}%
}\right\}  =\delta_{\left(  j_{1}j_{2}k_{1}k_{2}\right)  \left(  l_{1}%
l_{2}m_{1}m_{2}\right)  }%
\end{equation}
and%
\begin{equation}
\left\{  \left.  Q_{j_{1}j_{2}k_{1}k_{2}}\right\vert P_{A}\right\}  =\left\{
\left.  Q_{A}\right\vert P_{j_{1}j_{2}k_{1}k_{2}}\right\}  =A_{j_{1}j_{2}%
k_{1}k_{2}}.
\end{equation}
The biorthogonality is a result of the relationships
\begin{align}
\delta_{\left(  l_{1}l_{2}\right)  \left(  j_{1}j_{2}\right)  }\delta_{\left(
m_{1}m_{2}\right)  \left(  k_{1}k_{2}\right)  }  &  =\int\overline
{Q_{l_{1}l_{2}m_{1}m_{2}}\left(  \mathbf{z}\right)  }P_{j_{1}j_{2}k_{1}k_{2}%
}\left(  \mathbf{z}\right)  d\mu\left(  \mathbf{z}\right) \nonumber\\
&  =\int\overline{\operatorname*{Tr}\left\{  \left\vert \mathbf{z}%
\right\rangle \left\langle \mathbf{z}\right\vert \left\vert \varphi_{l_{1}%
}\varphi_{l_{2}}\right\rangle \left\langle \varphi_{m_{1}}\varphi_{m_{2}%
}\right\vert \right\}  }P_{j_{1}j_{2}k_{1}k_{2}}\left(  \mathbf{z}\right)
d\mu\left(  \mathbf{z}\right) \nonumber\\
&  =\operatorname*{Tr}\left\{  \left(  \int\left\vert \mathbf{z}\right\rangle
\left\langle \mathbf{z}\right\vert P_{j_{1}j_{2}k_{1}k_{2}}\left(
\mathbf{z}\right)  d\mu\left(  \mathbf{z}\right)  \right)  \left\vert
\varphi_{m_{1}}\varphi_{m_{2}}\right\rangle \left\langle \varphi_{l_{1}%
}\varphi_{l_{2}}\right\vert \right\}  ,
\end{align}
that define P-symbols, \emph{which shows that the integral defines the
following operator identity}
\begin{equation}
\left\vert \varphi_{j_{1}}\varphi_{j_{2}}\right\rangle \left\langle
\varphi_{k_{1}}\varphi_{k_{2}}\right\vert \equiv\int P_{j_{1}j_{2}k_{1}k_{2}%
}\left(  \mathbf{z}\right)  \left\vert \mathbf{z}\right\rangle \left\langle
\mathbf{z}\right\vert d\mu\left(  \mathbf{z}\right)
\end{equation}
\emph{and} \emph{that }$P_{j_{1}j_{2}k_{1}k_{2}}\left(  \mathbf{z}\right)  $
\emph{is the P}$\emph{-symbol}$ \emph{of} $\left\vert \varphi_{j_{1}}%
\varphi_{j_{2}}\right\rangle \left\langle \varphi_{k_{1}}\varphi_{k_{2}%
}\right\vert $, i.e.
\begin{equation}
P_{j_{1}j_{2}k_{1}k_{2}}\left(  \mathbf{z}\right)  =P_{\left\vert
\varphi_{j_{1}}\varphi_{j_{2}}\right\rangle \left\langle \varphi_{k_{1}%
}\varphi_{k_{2}}\right\vert }\left(  \mathbf{z}\right)  ,
\end{equation}
note that P-symbols are not unique.

In finite dimensions the P-symbols can be expanded in terms of the Q-symbols
\begin{equation}
P_{j_{1}j_{2}k_{1}k_{2}}\left(  \mathbf{z}\right)  =\sum_{1\leq p_{1}%
<p_{2}\leq r,1\leq q_{1}<q_{2}\leq r}Q_{p_{1}p_{2}q_{1}q_{2}}\left(
\mathbf{z}\right)  \Delta_{p_{1}p_{2}q_{1}q_{2}j_{1}j_{2}k_{1}k_{2}}^{-1}
\label{P in Q}%
\end{equation}
and the P-symbols can be expanded in terms of the Q-symbols%
\begin{equation}
Q_{j_{1}j_{2}k_{1}k_{2}}\left(  \mathbf{z}\right)  =\sum_{1\leq p_{1}%
<p_{2}\leq r,1\leq q_{1}<q_{2}\leq r}P_{p_{1}p_{2}q_{1}q_{2}}\left(
\mathbf{z}\right)  \Delta_{p_{1}p_{2}q_{1}q_{2}j_{1}j_{2}k_{1}k_{2}}
\label{Q in P}%
\end{equation}
giving in particular%
\begin{equation}
P_{A}\left(  \mathbf{z}\right)  =P_{A}\left(  \mathbf{\bar{z},z}\right)
=\sum_{\substack{1\leq l_{1}<l_{2}\leq r,1\leq m_{1}<m_{2}\leq r\\1\leq
p_{1}<p_{2}\leq r,1\leq q_{1}<q_{2}\leq r}}Q_{p_{1}p_{2}q_{1}q_{2}}\left(
\mathbf{z}\right)  \Delta_{p_{1}p_{2}q_{1}q_{2}j_{1}j_{2}k_{1}k_{2}}%
^{-1}A_{j_{1}j_{2}k_{1}k_{2}}.
\end{equation}
In infinite dimensions one can complete the space of Q-symbols in the
$L_{\infty}$ norm producing the banach space $L_{\infty}\left(  \mathbb{C}%
^{2\left(  r-2\right)  },\mu\right)  $ and the space of P-symbols in the
$L_{1}$ norm producing the space $L_{1}\left(  \mathbb{C}^{2\left(
r-2\right)  },\mu\right)  .$

\subsection{N-particle Energy Functional}

The trace of the product of a bounded operator $A\in\mathcal{B}\left(
\mathcal{H}^{2}\right)  $ with a trace class operator $B\in\mathcal{B}%
_{1}\left(  \mathcal{H}^{2}\right)  $ can be expressed in terms of Q and P
symbols $\left(  \text{If the P-symbol of }B\text{ exists}\right)  $ as
\begin{equation}
\operatorname*{Tr}\left\{  A^{\dagger}B\right\}  =\int\overline{Q_{A}\left(
\mathbf{z}\right)  }P_{B}\left(  \mathbf{z}\right)  d\mu\left(  \mathbf{z}%
\right)  =\left\{  \left.  Q_{A}\right\vert P_{B}\right\}  ,
\end{equation}
as
\begin{equation}
\operatorname*{Tr}\left\{  A^{\dagger}B\right\}  =\operatorname*{Tr}\left\{
A^{\dagger}\int P_{B}\left(  \mathbf{z}\right)  \left\vert \mathbf{z}%
\right\rangle \left\langle \mathbf{z}\right\vert d\mu\left(  \mathbf{z}%
\right)  \right\}  =\int\overline{\left\langle \mathbf{z}\left\vert
A\mathbf{z}\right.  \right\rangle }P_{B}\left(  \mathbf{z}\right)  d\mu\left(
\mathbf{z}\right)  .
\end{equation}
Then
\begin{equation}
\operatorname*{Tr}\left\{  A^{\dagger}B\right\}  =\operatorname*{Tr}\left\{
\int\overline{P_{A}}\left(  \mathbf{z}\right)  \left\vert \mathbf{z}%
\right\rangle \left\langle \mathbf{z}\right\vert d\mu\left(  \mathbf{z}%
\right)  B\right\}  =\int\left\langle \mathbf{z}\left\vert B\mathbf{z}\right.
\right\rangle \overline{P_{A}}\left(  \mathbf{z}\right)  d\mu\left(
\mathbf{z}\right)
\end{equation}
and one obtains the skew symmetric relationship
\begin{equation}
\operatorname*{Tr}\left\{  A^{\dagger}B\right\}  =\int\overline{Q_{A}\left(
\mathbf{z}\right)  }P_{B}\left(  \mathbf{z}\right)  d\mu\left(  \mathbf{z}%
\right)  =\int Q_{B}\left(  \mathbf{z}\right)  \overline{P_{A}}\left(
\mathbf{z}\right)  d\mu\left(  \mathbf{z}\right)  .
\end{equation}


The $N$-particle energy can be expressed in terms of the self adjoint Second
Order Reduced Hamiltonian (SORH) $K_{2N}$ and the SORDO\ $D_{2}$
\cite{Colemann} as
\begin{equation}
E_{N}\left(  D_{2}\right)  =\operatorname*{Tr}\left\{  K_{2N}D_{2}\right\}
\end{equation}
thus in terms of Q and P symbols as
\begin{equation}
E_{N}\left(  D_{2}\right)  =\int Q_{D_{2}}\left(  \mathbf{z}\right)
P_{K_{2N}}\left(  \mathbf{z}\right)  d\mu\left(  \mathbf{z}\right)  =\int
Q_{K_{2N}}\left(  \mathbf{z}\right)  P_{D_{2}}\left(  \mathbf{z}\right)
d\mu\left(  \mathbf{z}\right)  =\left\{  \left.  Q_{D_{2}}\right\vert
P_{K_{2N}}\right\}  .
\end{equation}
$\lceil$Hence $Q_{A}\mathbf{\in}L_{\infty}\left(  \mathbb{C}^{2\left(
r-2\right)  },\mu\right)  $ and $P_{A}\mathbf{\in}L_{\infty}\left(
\mathbb{C}^{2\left(  r-2\right)  },\mu\right)  ^{d}=$ space of measures on
$\mathbb{C}^{2\left(  r-2\right)  }.$ A subspace of $L_{\infty}\left(
\mathbb{C}^{2\left(  r-2\right)  },\mu\right)  ^{d}$ is $L_{1}\left(
\mathbb{C}^{2\left(  r-2\right)  },\mu\right)  $, which is the predual of
$L_{\infty}\left(  \mathbb{C}^{2\left(  r-2\right)  },\mu\right)  .\rfloor$

The Q-symbol of the Reduced Hamiltonian is given by
\begin{align}
Q_{K_{2N}}\left(  \mathbf{z}\right)   &  =\frac{1}{4}\sum_{1\leq j_{1}%
,j_{2},k_{1},k_{2}\leq r}K_{2Nj_{1}j_{2}k_{1}k_{2}}\left\langle \mathbf{z}%
\left\vert a_{j_{1}}^{\dagger}a_{j_{2}}^{\dagger}a_{k_{2}}a_{k_{1}}%
\mathbf{z}\right.  \right\rangle \nonumber\\
&  =\frac{1}{4}\left\langle \mathbf{z}\left\vert \mathbf{z}\right.
\right\rangle \sum_{1\leq j_{1},j_{2},k_{1},k_{2}\leq r}K_{2Nj_{1}j_{2}%
k_{1}k_{2}}\left\{  \gamma_{1}\left(  \mathbf{z}\right)  _{j_{1}k_{1}}%
\gamma_{1}\left(  \mathbf{z}\right)  _{j_{2}k_{2}}-\gamma_{1}\left(
\mathbf{z}\right)  _{j_{1}k_{2}}\gamma_{1}\left(  \mathbf{z}\right)
_{j_{2}k_{1}}\right\}  ,
\end{align}
where%
\begin{align}
K_{2Nj_{1}j_{2}k_{1}k_{2}}  &  =\frac{2}{N\left(  N-1\right)  }\eta\left(
\delta_{j_{1}k_{1}}\delta_{j_{2}k_{2}}-\delta_{j_{1}k_{2}}\delta_{j_{2}k_{1}%
}\right) \nonumber\\
&  +\frac{1}{\left(  N-1\right)  }\left(  h_{1j_{1}k_{1}}\delta_{j_{2}k_{2}%
}+h_{1j_{2}k_{2}}\delta_{j_{1}k_{1}}\right)  +\frac{1}{2}h_{2j_{1}j_{2}%
k_{1}k_{2}},
\end{align}
while the P-symbol of the Reduced Hamiltonian can be obtained by using the
relationship eq. $\left(  \ref{P in Q}\right)  ,$ which gives
\begin{align}
P_{K_{2N}}\left(  \mathbf{z}\right)   &  =\sum_{\substack{1\leq l_{1}%
<l_{2}\leq r,1\leq m_{1}<m_{2}\leq r\\1\leq p_{1}<p_{2}\leq r,1\leq
q_{1}<q_{2}\leq r}}Q_{p_{1}p_{2}q_{1}q_{2}}\left(  \mathbf{z}\right)
\Delta_{p_{1}p_{2}q_{1}q_{2}j_{1}j_{2}k_{1}k_{2}}^{-1}K_{2Nj_{1}j_{2}%
k_{1}k_{2}}\\
&  =\left\langle \mathbf{z\left\vert \mathbf{z}\right.  }\right\rangle
\sum_{\substack{1\leq l_{1}<l_{2}\leq r,1\leq m_{1}<m_{2}\leq r\\1\leq
p_{1}<p_{2}\leq r,1\leq q_{1}<q_{2}\leq r}}\gamma_{2}\left(  \mathbf{z}%
\right)  _{p_{1}p_{2}q_{1}q_{2}}\Delta_{p_{1}p_{2}q_{1}q_{2}j_{1}j_{2}%
k_{1}k_{2}}^{-1}K_{2Nj_{1}j_{2}k_{1}k_{2}}\\
&  =\frac{\left\langle \mathbf{z\left\vert \mathbf{z}\right.  }\right\rangle
}{2}\sum_{\substack{1\leq l_{1}<l_{2}\leq r,1\leq m_{1}<m_{2}\leq r\\1\leq
p_{1}<p_{2}\leq r,1\leq q_{1}<q_{2}\leq r}}\left\{  \gamma_{1}\left(
\mathbf{z}\right)  _{p_{1}q_{1}}\gamma_{1}\left(  \mathbf{z}\right)
_{p_{2}q_{2}}-\gamma_{1}\left(  \mathbf{z}\right)  _{p_{1}q_{2}}\gamma
_{1}\left(  \mathbf{z}\right)  _{p_{2}q_{1}}\right\}  \Delta_{p_{1}p_{2}%
q_{1}q_{2}j_{1}j_{2}k_{1}k_{2}}^{-1}K_{2Nj_{1}j_{2}k_{1}k_{2}}.
\end{align}


The Q-symbol of any SORDO is given by
\begin{align}
Q_{D_{2}}\left(  \mathbf{z}\right)   &  =\frac{1}{4}\left\langle
\mathbf{z}\left\vert \mathbf{z}\right.  \right\rangle \sum_{1\leq j_{1}%
,j_{2},k_{1},k_{2}\leq r}D_{2j_{1}j_{2}k_{1}k_{2}}\gamma_{2}\left(
\mathbf{z}\right)  _{k_{1}k_{2}j_{1}j_{2}}\nonumber\\
&  =\frac{1}{4}\left\langle \mathbf{z}\left\vert \mathbf{z}\right.
\right\rangle \sum_{1\leq j_{1},j_{2},k_{1},k_{2}\leq r}D_{2j_{1}j_{2}%
k_{1}k_{2}}\left\{  \gamma_{1}\left(  \mathbf{z}\right)  _{j_{1}k_{1}}%
\gamma_{1}\left(  \mathbf{z}\right)  _{j_{2}k_{2}}-\gamma_{1}\left(
\mathbf{z}\right)  _{j_{1}k_{2}}\gamma_{1}\left(  \mathbf{z}\right)
_{j_{2}k_{1}}\right\}  ,
\end{align}
which we recall can also be expressed as
\begin{align}
&  Q_{D_{2}}\left(  \mathbf{z}\right)  =\left\langle \varphi_{1}\left(
\mathbf{z}\right)  \varphi_{2}\left(  \mathbf{z}\right)  \left\vert
D_{2}\varphi_{1}\left(  \mathbf{z}\right)  \varphi_{2}\left(  \mathbf{z}%
\right)  \right.  \right\rangle \nonumber\\
&  =\operatorname*{Tr}\left\{  \left\vert \mathbf{z}\right\rangle \left\langle
\mathbf{z}\right\vert D_{2}\right\}  =\operatorname*{Tr}\left\{  c\left(
\mathbf{z}\right)  \left\vert \varphi_{1}\varphi_{2}\right\rangle \left\langle
\varphi_{1}\varphi_{2}\right\vert c\left(  \mathbf{z}\right)  ^{\dagger
}\right\}  =f_{D_{2}}\left(  c\left(  \mathbf{z}\right)  a_{1}^{\dagger}%
a_{2}^{\dagger}a_{2}a_{1}c\left(  \mathbf{z}\right)  ^{\dagger}\right)  .
\end{align}


\subsection{Minimization Problem}

The basic minimization problem to find the energy and the SORDO of the ground
state is given by
\begin{equation}
E_{N}\left(  D_{2GS}\right)  =\min_{D_{2}\in S_{2N}}\left\{
\operatorname*{Tr}\left\{  K_{2N}D_{2}\right\}  \right\}  =\min_{Q_{D_{2}}%
\in\mathfrak{S}_{2N}}\left\{  \int Q_{D_{2}}\left(  \mathbf{z}\right)
P_{K_{2N}}\left(  \mathbf{z}\right)  d\mu\left(  \mathbf{z}\right)  \right\}
,
\end{equation}
where $\mathfrak{S}_{2N}$ is the set of N-Representable Q symbols, which we
chacterize in the next section.

\section{N-Representability Conditions for Q-Symbols}

The N-Representability conditions expressed eqs.$\left(  \ref{NormCon}\right)
-\left(  \ref{GCond}\right)  $ can be recast in terms of Q-symbols as

\begin{enumerate}
\item \bigskip the normalization condition
\begin{equation}
\mathcal{L}_{N}\left(  Q_{D_{2}}\right)  =\int Q_{D_{2}}\left(  \mathbf{z}%
\right)  \left\langle \mathbf{z\left\vert \mathbf{z}\right.  }\right\rangle
d\mu\left(  \mathbf{z}\right)  =\binom{N}{2} \label{QsymNorm}%
\end{equation}


\item the generalized P-condition%
\begin{equation}
1\geq Q_{D_{2}}\left(  \mathbf{z}\right)  \geq0 \label{QsymP}%
\end{equation}


\item the generalized Q-condition
\begin{equation}
1\geq\left[  \mathcal{L}_{11}\left(  Q_{D_{2}}\right)  \right]  \left(
\mathbf{z}\right)  -Q_{D_{2}}\left(  \mathbf{z}\right)  \geq0 \label{QsymQ1}%
\end{equation}
or%
\begin{equation}
1\geq\left[  \mathcal{L}_{22}\left(  Q_{D_{2}}\right)  \right]  \left(
\mathbf{z}\right)  -Q_{D_{2}}\left(  \mathbf{z}\right)  \geq0 \label{QsymQ2}%
\end{equation}


\item the generalized G-condition%
\begin{equation}
1\geq\left[  \mathcal{L}_{11}\left(  Q_{D_{2}}\right)  \right]  \left(
\mathbf{z}\right)  +\left[  \mathcal{L}_{22}\left(  Q_{D_{2}}\right)  \right]
\left(  \mathbf{z}\right)  -Q_{D_{2}}\left(  \mathbf{z}\right)  \geq0.
\label{QsymG}%
\end{equation}

\end{enumerate}

The linear functionals $\mathcal{L}_{jk}$ produce the FORDM elements
\begin{equation}
D_{1}\left(  \mathbf{z}\right)  _{jk}=\left[  \mathcal{L}_{jk}\left(
Q_{D_{2}}\right)  \right]  \left(  \mathbf{z}\right)  =\int\mathcal{L}%
_{jk}\left(  \mathbf{z,z}^{\prime}\right)  Q_{D_{2}}\left(  \mathbf{z}%
^{\prime}\right)  d\mu\left(  \mathbf{z}^{\prime}\right)  \label{Lij}%
\end{equation}
and are defined as
\[
\mathcal{L}_{jk}\left(  \mathbf{z,z}^{\prime}\right)  =\frac{1}{2\left(
N-1\right)  }\sum_{1\leq j_{1},k_{1},l\leq r}P_{j_{1}lk_{1}l}\left(
\mathbf{z}^{\prime}\right)  \operatorname*{Tr}\left\{  \left\vert
\tilde{\varphi}_{j}\left(  \mathbf{z}\right)  \right\rangle \left\langle
\tilde{\varphi}_{k}\left(  \mathbf{z}\right)  \right\vert
\overset{}{\left\vert \varphi_{j_{1}}\right\rangle \left\langle \varphi
_{k_{1}}\right\vert }\right\}  K_{Q}\left(  \mathbf{z,z}^{\prime}\right)  ,
\]
where the reproducing kernel $K_{Q}\left(  \mathbf{z,z}^{\prime}\right)  $
that forms the projector onto the space of Q-symbols is given by
\begin{align}
\left\{  \mathbf{z}\left\vert Q\right.  \right\}  \left\{  P\left\vert
\mathbf{z}^{\prime}\right.  \right\}   &  =K_{Q}\left(  \mathbf{z,z}^{\prime
}\right)  =\sum_{1\leq j_{1}<j_{2},k_{1}<k_{2}\leq r}Q_{j_{1}j_{2}k_{1}k_{2}%
}\left(  \mathbf{z}\right)  \overline{P_{j_{1}j_{2}k_{1}k_{2}}\left(
\mathbf{z}^{\prime}\right)  }\nonumber\\
&  =\sum_{\substack{1\leq p_{1}<p_{2}\leq r,1\leq q_{1}<q_{2}\leq r\\1\leq
j_{1}<j_{2}\leq r,1\leq k_{1}<k_{2}\leq r}}\overline{Q_{p_{1}p_{2}q_{1}q_{2}%
}\left(  \mathbf{z}^{\prime}\right)  \Delta_{p_{1}p_{2}q_{1}q_{2}j_{1}%
j_{2}k_{1}k_{2}}^{-1}}Q_{j_{1}j_{2}k_{1}k_{2}}\left(  \mathbf{z}\right)  .
\end{align}
The condition eq. $\left(  \ref{QsymNorm}\right)  $ can be expressed in terms
of a $L_{1}\left(  \mathbb{C}^{r\left(  r-2\right)  },\tilde{\mu}\right)  $
norm as
\begin{equation}
\left\Vert Q_{D_{2}}\right\Vert _{L_{1}\left(  \mathbb{C}^{r\left(
r-2\right)  },\tilde{\mu}\right)  }=\int Q_{D_{2}}\left(  \mathbf{z}\right)
d\tilde{\mu}\left(  \mathbf{z}\right)  =\binom{N}{2},
\end{equation}
where
\begin{equation}
\tilde{\mu}\left(  \mathbf{z}\right)  =\left\langle \mathbf{z\left\vert
\mathbf{z}\right.  }\right\rangle d\mu\left(  \mathbf{z}\right)  .
\end{equation}


The condition eq. $\left(  \ref{QsymP}\right)  $ can be expressed in terms of
a $L_{\infty}\left(  \mathbb{C}^{2\left(  r-2\right)  },\mu\right)  $ norm as
\begin{equation}
\left\Vert Q_{D_{2}}\right\Vert _{L_{\infty}\left(  \mathbb{C}^{2\left(
r-2\right)  },\mu\right)  }=\sup_{\mathbf{z\in}\mathbb{C}^{2\left(
r-2\right)  }}\left\{  Q_{D_{2}}\right\}  =\sup_{\mathbf{z\in}\mathbb{C}%
^{r\left(  r-2\right)  }}\left\{  \left\langle \varphi_{1}\left(
\mathbf{z}\right)  \varphi_{2}\left(  \mathbf{z}\right)  \left\vert
D_{2}\varphi_{1}\left(  \mathbf{z}\right)  \varphi_{2}\left(  \mathbf{z}%
\right)  \right.  \right\rangle \right\}  \leq1,
\end{equation}
and hence in terms of a $L_{\infty}\left(  \mathbb{C}^{2\left(  r-2\right)
},\tilde{\mu}\right)  $ norm as
\begin{equation}
\left\Vert Q_{D_{2}}\right\Vert _{L_{\infty}\left(  \mathbb{C}^{2\left(
r-2\right)  },\tilde{\mu}\right)  }=\sup_{\mathbf{z\in}\mathbb{C}^{2\left(
r-2\right)  }}\left\{  \tilde{Q}_{D_{2}}\right\}  =\sup_{\mathbf{z\in
}\mathbb{C}^{r\left(  r-2\right)  }}\left\{  \frac{\left\langle \varphi
_{1}\left(  \mathbf{z}\right)  \varphi_{2}\left(  \mathbf{z}\right)
\left\vert D_{2}\varphi_{1}\left(  \mathbf{z}\right)  \varphi_{2}\left(
\mathbf{z}\right)  \right.  \right\rangle }{\det\left\{  \mathbb{I}%
_{2}-\mathbf{z}^{\dagger}\mathbf{z}\right\}  }\right\}  \leq1.
\end{equation}


\section{Approximations}

If $r=\infty$ then space $\mathbb{C}^{2\left(  r-2\right)  }$ becomes
$\mathbb{C}^{\infty}$ and has to be replaced by a function space i.e.
consider
\begin{equation}
\left(  z_{1},\cdots,z_{2\left(  r-2\right)  }\right)  \equiv\left(  z\left(
1\right)  ,\cdots,z\left(  2\left(  r-2\right)  \right)  \right)
\end{equation}%
\begin{equation}
z:\left\{  1,\cdots,2\left(  r-2\right)  \right\}  \mathbb{\rightarrow C}%
\end{equation}
so as $r\rightarrow\infty$ the function $z$ becomes
\begin{equation}
z:\mathbb{N\rightarrow C}%
\end{equation}
and can be thought as the values of a function
\begin{equation}
z:\mathbb{R\rightarrow C}%
\end{equation}
restricted to a countably infinte points in $\mathbb{R}.$ Then one use the
Wierstrauss theorem and $\mathbf{z}$ is replaced by a function $z$. If
$r<\infty$ then one just uses the parameter values.

Even if $r<\infty$ the spaces $L_{\infty}\left(  \mathbb{C}^{2\left(
r-2\right)  },\mu\right)  $ and $L_{1}\left(  \mathbb{C}^{2\left(  r-2\right)
},\mu\right)  $ are infinite dimensional and one can consider further
approximations by finite dimensional subspaces $\mathbb{C}^{\kappa}$ generated
by $\kappa$ basis functions.

\section{Discussion}

\appendix


\section{Mathematical Essentials}

\label{Math}

In the following the dimension of all spaces are considered to be finite and
can be thought of as approximations to the infinite dimensional case, which
must be used to obtain the exact solutions of quantum mechanical equations.
However notation that is appropriate for the infinite dimensional case will be
used even though it becomes redundant in the finite dimensional case.

\subsection{Fermion Fock Space \label{Fermi Fock Space}}

We consider the Fermi Fock space, $\mathcal{H}_{\mathcal{F}},$ based on the
one electron Hilbert space $\mathcal{H}^{1},$ defined by the direct sum
\begin{equation}
\mathcal{H}_{\mathcal{F}}=\underset{0\leq N\leq r}{\bigoplus}\mathcal{H}^{N}%
\end{equation}
of the $N$-electron Hilbert spaces
\begin{equation}
\mathcal{H}^{N}=\bigwedge\limits^{N}\mathcal{H}^{1},
\end{equation}
where $\wedge$ denotes antisymmetric tensor product and $\mathcal{H}%
^{0}=\left\{  \left.  \lambda\left\vert \phi\right\rangle \right\vert
\lambda\in\mathbb{C}\right\}  ,$ and $\left\vert \phi\right\rangle $ is the
vacuum state. The inner product in the Hilbert space $\mathcal{H}%
_{\mathcal{F}}$ is defined as
\begin{equation}
\left\langle \Psi\left\vert \Phi\right.  \right\rangle _{\mathcal{F}}%
=\sum_{0\leq N\leq r}\left\langle \Psi_{N}\left\vert \Phi_{N}\right.
\right\rangle _{\mathcal{H}^{N}}=\sum_{0\leq N\leq r}\left\langle
\Psi\left\vert P_{\mathcal{H}^{N}}\Phi\right.  \right\rangle _{\mathcal{H}%
^{N}},
\end{equation}
where $P_{\mathcal{H}^{N}}$ is the orthogonal projector onto $\mathcal{H}^{N}$
and $\left\langle \cdot\left\vert \cdot\right.  \right\rangle _{\mathcal{H}%
^{N}}$ is the inner product in this Hilbert space. (Notational note: If the
domain of $\left\langle \cdot\left\vert \cdot\right.  \right\rangle
_{\mathcal{H}^{N}}$ is obvious from the context the subscript "$\mathcal{H}%
^{N}"$ will not be explicitly shown.)Fermion Fock Algebra
\label{Fermi Fock Alg}

The fermion Fock algebra, $\mathcal{F},$ is the space of bounded operators
$\left(  \text{Appendix \ref{Operators}}\right)  $ acting in $\mathcal{H}%
_{\mathcal{F}}$ i.e.
\begin{equation}
\mathcal{F}=\mathcal{B}\left(  \mathcal{H}_{\mathcal{F}}\right)  .
\end{equation}
This space can be generated by polynomials
\begin{equation}
X=\sum_{0\leq P,Q\leq r}\sum_{\substack{1\leq j_{1}<\cdots<j_{P}\leq r\\1\leq
k_{1}<\cdots<k_{Q}\leq r}}X_{j_{1}\cdots j_{P}k_{1}\cdots k_{P}}a_{j_{1}%
}^{\dagger}\cdots a_{j_{P}}^{\dagger}a_{k_{Q}}\cdots a_{k_{1}}%
\end{equation}
of second quantized operators, $\left\{  \left.  a_{j}^{\dagger}\right\vert
1\leq j\leq r\right\}  ,$ which are defined by their action on the vacuum
vector $\left\vert \phi\right\rangle $ by
\begin{equation}
a_{j}^{\dagger}\left\vert \phi\right\rangle =\left\vert \varphi_{j}%
\right\rangle ,1\leq j\leq r,
\end{equation}
where $\left\{  \left.  \left\vert \varphi_{j}\right\rangle \right\vert 1\leq
j\leq r\right\}  $ is a complete orthonormal basis of the one-electron space
$\mathcal{H}^{1}$ and satisfy the anti-commutation relationships given by
\begin{equation}
\left[  a_{j}^{\dagger},a_{k}\right]  _{+}=a_{j}^{\dagger}a_{k}+a_{k}%
a_{j}^{\dagger}=\delta_{jk}I_{\mathcal{H}_{\mathcal{F}}};1\leq j,k\leq r.
\label{CAR}%
\end{equation}


The fermion Fock algebra can be decomposed as a vector space direct sum%
\begin{equation}
\mathcal{F=F}_{e}\mathcal{\oplus F}_{0}%
\end{equation}
of a electron conserving subalgebra $\mathcal{F}_{e}$ and a electron non
conserving subspace of operators $\mathcal{F}_{0}$. The subalgebra
$\mathcal{F}_{e}$ can be further decomposed as a direct sum of subspaces
\begin{equation}
\mathcal{F}_{e}=%
%TCIMACRO{\dbigoplus \limits_{0\leq N\leq r}}%
%BeginExpansion
{\displaystyle\bigoplus\limits_{0\leq N\leq r}}
%EndExpansion
\mathcal{B}\left(  \mathcal{H}_{\mathcal{F}}\right)  _{N}, \label{evenvect}%
\end{equation}
where
\begin{equation}
\mathcal{B}\left(  \mathcal{H}_{\mathcal{F}}\right)  _{N}=\left\{
\sum_{\substack{1\leq j_{1}<\cdots<j_{N}\leq r\\1\leq k_{1}<\cdots<k_{N}\leq
r}}X_{j_{1}\cdots j_{N}k_{1}\cdots k_{N}}a_{j_{1}}^{\dagger}\cdots a_{j_{N}%
}^{\dagger}a_{k_{N}}\cdots a_{k_{1}}\right\}
\end{equation}
or as a direct sum of subalgebras
\begin{equation}
\mathcal{F}_{e}=%
%TCIMACRO{\dbigoplus \limits_{0\leq N\leq r}}%
%BeginExpansion
{\displaystyle\bigoplus\limits_{0\leq N\leq r}}
%EndExpansion
\mathcal{B}\left(  \mathcal{H}^{N}\right)  , \label{evenalg}%
\end{equation}
where $\mathcal{B}\left(  \mathcal{H}^{N}\right)  $ is the space of bounded
operators acting in $\mathcal{H}^{N}.$ The subspace $\mathcal{F}_{0}$ can be
decomposed as the vector space direct sum
\begin{equation}
\mathcal{F}_{o}=%
%TCIMACRO{\dbigoplus \limits_{0\leq N\neq M\leq r}}%
%BeginExpansion
{\displaystyle\bigoplus\limits_{0\leq N\neq M\leq r}}
%EndExpansion
\mathcal{B}\left(  \mathcal{H}_{\mathcal{F}}\right)  _{N,M}, \label{oddsecond}%
\end{equation}
where
\begin{equation}
\mathcal{B}\left(  \mathcal{H}_{\mathcal{F}}\right)  _{N,M}=\left\{
\sum_{\substack{1\leq j_{1}<\cdots<j_{N}\leq r\\1\leq k_{1}<\cdots<k_{M}\leq
r}}X_{j_{1}\cdots j_{N}k_{1}\cdots k_{M}}a_{j_{1}}^{\dagger}\cdots a_{j_{N}%
}^{\dagger}a_{k_{M}}\cdots a_{k_{1}}\right\}
\end{equation}
or as the vector space direct sum
\begin{equation}
\mathcal{F}_{o}=%
%TCIMACRO{\dbigoplus \limits_{0\leq N\neq M\leq r}}%
%BeginExpansion
{\displaystyle\bigoplus\limits_{0\leq N\neq M\leq r}}
%EndExpansion
\mathcal{B}\left(  \mathcal{H}^{N},\mathcal{H}^{M}\right)  , \label{oddfirst}%
\end{equation}
where $\mathcal{B}\left(  \mathcal{H}^{N},\mathcal{H}^{M}\right)  $ is the
vector space of bounded linear maps from the Hilbert space $\mathcal{H}^{N}$
to the Hilbert space $\mathcal{H}^{M}.$ The decompositions eqs. $\left(
\ref{evenvect}\right)  $ and $\left(  \ref{oddsecond}\right)  $ express the
fermion Fock algebra in \emph{second quantized} form as polynomials of second
quantized operators, while the decompositions eqs. $\left(  \ref{evenalg}%
\right)  $ and $\left(  \ref{oddfirst}\right)  $ express it in \emph{first
quantized} form as expansions of the ket-bras
\begin{equation}
\left\{  \left.  \left\vert \varphi_{j_{1}}\cdots\varphi_{j_{N}}\right\rangle
\left\langle \varphi_{k_{1}}\cdots\varphi_{k_{M}}\right\vert \right\vert 1\leq
j_{1}<\cdots<j_{N}\leq r,1\leq k_{1}<\cdots<k_{M}\leq r;0\leq N,M\leq
r\right\}
\end{equation}


\subsection{Operator Inner Product \label{Inner prod}}

The Fock space trace operation is defined as
\begin{equation}
\operatorname*{Tr}\left\{  X\right\}  =\sum_{0\leq N\leq r}\sum_{1\leq
j_{1}<\cdots<j_{N}\leq r}\left\langle \Psi_{Nj_{1}\cdots j_{N}}\left\vert
X\Psi_{Nj_{1}\cdots j_{N}}\right.  \right\rangle ,
\end{equation}
where $\left\{  \left.  \Psi_{Nj_{1}\cdots j_{N}}\right\vert 1\leq
j_{1}<\cdots<j_{N}\leq r;0\leq N\leq r\right\}  $ are complete orthonormal
bases for $\mathcal{H}^{N},0\leq N\leq r.$ One can use this trace operation to
define the space of trace class operators $\mathcal{B}_{1}\left(
\mathcal{H}_{\mathcal{F}}\right)  $ that have finite trace for all bases and
the space of Hilbert Schmidt operators $\mathcal{B}_{2}\left(  \mathcal{H}%
_{\mathcal{F}}\right)  $ that form a Hilbert space with inner product
\begin{equation}
\left(  X\left\vert Y\right.  \right)  =\operatorname*{Tr}\left\{  X^{\dagger
}Y\right\}  .
\end{equation}
In finite dimensions, i.e. $r<\infty,$ there is no distinction between these
spaces of operators i.e. $\mathcal{B}\left(  \mathcal{H}_{\mathcal{F}}\right)
=\mathcal{B}_{1}\left(  \mathcal{H}_{\mathcal{F}}\right)  =\mathcal{B}%
_{2}\left(  \mathcal{H}_{\mathcal{F}}\right)  .$

\section{Expansion and Contraction Maps \label{Expansion and Contraction}}

The expansion and contraction maps relate the first and second quantization
formulation of quantum mechanics and are defined by the following:

\subsection{Contraction Map \label{Contraction}}

Employing the work of \cite{Kummar70} one can define a \emph{non singular
}generalized contraction map
\begin{equation}
L:\mathcal{B}_{1}\left(  \mathcal{H}_{\mathcal{F}}\right)  \rightarrow
\mathcal{B}_{1}\left(  \mathcal{H}_{\mathcal{F}}\right)
\end{equation}
by
\begin{equation}
L\left(  Y\right)  =\sum_{0\leq N,M\leq r}\sum_{\substack{1\leq j_{1}%
<\cdots<j_{M}\leq r\\1\leq k_{1}<\cdots<k_{N}\leq r}}\left(  \left.  a_{k_{1}%
}^{\dagger}\cdots a_{k_{N}}^{\dagger}a_{j_{M}}\cdots a_{j_{1}}\right\vert
Y\right)  \left\vert \varphi_{j_{1}}\cdots\varphi_{j_{M}}\right\rangle
\left\langle \varphi_{k_{1}}\cdots\varphi_{k_{N}}\right\vert . \label{CONT}%
\end{equation}


\subsection{Expansion Map \label{Expansion}}

The expansion map
\begin{equation}
\Gamma:\mathcal{B}\left(  \mathcal{H}_{\mathcal{F}}\right)  \rightarrow
\mathcal{B}\left(  \mathcal{H}_{\mathcal{F}}\right)
\end{equation}
is adjoint to the contraction map, i.e. $\Gamma=L^{\dagger},$ with respect to
inner product $\left(  \cdot\left\vert \cdot\right.  \right)  $ and is defined
by
\begin{equation}
\left(  \Gamma\left(  X\right)  \left\vert Y\right.  \right)  =\left(
X\left\vert L\left(  Y\right)  \right.  \right)  \quad\forall Y\in
\mathcal{B}_{1}\left(  \mathcal{H}_{\mathcal{F}}\right)  .
\end{equation}
One can show (Appendix \ref{SecQuant}) that $\Gamma$ is \emph{actually the
second quantization map}
\begin{equation}
a_{j_{1}}^{\dagger}\cdots a_{j_{N}}^{\dagger}a_{k_{P}}\cdots a_{k_{M}}%
=\Gamma\left(  \left\vert \varphi_{j_{1}}\cdots\varphi_{j_{M}}\right\rangle
\left\langle \varphi_{k_{1}}\cdots\varphi_{k_{N}}\right\vert \right)  ,
\end{equation}
which can be expressed as
\begin{equation}
\Gamma\left(  \left\vert \varphi_{j_{1}}\cdots\varphi_{j_{M}}\right\rangle
\left\langle \varphi_{k_{1}}\cdots\varphi_{k_{N}}\right\vert \right)
=\sum_{1\leq l_{1}<\cdots<l_{P}\leq r}\left\vert \varphi_{j_{1}}\cdots
\varphi_{j_{M}}\varphi_{l_{1}}\cdots\varphi_{l_{P}}\right\rangle \left\langle
\varphi_{k_{1}}\cdots\varphi_{k_{N}}\varphi_{l_{1}}\cdots\varphi_{l_{P}%
}\right\vert .
\end{equation}
The expansion and contraction maps transform naturally between the first and
second quantization direct sum decompositions of $\mathcal{F}$ i.e.
\begin{equation}
\Gamma:%
%TCIMACRO{\dbigoplus \limits_{0\leq N,M\leq r}}%
%BeginExpansion
{\displaystyle\bigoplus\limits_{0\leq N,M\leq r}}
%EndExpansion
\mathcal{B}\left(  \mathcal{H}^{N},\mathcal{H}^{M}\right)  \rightarrow%
%TCIMACRO{\dbigoplus \limits_{0\leq N,M\leq r}}%
%BeginExpansion
{\displaystyle\bigoplus\limits_{0\leq N,M\leq r}}
%EndExpansion
\mathcal{B}\left(  \mathcal{H}_{\mathcal{F}}\right)  _{N,M}%
\end{equation}
and
\begin{equation}
L:%
%TCIMACRO{\dbigoplus \limits_{0\leq N,M\leq r}}%
%BeginExpansion
{\displaystyle\bigoplus\limits_{0\leq N,M\leq r}}
%EndExpansion
\mathcal{B}\left(  \mathcal{H}_{\mathcal{F}}\right)  _{N,M}\rightarrow%
%TCIMACRO{\dbigoplus \limits_{0\leq N,M\leq r}}%
%BeginExpansion
{\displaystyle\bigoplus\limits_{0\leq N,M\leq r}}
%EndExpansion
\mathcal{B}\left(  \mathcal{H}^{N},\mathcal{H}^{M}\right)  .
\end{equation}
The normalized Q-symbol of the two electron reference state is
\begin{align}
\tilde{Q}_{\gamma_{2\varphi}}\left(  \mathbf{z}\right)   &  =\frac
{\left\langle \mathbf{z}\left\vert \left\vert \varphi_{1}\varphi
_{2}\right\rangle \left\langle \varphi_{1}\varphi_{2}\right\vert \right.
\mathbf{z}\right\rangle }{\left\langle \mathbf{z\left\vert \mathbf{z}\right.
}\right\rangle }=\left\langle \varphi_{1}\varphi_{2}\left\vert \frac
{\left\vert \mathbf{z}\right\rangle \left\langle \mathbf{z}\right\vert
}{\left\langle \mathbf{z\left\vert \mathbf{z}\right.  }\right\rangle }\right.
\varphi_{1}\varphi_{2}\right\rangle =\gamma_{2}\left(  \mathbf{z}\right)
_{1212}\nonumber\\
&  =\left\{  \gamma_{1}\left(  \mathbf{z}\right)  _{11}\gamma_{1}\left(
\mathbf{z}\right)  _{22}-\gamma_{1}\left(  \mathbf{z}\right)  _{12}\gamma
_{1}\left(  \mathbf{z}\right)  _{21}\right\}  =\left\vert
\begin{array}
[c]{cc}%
\gamma_{1}\left(  \mathbf{z}\right)  _{11} & \gamma_{1}\left(  \mathbf{z}%
\right)  _{12}\\
\gamma_{1}\left(  \mathbf{z}\right)  _{21} & \gamma_{1}\left(  \mathbf{z}%
\right)  _{22}%
\end{array}
\right\vert \nonumber\\
&  =\left\vert
\begin{array}
[c]{cc}%
\left[  \left(  \mathbb{I}_{2}+\mathbf{z}^{\dagger}\mathbf{z}\right)
^{-1}\right]  _{11} & \left[  \left(  \mathbb{I}_{2}+\mathbf{z}^{\dagger
}\mathbf{z}\right)  ^{-1}\right]  _{12}\\
\left[  \left(  \mathbb{I}_{2}+\mathbf{z}^{\dagger}\mathbf{z}\right)
^{-1}\right]  _{21} & \left[  \left(  \mathbb{I}_{2}+\mathbf{z}^{\dagger
}\mathbf{z}\right)  ^{-1}\right]  _{22}%
\end{array}
\right\vert =\det\left(  \mathbb{I}_{2}+\mathbf{z}^{\dagger}\mathbf{z}\right)
^{-1},
\end{align}
The SORDO $D_{2}\left(  \mathbf{z}\right)  $ of $D\left(  \mathbf{z}\right)
=\widehat{u\left(  \mathbf{z}\right)  }^{\dagger}\left(  D\right)  $ is
defined by%
\begin{equation}
D_{2}\left(  \mathbf{z}\right)  =\sum_{1\leq j_{1}<j_{2}\leq r,1\leq
k_{1}<k_{2}\leq r}D_{2}\left(  \mathbf{z}\right)  {}_{j_{1}j_{2}k_{1}k_{2}%
}\left\vert \varphi_{j_{1}}\varphi_{j_{2}}\right\rangle \left\langle
\varphi_{k_{1}}\varphi_{k_{2}}\right\vert ,
\end{equation}
where%
\begin{align}
&  f_{D}\left(  \Phi\left(  \tilde{\varphi}_{j_{1}}\left(  \mathbf{z}\right)
\right)  ^{\dagger}\Phi\left(  \tilde{\varphi}_{j_{2}}\left(  \mathbf{z}%
\right)  \right)  ^{\dagger}\Phi\left(  \tilde{\varphi}_{k_{2}}\left(
\mathbf{z}\right)  \right)  \Phi\left(  \tilde{\varphi}_{k_{1}}\left(
\mathbf{z}\right)  \right)  \right)  =f_{D}\left(  \widehat{u\left(
\mathbf{z}\right)  }\left(  a_{j_{1}}^{\dagger}a_{j_{2}}^{\dagger}a_{k_{2}%
}a_{k_{1}}\right)  \right) \nonumber\\
&  =\left(  \left.  \overset{}{\widehat{u\left(  \mathbf{z}\right)  }\left(
a_{k_{1}}^{\dagger}a_{k_{2}}^{\dagger}a_{j_{2}}a_{j_{1}}\right)  }\right\vert
D\right)  =\left(  \left.  \overset{}{\widehat{u\left(  \mathbf{z}\right)
}\left(  \Gamma\left(  \left\vert \varphi_{k_{1}}\varphi_{k_{2}}\right\rangle
\left\langle \varphi_{j_{1}}\varphi_{j_{2}}\right\vert \right)  \right)
}\right\vert D\right) \nonumber\\
&  =\left(  \left.  \overset{}{\Gamma\left(  \left\vert \varphi_{k_{1}}%
\varphi_{k_{2}}\right\rangle \left\langle \varphi_{j_{1}}\varphi_{j_{2}%
}\right\vert \right)  }\right\vert \widehat{u\left(  \mathbf{z}\right)
}^{\dagger}\left(  D\right)  \right)  =\left(  \left.  \overset{}{\left\vert
\varphi_{k_{1}}\varphi_{k_{2}}\right\rangle \left\langle \varphi_{j_{1}%
}\varphi_{j_{2}}\right\vert }\right\vert L_{2}\left(  \widehat{u\left(
\mathbf{z}\right)  }^{\dagger}\left(  D\right)  \right)  \right) \nonumber\\
&  =\left(  \left.  \overset{}{\left\vert \varphi_{k_{1}}\varphi_{k_{2}%
}\right\rangle \left\langle \varphi_{j_{1}}\varphi_{j_{2}}\right\vert
}\right\vert \widehat{u\left(  \mathbf{z}\right)  }^{\dagger}L_{2}\left(
\left(  D\right)  \right)  \right)  =D_{2}\left(  \mathbf{z}\right)
_{j_{1}j_{2}k_{1}k_{2}}.
\end{align}
The Q-symbol of $D_{2}$ is given as
\begin{equation}
Q_{D_{2}}\left(  \mathbf{z}\right)  =D_{2}\left(  \mathbf{z}\right)  _{1212},
\end{equation}
i.e. it is the $\left(  12,12\right)  $ matrix element of the SORDO
$\widehat{u\left(  \mathbf{z}\right)  }^{\dagger}\left(  D_{2}\right)  $ and
as $\mathbf{z\in}\mathbb{C}^{2\left(  r-2\right)  }$ varies $\left(
12,12\right)  $ matrix elements are generated for \emph{all }possible SORDO's.
Thus
\begin{align}
&  \widehat{u\left(  \mathbf{z}\right)  }^{\dagger}\left(  D_{2}\right)
\nonumber\\
&  \equiv\sum_{1\leq j_{1}<j_{2}\leq r,1\leq k_{1}<k_{2}\leq r}\left\{
\left.  \overset{}{Q_{D_{2}}}\right\vert \overset{}{P_{j_{1}j_{2}k_{1}k_{2}}%
}\right\}  \left\vert \widehat{u\left(  \mathbf{z}\right)  }^{\dagger}\left(
\overset{}{\left\vert \varphi_{j_{1}}\varphi_{j_{2}}\right\rangle \left\langle
\varphi_{k_{1}}\varphi_{k_{2}}\right\vert }\right)  \right)  =\left\vert
\widehat{u\left(  \mathbf{z}\right)  }^{\dagger}\left(  D_{2}\right)  \right)
\end{align}
and
\begin{align}
Q_{D_{2}}\left(  \mathbf{z}\right)   &  =\sum_{1\leq j_{1}<j_{2}\leq r,1\leq
k_{1}<k_{2}\leq r}\left\{  \left.  \overset{}{Q_{D_{2}}}\right\vert
\overset{}{P_{j_{1}j_{2}k_{1}k_{2}}}\right\}  \left(  \left\vert \varphi
_{1}\varphi_{2}\right\rangle \left\langle \varphi_{1}\varphi_{2}\right\vert
\left\vert \widehat{u\left(  \mathbf{z}\right)  }^{\dagger}\left(
\overset{}{\left\vert \varphi_{j_{1}}\varphi_{j_{2}}\right\rangle \left\langle
\varphi_{k_{1}}\varphi_{k_{2}}\right\vert }\right)  \right.  \right)
\nonumber\\
&  =\left(  \left\vert \varphi_{1}\varphi_{2}\right\rangle \left\langle
\varphi_{1}\varphi_{2}\right\vert \left\vert \widehat{u\left(  \mathbf{z}%
\right)  }^{\dagger}\left(  D_{2}\right)  \right.  \right)  .
\end{align}
the FORDM\ by
\begin{equation}
D_{1}\left(  \mathbf{z}\right)  _{j_{1}k_{1}}=\frac{1}{2\left(  N-1\right)
}\left\{  Q_{D_{2}}\left\vert \sum_{1\leq m,n,l\leq r}P_{j_{1}lk_{1}l}\right.
\right\}  \left(  \left\vert \varphi_{j_{1}}\right\rangle \left\langle
\varphi_{k_{1}}\right\vert \left\vert \widehat{u\left(  \mathbf{z}\right)
}^{\dagger}\left(  \overset{}{\left\vert \varphi_{m}\right\rangle \left\langle
\varphi_{n}\right\vert }\right)  \right.  \right)
\end{equation}
and in particular
\begin{align}
D_{1}\left(  \mathbf{z}\right)  _{11}  &  =\left(  \left\vert \varphi
_{1}\right\rangle \left\langle \varphi_{1}\right\vert \left\vert
\widehat{u\left(  \mathbf{z}\right)  }^{\dagger}\left(  D_{1}\right)  \right.
\right) \nonumber\\
&  =\frac{1}{2\left(  N-1\right)  }\left\{  Q_{D_{2}}\left\vert \sum_{1\leq
m,n,l\leq r}P_{j_{1}lk_{1}l}\right.  \right\}  \left(  \left\vert \varphi
_{1}\right\rangle \left\langle \varphi_{1}\right\vert \left\vert
\widehat{u\left(  \mathbf{z}\right)  }^{\dagger}\left(  \overset{}{\left\vert
\varphi_{m}\right\rangle \left\langle \varphi_{n}\right\vert }\right)
\right.  \right)  .
\end{align}
Thus one can observe that $D_{1}\left(  \mathbf{z}\right)  _{11}$ is a linear
functional of $Q_{D_{2}}$ i.e.
\begin{equation}
D_{1}\left(  \mathbf{z}\right)  _{11}=\left[  \mathcal{L}_{11}\left(
Q_{D_{2}}\right)  \right]  \left(  \mathbf{z}\right)  =\int\mathcal{L}%
_{11}\left(  \mathbf{z,z}^{\prime}\right)  Q_{D_{2}}\left(  \mathbf{z}%
^{\prime}\right)  d\mu\left(  \mathbf{z}^{\prime}\right)
\end{equation}
and $D_{1}\left(  \mathbf{z}\right)  _{11}$ is a Q-symbol$.$

\section{Proof of Generalized P,Q and G Conditions. \label{Proof}}

\subsection{Idempotency of $\lambda_{1}+\lambda_{2}$}

The general operator is given by%
\begin{equation}
\lambda_{1}+\lambda_{2}=\sum_{1\leq l\leq r}\alpha_{l}n_{l}+\sum
_{\substack{1\leq l_{1}<l_{2}\leq r\\1\leq l_{3}<l_{4}\leq r}}\lambda
_{2l_{1}l_{2}l_{3}l_{4}}a_{l_{1}}^{\dagger}a_{l_{2}}^{\dagger}a_{l_{4}%
}a_{l_{3}}.
\end{equation}
The idempotency condition leads to%
\begin{equation}
\left(  \lambda_{1}+\lambda_{2}\right)  ^{2}-\lambda_{1}-\lambda_{2}%
=\lambda_{1}^{2}-\lambda_{1}+\lambda_{2}^{2}-\lambda_{2}+\lambda_{1}%
\lambda_{2}+\lambda_{2}\lambda_{1}=0.
\end{equation}
As
\begin{equation}
\lambda=\sum_{0\leq P\leq r}\Gamma\left(  P_{\mathcal{H}^{P}}\lambda
P_{\mathcal{H}^{P}}\right)
\end{equation}
one obtains
\begin{gather}
\sum_{0\leq P\leq4}\Gamma\left(  P_{\mathcal{H}^{P}}\left(  \lambda_{1}%
^{2}-\lambda_{1}+\lambda_{2}^{2}-\lambda_{2}+\lambda_{1}\lambda_{2}%
+\lambda_{2}\lambda_{1}\right)  P_{\mathcal{H}^{P}}\right)  =0\\
\Updownarrow\\
\sum_{0\leq P\leq4}P_{\mathcal{H}^{P}}\left(  \lambda_{1}^{2}-\lambda
_{1}+\lambda_{2}^{2}-\lambda_{2}+\lambda_{1}\lambda_{2}+\lambda_{2}\lambda
_{1}\right)  P_{\mathcal{H}^{P}}=0.
\end{gather}
Bringing the operator products to normal form shows that the terms produced by
the one electron operators are
\[
\lambda_{1}^{2}-\lambda_{1}=\sum_{1\leq l\leq r}\alpha_{l}\left(  \alpha
_{l}-1\right)  a_{l}^{\dagger}a_{l}+\sum_{1\leq l_{1}\neq l_{2}\leq r}%
\alpha_{l_{1}}\alpha_{l_{2}}a_{l_{1}}^{\dagger}a_{l_{2}}^{\dagger}a_{l_{2}%
}a_{l_{1}}=\Lambda_{111}+\Lambda_{112}=0,
\]
the terms due to the two electron operators are
\begin{equation}
\lambda_{2}^{2}-\lambda_{2}=\Lambda_{224}+\Lambda_{223}+\Lambda_{222},
\end{equation}
where
\begin{equation}
\Lambda_{224}=36\sum_{\substack{1\leq l_{1}<l_{2}<k_{1}<k_{2}\leq r\\1\leq
l_{3}<l_{4}<k_{3}<k_{4}\leq r\\\sigma,\eta\in S_{4}}}\left(  -1\right)
^{\pi\left(  \sigma\right)  }\left(  -1\right)  ^{\pi\left(  \eta\right)
}\lambda_{2\sigma\left(  l_{1}\right)  \sigma\left(  l_{2}\right)  \eta\left(
l_{3}\right)  \eta\left(  l_{4}\right)  }\lambda_{2\sigma\left(  k_{1}\right)
\sigma\left(  k_{2}\right)  \eta\left(  k_{3}\right)  \eta\left(
k_{4}\right)  }a_{l_{1}}^{\dagger}a_{l_{2}}^{\dagger}a_{k_{1}}^{\dagger
}a_{k_{2}}^{\dagger}a_{l_{4}}a_{l_{3}}a_{k_{4}}a_{k_{3}},
\end{equation}%
\begin{equation}
\Lambda_{223}=9\sum_{\substack{1\leq l_{1}<l_{2}<k_{2}\leq r\\1\leq
l_{3}<k_{3}<k_{4}\leq r}}\sum_{\substack{\sigma,\eta\in S_{3}\\1\leq l_{4}\leq
r}}\left(  -1\right)  ^{\pi\left(  \sigma\right)  }\left(  -1\right)
^{\pi\left(  \eta\right)  }\lambda_{2\sigma\left(  l_{1}\right)  \sigma\left(
l_{2}\right)  \eta\left(  l_{3}\right)  l_{4}}\lambda_{2l_{4}\sigma\left(
k_{2}\right)  \eta\left(  k_{3}\right)  \eta\left(  k_{4}\right)  }a_{l_{1}%
}^{\dagger}a_{l_{2}}^{\dagger}a_{k_{2}}^{\dagger}a_{l_{3}}a_{k_{4}}a_{k_{3}}%
\end{equation}
and
\begin{equation}
\Lambda_{222}=\sum_{\substack{1\leq l_{1}<l_{2}\leq r\\1\leq k_{3}<k_{4}\leq
r}}\left\{  \lambda_{2l_{1}l_{2}k_{3}k_{4}}^{2}-\lambda_{2l_{1}l_{2}k_{3}%
k_{4}}\right\}  a_{l_{1}}^{\dagger}a_{l_{2}}^{\dagger}a_{k_{4}}a_{k_{3}%
}.\nonumber
\end{equation}
The mixed term is the operator
\begin{equation}
\lambda_{1}\lambda_{2}+\lambda_{2}\lambda_{1},
\end{equation}
bringing this operator to normal form produces
\begin{align}
\lambda_{1}\lambda_{2}+\lambda_{2}\lambda_{1}  &  =\frac{1}{2}\sum
_{\substack{1\leq l_{1},l_{2}\leq r\\1\leq l_{3},l_{4}\leq r\\1\leq l\leq
r}}\lambda_{2l_{1}l_{2}l_{3}l_{4}}\alpha_{l}a_{l_{1}}^{\dagger}a_{l_{2}%
}^{\dagger}a_{l}^{\dagger}a_{l}a_{l_{4}}a_{l_{3}}+\frac{1}{2}\sum
_{\substack{1\leq l_{1},l_{2}\leq r\\1\leq l_{3},l_{4}\leq r}}\lambda
_{2l_{1}l_{2}l_{3}l_{4}}\left(  \alpha_{l_{1}}+\alpha_{l_{3}}\right)
a_{l_{1}}^{\dagger}a_{l_{2}}^{\dagger}a_{l_{4}}a_{l_{3}}\nonumber\\
&  =2\sum_{\substack{1\leq l_{1}<l_{2}\leq r\\1\leq l_{3}<l_{4}\leq r\\1\leq
l\leq r}}\lambda_{2l_{1}l_{2}l_{3}l_{4}}\alpha_{l}a_{l_{1}}^{\dagger}a_{l_{2}%
}^{\dagger}a_{l}^{\dagger}a_{l}a_{l_{4}}a_{l_{3}}+\sum_{\substack{1\leq
l_{1}<l_{2}\leq r\\1\leq l_{3}<l_{4}\leq r}}\lambda_{2l_{1}l_{2}l_{3}l_{4}%
}\left(  \alpha_{l_{1}}+\alpha_{l_{2}}+\alpha_{l_{3}}+\alpha_{l_{4}}\right)
a_{l_{1}}^{\dagger}a_{l_{2}}^{\dagger}a_{l_{4}}a_{l_{3}}%
\end{align}
and we can define
\begin{align}
\Lambda_{123}  &  =\,2\sum_{\substack{1\leq l_{1}<l_{2}\leq r\\1\leq
l_{3}<l_{4}\leq r\\1\leq l\leq r}}\lambda_{2l_{1}l_{2}l_{3}l_{4}}\alpha
_{l}a_{l_{1}}^{\dagger}a_{l_{2}}^{\dagger}a_{l}^{\dagger}a_{l}a_{l_{4}%
}a_{l_{3}}\\
\Lambda_{122}  &  =\sum_{\substack{1\leq l_{1}<l_{2}\leq r\\1\leq l_{3}%
<l_{4}\leq r}}\lambda_{2l_{1}l_{2}l_{3}l_{4}}\left(  \alpha_{l_{1}}%
+\alpha_{l_{2}}+\alpha_{l_{3}}+\alpha_{l_{4}}\right)  a_{l_{1}}^{\dagger
}a_{l_{2}}^{\dagger}a_{l_{4}}a_{l_{3}}.
\end{align}
The idempotency condition becomes
\[
\left(  \lambda_{1}+\lambda_{2}\right)  ^{2}-\left(  \lambda_{1}+\lambda
_{2}\right)  =\Lambda_{111}+\Lambda_{112}+\Lambda_{224}+\Lambda_{223}%
+\Lambda_{222}+\Lambda_{122}+\Lambda_{123}=0
\]
and
\begin{subequations}
\begin{align}
\Lambda_{111}  &  =0\\
\Lambda_{112}+\Lambda_{222}+\Lambda_{122}  &  =0\\
\Lambda_{223}+\Lambda_{123}  &  =0\\
\Lambda_{224}  &  =0.
\end{align}
Using the FQ\ equivalence this also gives
\end{subequations}
\begin{subequations}
\begin{align}
P_{\mathcal{H}^{1}}\Lambda_{111}P_{\mathcal{H}^{1}}  &  =0\\
P_{\mathcal{H}^{2}}\Lambda_{112}P_{\mathcal{H}^{2}}+P_{\mathcal{H}^{2}}%
\Lambda_{222}P_{\mathcal{H}^{2}}+P_{\mathcal{H}^{2}}\Lambda_{122}%
P_{\mathcal{H}^{2}}  &  =0\\
P_{\mathcal{H}^{3}}\Lambda_{223}P_{\mathcal{H}^{3}}+P_{\mathcal{H}^{3}}%
\Lambda_{123}P_{\mathcal{H}^{3}}  &  =0\\
P_{\mathcal{H}^{4}}\Lambda_{224}P_{\mathcal{H}^{4}}  &  =0.
\end{align}


Collecting together the operators of the same order one obtains by considering
linear dependence
\end{subequations}
\begin{subequations}
\begin{align}
&  \Lambda_{111}=\sum_{1\leq l\leq r}\alpha_{l}\left(  \alpha_{l}-1\right)
a_{l}^{\dagger}a_{l}=0\\
&  \Lambda_{112}+\Lambda_{222}+\Lambda_{122}=2\sum_{1\leq l_{1}<l_{2}\leq
r}\alpha_{l_{1}}\alpha_{l_{2}}a_{l_{1}}^{\dagger}a_{l_{2}}^{\dagger}a_{l_{2}%
}a_{l_{1}}+\sum_{\substack{1\leq l_{1}<l_{2}\leq r\\1\leq k_{3}<k_{4}\leq
r}}\left\{  \lambda_{2l_{1}l_{2}k_{3}k_{4}}^{2}-\lambda_{2l_{1}l_{2}k_{3}%
k_{4}}\right\}  a_{l_{1}}^{\dagger}a_{l_{2}}^{\dagger}a_{k_{4}}a_{k_{3}}\\
&  +\sum_{\substack{1\leq l_{1}<l_{2}\leq r\\1\leq l_{3}<l_{4}\leq r}%
}\lambda_{2l_{1}l_{2}l_{3}l_{4}}\left(  \alpha_{l_{1}}+\alpha_{l_{2}}%
+\alpha_{l_{3}}+\alpha_{l_{4}}\right)  a_{l_{1}}^{\dagger}a_{l_{2}}^{\dagger
}a_{l_{4}}a_{l_{3}}=0\\
&  \Lambda_{223}+\Lambda_{123}=9\sum_{\substack{1\leq l_{1}<l_{2}<k_{2}\leq
r\\1\leq l_{3}<k_{3}<k_{4}\leq r}}\sum_{\substack{\sigma,\eta\in S_{3}\\1\leq
l_{4}\leq r}}\left(  -1\right)  ^{\pi\left(  \sigma\right)  }\left(
-1\right)  ^{\pi\left(  \eta\right)  }\lambda_{2\sigma\left(  l_{1}\right)
\sigma\left(  l_{2}\right)  \eta\left(  l_{3}\right)  l_{4}}\lambda
_{2l_{4}\sigma\left(  k_{2}\right)  \eta\left(  k_{3}\right)  \eta\left(
k_{4}\right)  }a_{l_{1}}^{\dagger}a_{l_{2}}^{\dagger}a_{k_{2}}^{\dagger
}a_{l_{3}}a_{k_{4}}a_{k_{3}}\\
&  +\sum_{\substack{1\leq l_{1}<l_{2}\leq r\\1\leq l_{3}<l_{4}\leq r\\1\leq
l\leq r}}\lambda_{2l_{1}l_{2}l_{3}l_{4}}\alpha_{l}a_{l_{1}}^{\dagger}a_{l_{2}%
}^{\dagger}a_{l}^{\dagger}a_{l}a_{l_{4}}a_{l_{3}}\left.  =\right.  0\\
&  \Lambda_{224}=36\sum_{\substack{1\leq l_{1}<l_{2}<k_{1}<k_{2}\leq r\\1\leq
l_{3}<l_{4}<k_{3}<k_{4}\leq r\\\sigma,\eta\in S_{4}}}\left(  -1\right)
^{\pi\left(  \sigma\right)  }\left(  -1\right)  ^{\pi\left(  \eta\right)
}\lambda_{2\sigma\left(  l_{1}\right)  \sigma\left(  l_{2}\right)  \eta\left(
l_{3}\right)  \eta\left(  l_{4}\right)  }\lambda_{2\sigma\left(  k_{1}\right)
\sigma\left(  k_{2}\right)  \eta\left(  k_{3}\right)  \eta\left(
k_{4}\right)  }a_{l_{1}}^{\dagger}a_{l_{2}}^{\dagger}a_{k_{1}}^{\dagger
}a_{k_{2}}^{\dagger}a_{l_{4}}a_{l_{3}}a_{k_{4}}a_{k_{3}}=0.
\end{align}
One can then use linear independency on the indivual terms of the summation to
obtain%
\end{subequations}
\begin{subequations}
\begin{align}
&  \alpha_{l}\left(  \alpha_{l}-1\right)  a_{l}^{\dagger}a_{l}=0,\\
&  \left\{  2\alpha_{l_{1}}\alpha_{l_{2}}+\left(  \lambda_{2l_{1}l_{2}%
l_{1}l_{2}}^{2}-\lambda_{2l_{1}l_{2}l_{1}l_{2}}\right)  +2\lambda_{2l_{1}%
l_{2}l_{3}l_{4}}\left(  \alpha_{l_{1}}+\alpha_{l_{2}}\right)  \right\}
a_{l_{1}}^{\dagger}a_{l_{2}}^{\dagger}a_{l_{2}}a_{l_{1}}=0,\\
&  \left\{  \lambda_{2l_{1}l_{2}k_{3}k_{4}}^{2}-\lambda_{2l_{1}l_{2}k_{3}%
k_{4}}+\lambda_{2l_{1}l_{2}l_{3}l_{4}}\left(  \alpha_{l_{1}}+\alpha_{l_{2}%
}+\alpha_{l_{3}}+\alpha_{l_{4}}\right)  \right\}  a_{l_{1}}^{\dagger}a_{l_{2}%
}^{\dagger}a_{k_{4}}a_{k_{3}}=0,l_{1},l_{2}\neq k_{3},k_{4}\\
&  \sum_{\substack{\sigma,\eta\in S_{3}\\1\leq l_{4}\leq r}}\left(  -1\right)
^{\pi\left(  \sigma\right)  }\left(  -1\right)  ^{\pi\left(  \eta\right)
}\lambda_{2\sigma\left(  l_{1}\right)  \sigma\left(  l_{2}\right)  \eta\left(
l_{3}\right)  l_{4}}\lambda_{2l_{4}\sigma\left(  l\right)  \eta\left(
l_{3}\right)  \eta\left(  l_{4}\right)  }a_{l_{1}}^{\dagger}a_{l_{2}}%
^{\dagger}a_{l}^{\dagger}a_{l}a_{l_{4}}a_{l_{3}}+\lambda_{2l_{1}l_{2}%
l_{3}l_{4}}\alpha_{l}a_{l_{1}}^{\dagger}a_{l_{2}}^{\dagger}a_{l}^{\dagger
}a_{l}a_{l_{4}}a_{l_{3}}\left.  =\right.  0\\
&  \sum_{\sigma,\eta\in S_{3}}\left(  -1\right)  ^{\pi\left(  \sigma\right)
}\left(  -1\right)  ^{\pi\left(  \eta\right)  }\lambda_{2\sigma\left(
l_{1}\right)  \sigma\left(  l_{2}\right)  \eta\left(  l_{3}\right)
\eta\left(  l_{4}\right)  }\lambda_{2\sigma\left(  k_{1}\right)  \sigma\left(
k_{2}\right)  \eta\left(  k_{3}\right)  \eta\left(  k_{4}\right)  }a_{l_{1}%
}^{\dagger}a_{l_{2}}^{\dagger}a_{k_{1}}^{\dagger}a_{k_{2}}^{\dagger}a_{l_{4}%
}a_{l_{3}}a_{k_{4}}a_{k_{3}}=0
\end{align}


\subsubsection{Relation to First Quantization}

First one can note that
\end{subequations}
\begin{equation}
\left(  \lambda_{1}+\lambda_{2}\right)  ^{2}=\lambda_{1}+\lambda
_{2}\Rightarrow\lambda_{1}+\lambda_{2}\geq0,
\end{equation}
which shows
\begin{equation}
\hat{\lambda}_{1}\geq0\Leftrightarrow\lambda_{1}\geq0,
\end{equation}
where
\begin{equation}
\hat{\lambda}_{1}=\Gamma^{-1}\left(  \lambda_{1}\right)
\end{equation}
and
\begin{equation}
\lambda_{2}=\lambda_{2}^{\dagger}.
\end{equation}
Then noting that
\begin{align}
\lambda_{1}+\lambda_{2}  &  =\Gamma\left(  \hat{\lambda}_{1}+\hat{\lambda}%
_{2}\right)  ;\hat{\lambda}_{1}\in\mathcal{B}\left(  \mathcal{H}^{1}\right)
,\hat{\lambda}_{2}\in\mathcal{B}\left(  \mathcal{H}^{2}\right) \nonumber\\
\hat{\lambda}_{1}+\hat{\lambda}_{2}  &  =\Gamma^{-1}\left(  \lambda
_{1}+\lambda_{2}\right)
\end{align}
one obtains%
\begin{align}
&  \left(  \lambda_{1}\right)  ^{2}+\lambda_{1}\lambda_{2}+\lambda_{2}%
\lambda_{1}+\left(  \lambda_{2}\right)  ^{2}=\nonumber\\
&  \Gamma\left(  \hat{\lambda}_{1}^{2}\right)  +\Gamma\left(  \hat{\lambda
}_{1}\wedge\hat{\lambda}_{1}\right)  +2\Gamma\left(  \hat{\lambda}_{1}%
\wedge\hat{\lambda}_{2}\right)  +2\Gamma\left(  L_{3}^{2}\left(  \hat{\lambda
}_{1}\wedge\hat{\lambda}_{2}\right)  \right)  +\Gamma\left(  \hat{\lambda}%
_{2}\wedge\hat{\lambda}_{2}\right)  +\Gamma\left(  L_{3}^{4}\left(
\hat{\lambda}_{2}\wedge\hat{\lambda}_{2}\right)  \right)  +\Gamma\left(
\hat{\lambda}_{2}^{2}\right)  =\nonumber\\
&  \Gamma\left(  \hat{\lambda}_{1}^{2}+\hat{\lambda}_{1}\wedge\hat{\lambda
}_{1}+2\hat{\lambda}_{1}\wedge\hat{\lambda}_{2}+2L_{3}^{2}\left(  \hat
{\lambda}_{1}\wedge\hat{\lambda}_{2}\right)  +\hat{\lambda}_{2}\wedge
\hat{\lambda}_{2}+L_{3}^{4}\left(  \hat{\lambda}_{2}\wedge\hat{\lambda}%
_{2}\right)  +\hat{\lambda}_{2}^{2}\right)
\end{align}
This gives
\begin{gather}
\left(  \lambda_{1}+\lambda_{2}\right)  ^{2}-\left(  \lambda_{1}+\lambda
_{2}\right)  =\nonumber\\
\Gamma\left(  \hat{\lambda}_{1}^{2}+\hat{\lambda}_{1}\wedge\hat{\lambda}%
_{1}+2\hat{\lambda}_{1}\wedge\hat{\lambda}_{2}+\Gamma_{1}^{2}\left(
\hat{\lambda}_{1}\right)  \hat{\lambda}_{2}+\hat{\lambda}_{2}\Gamma_{1}%
^{2}\left(  \hat{\lambda}_{1}\right)  \right.  =0\nonumber\\
\left.  +\hat{\lambda}_{2}\wedge\hat{\lambda}_{2}+L_{4}^{3}\left(
\hat{\lambda}_{2}\wedge\hat{\lambda}_{2}\right)  +\hat{\lambda}_{2}^{2}%
-\hat{\lambda}_{1}-\hat{\lambda}_{2}\right) \\
\Updownarrow\nonumber
\end{gather}%
\begin{subequations}
\begin{align}
\hat{\lambda}_{1}^{2}-\hat{\lambda}_{1}  &  =0\label{1ord1}\\
\hat{\lambda}_{1}\wedge\hat{\lambda}_{1}+\hat{\lambda}_{2}^{2}-\hat{\lambda
}_{2}+\Gamma_{1}^{2}\left(  \hat{\lambda}_{1}\right)  \hat{\lambda}_{2}%
+\hat{\lambda}_{2}\Gamma_{1}^{2}\left(  \hat{\lambda}_{1}\right)   &
=0\label{2ord2}\\
2\hat{\lambda}_{1}\wedge\hat{\lambda}_{2}+L_{4}^{3}\left(  \hat{\lambda}%
_{2}\wedge\hat{\lambda}_{2}\right)   &  =0\Rightarrow\hat{\lambda}_{1}%
\wedge\hat{\lambda}_{2}=0\label{3ord3}\\
\hat{\lambda}_{2}\wedge\hat{\lambda}_{2}  &  =0, \label{4ord4}%
\end{align}
which can be expressed as
\end{subequations}
\begin{subequations}
\begin{align}
P_{\mathcal{H}^{1}}\Lambda_{111}P_{\mathcal{H}^{1}}  &  =0\label{1ord}\\
P_{\mathcal{H}^{2}}\Lambda_{112}P_{\mathcal{H}^{2}}+P_{\mathcal{H}^{2}}%
\Lambda_{222}P_{\mathcal{H}^{2}}+P_{\mathcal{H}^{2}}\Lambda_{122}%
P_{\mathcal{H}^{2}}  &  =0\label{2ord}\\
P_{\mathcal{H}^{3}}\Lambda_{223}P_{\mathcal{H}^{3}}+P_{\mathcal{H}^{3}}%
\Lambda_{123}P_{\mathcal{H}^{3}}  &  =0\label{3ord}\\
P_{\mathcal{H}^{4}}\Lambda_{224}P_{\mathcal{H}^{4}}  &  =0. \label{4ord}%
\end{align}


The condition eq. $\left(  \ref{1ord}\right)  $
\end{subequations}
\begin{equation}
\hat{\lambda}_{1}^{2}-\hat{\lambda}_{1}=0
\end{equation}
shows that $\hat{\lambda}_{1}$ has eigenvalues $\left\{  0,1\right\}  $ and
$\hat{\lambda}_{1}$ is a projector which can be expanded as
\begin{equation}
\hat{\lambda}_{1}=P_{\nu}=\sum_{1\leq j\leq v}\left\vert \psi_{j}\right\rangle
\left\langle \psi_{j}\right\vert ;\nu\leq r
\end{equation}
for any basis that spans the range of $\hat{\lambda}_{1}$.

The conditions eq. $\left(  \ref{4ord}\right)  $
\begin{equation}
\hat{\lambda}_{2}\wedge\hat{\lambda}_{2}=0
\end{equation}
and eq. $\left(  \ref{3ord}\right)  $ can only hold if
\begin{equation}
\hat{\lambda}_{1}\wedge\hat{\lambda}_{2}=0.
\end{equation}


The condition $\hat{\lambda}_{1}\wedge\hat{\lambda}_{2}=0$ shows that
\begin{equation}
\hat{\lambda}_{1}\wedge\hat{\lambda}_{2}=\sum_{1\leq j\leq v}\left\vert
\psi_{j}\right\rangle \left\langle \psi_{j}\right\vert \wedge\sum
_{\substack{1\leq k<l\leq r\\1\leq m<n\leq r}}\lambda_{klmn}\left\vert
\psi_{k}\psi_{l}\right\rangle \left\langle \psi_{m}\psi_{n}\right\vert
=\sum_{\substack{1\leq k<l\leq r\\1\leq m<n\leq r\\1\leq j\leq v}%
}\lambda_{klmn}\left\vert \psi_{j}\psi_{k}\psi_{l}\right\rangle \left\langle
\psi_{j}\psi_{m}\psi_{n}\right\vert =0.
\end{equation}


\begin{enumerate}
\item \bigskip Let
\begin{equation}
\hat{\lambda}_{1}=\left\vert \psi_{1}\right\rangle \left\langle \psi
_{1}\right\vert
\end{equation}
then
\begin{gather}
\hat{\lambda}_{1}\wedge\hat{\lambda}_{2}=0\\
\Downarrow\\
\hat{\lambda}_{2}=\sum_{\substack{2\leq j<k\leq r\\2\leq l\leq r}%
}\lambda_{1jkl}\left\vert \psi_{1}\psi_{l}\right\rangle \left\langle \psi
_{j}\psi_{k}\right\vert +\sum_{\substack{2\leq j<k\leq r\\2\leq l\leq
r}}\lambda_{jk1l}\left\vert \psi_{j}\psi_{k}\right\rangle \left\langle
\psi_{1}\psi_{l}\right\vert +\sum_{2\leq j,k\leq r}\lambda_{1j1k}\left\vert
\psi_{1}\psi_{j}\right\rangle \left\langle \psi_{1}\psi_{k}\right\vert
\end{gather}
and
\begin{gather}
\hat{\lambda}_{1}\wedge\hat{\lambda}_{1}+\hat{\lambda}_{2}^{2}-\hat{\lambda
}_{2}+\Gamma_{1}^{2}\left(  \hat{\lambda}_{1}\right)  \hat{\lambda}_{2}%
+\hat{\lambda}_{2}\Gamma_{1}^{2}\left(  \hat{\lambda}_{1}\right)  =0\\
\Downarrow\\
\hat{\lambda}_{2}^{2}-\hat{\lambda}_{2}+\sum_{2\leq i<r}\left\vert \psi
_{1}\psi_{i}\right\rangle \left\langle \psi_{1}\psi_{i}\right\vert
\hat{\lambda}_{2}+\hat{\lambda}_{2}\sum_{2\leq i<r}\left\vert \psi_{1}\psi
_{i}\right\rangle \left\langle \psi_{1}\psi_{i}\right\vert =0
\end{gather}
giving%
\begin{align}
&  \hat{\lambda}_{2}^{2}-\hat{\lambda}_{2}+\sum_{2\leq i<r}\left\vert \psi
_{1}\psi_{i}\right\rangle \left\langle \psi_{1}\psi_{i}\right\vert \left(
\sum_{2\leq j,l\leq r}\lambda_{1jkl}\left\vert \psi_{1}\psi_{j}\right\rangle
\left\langle \psi_{1}\psi_{l}\right\vert +\sum_{\substack{2\leq k<l\leq
r\\2\leq j\leq r}}\lambda_{1jkl}\left\vert \psi_{1}\psi_{j}\right\rangle
\left\langle \psi_{k}\psi_{l}\right\vert \right)  +\\
&  \left(  \sum_{2\leq j,l\leq r}\lambda_{1jkl}\left\vert \psi_{1}\psi
_{j}\right\rangle \left\langle \psi_{1}\psi_{l}\right\vert +\sum
_{\substack{2\leq k<l\leq r\\2\leq j\leq r}}\lambda_{kl1j}\left\vert \psi
_{k}\psi_{l}\right\rangle \left\langle \psi_{1}\psi_{j}\right\vert \right)
\sum_{2\leq i<r}\left\vert \psi_{1}\psi_{i}\right\rangle \left\langle \psi
_{1}\psi_{i}\right\vert =\\
&  \hat{\lambda}_{2}^{2}-\hat{\lambda}_{2}+\hat{\lambda}_{2}+\hat{\lambda}%
_{2}=\hat{\lambda}_{2}^{2}+\hat{\lambda}_{2}=0\Rightarrow\hat{\lambda}_{2}%
^{2}=-\hat{\lambda}_{2}.
\end{align}
As
\begin{equation}
\hat{\lambda}_{2}^{2}=\sum_{1\leq j\leq\nu}\beta_{j}^{2}\left\vert
g_{j}\right\rangle \left\langle g_{j}\right\vert =-\sum_{1\leq j\leq\nu}%
\beta_{j}\left\vert g_{j}\right\rangle \left\langle g_{j}\right\vert
\end{equation}%
\begin{equation}
\beta_{j}=-1\text{ or }0.
\end{equation}
Let $\nu=2$ then using $\hat{\lambda}_{2}\wedge\hat{\lambda}_{2}=0$ gives
\begin{equation}
\left\vert g_{1}\right\rangle \left\langle g_{1}\right\vert \wedge\left\vert
g_{1}\right\rangle \left\langle g_{1}\right\vert +2\left\vert g_{1}%
\right\rangle \left\langle g_{1}\right\vert \wedge\left\vert g_{2}%
\right\rangle \left\langle g_{2}\right\vert +\left\vert g_{2}\right\rangle
\left\langle g_{2}\right\vert \wedge\left\vert g_{2}\right\rangle \left\langle
g_{2}\right\vert =0,
\end{equation}
which can expressed as
\begin{equation}
\left\vert g_{1}^{2}\right\rangle \left\langle g_{1}^{2}\right\vert
+2\left\vert g_{1}\wedge g_{2}\right\rangle \left\langle g_{1}\wedge
g_{2}\right\vert +\left\vert g_{2}^{2}\right\rangle \left\langle g_{2}%
^{2}\right\vert =0
\end{equation}
and as each term is a positive operator
\begin{align}
\left\vert g_{1}^{2}\right\rangle \left\langle g_{1}^{2}\right\vert  &  =0\\
\left\vert g_{1}\wedge g_{2}\right\rangle \left\langle g_{1}\wedge
g_{2}\right\vert  &  =0\\
\left\vert g_{2}^{2}\right\rangle \left\langle g_{2}^{2}\right\vert  &  =0.
\end{align}
This further implies (by taking expectation values wrt all states in
$\mathcal{H}^{2})$ that
\begin{align}
\left\vert g_{1}^{2}\right\rangle  &  =0\\
\left\vert g_{1}\wedge g_{2}\right\rangle  &  =0\\
\left\vert g_{2}^{2}\right\rangle  &  =0.
\end{align}
Hence the geminals $\left\vert g_{1}\right\rangle $ and $\left\vert
g_{2}\right\rangle $ must have rank equal to one i.e
\begin{align}
\left\vert g_{1}\right\rangle  &  =\left\vert \varphi_{1}\varphi
_{2}\right\rangle \\
\left\vert g_{2}\right\rangle  &  =\left\vert \psi_{1}\psi_{2}\right\rangle
\end{align}
and in order that
\begin{equation}
\left\vert \varphi_{1}\wedge\varphi_{2}\wedge\psi_{1}\wedge\psi_{2}%
\right\rangle =0
\end{equation}
the vectors $\left\{  \varphi_{1},\varphi_{2},\psi_{1},\psi_{2}\right\}  $
must be linearly \emph{dependent}, so their gram determinant
\begin{equation}
\left\vert
\begin{array}
[c]{cccc}%
\left\langle \left.  \varphi_{1}\right\vert \varphi_{1}\right\rangle  &
\left\langle \left.  \varphi_{1}\right\vert \varphi_{2}\right\rangle  &
\left\langle \left.  \varphi_{1}\right\vert \psi_{1}\right\rangle  &
\left\langle \left.  \varphi_{1}\right\vert \psi_{2}\right\rangle \\
\left\langle \left.  \varphi_{2}\right\vert \varphi_{1}\right\rangle  &
\left\langle \left.  \varphi_{2}\right\vert \varphi_{2}\right\rangle  &
\left\langle \left.  \varphi_{2}\right\vert \psi_{1}\right\rangle  &
\left\langle \left.  \varphi_{2}\right\vert \psi_{2}\right\rangle \\
\left\langle \left.  \psi_{1}\right\vert \varphi_{1}\right\rangle  &
\left\langle \left.  \psi_{1}\right\vert \varphi_{2}\right\rangle  &
\left\langle \left.  \psi_{1}\right\vert \psi_{1}\right\rangle  & \left\langle
\left.  \psi_{1}\right\vert \psi_{2}\right\rangle \\
\left\langle \left.  \psi_{2}\right\vert \varphi_{1}\right\rangle  &
\left\langle \left.  \psi_{2}\right\vert \varphi_{2}\right\rangle  &
\left\langle \left.  \psi_{2}\right\vert \psi_{1}\right\rangle  & \left\langle
\left.  \psi_{2}\right\vert \psi_{2}\right\rangle
\end{array}
\right\vert =0.
\end{equation}
The geminals $\left\vert g_{1}\right\rangle $ and $\left\vert g_{2}%
\right\rangle $ are orthogonal thus
\begin{equation}
\left\langle \left.  \varphi_{1}\wedge\varphi_{2}\right\vert \psi_{1}%
\wedge\psi_{2}\right\rangle =0,
\end{equation}
which can expanded as
\begin{equation}
\left\vert
\begin{array}
[c]{cc}%
\left\langle \left.  \varphi_{1}\right\vert \psi_{1}\right\rangle  &
\left\langle \left.  \varphi_{1}\right\vert \psi_{2}\right\rangle \\
\left\langle \left.  \varphi_{2}\right\vert \psi_{1}\right\rangle  &
\left\langle \left.  \varphi_{2}\right\vert \psi_{2}\right\rangle
\end{array}
\right\vert =\left\langle \left.  \varphi_{1}\right\vert \psi_{1}\right\rangle
\left\langle \left.  \varphi_{2}\right\vert \psi_{2}\right\rangle
-\left\langle \left.  \varphi_{1}\right\vert \psi_{2}\right\rangle
\left\langle \left.  \varphi_{2}\right\vert \psi_{1}\right\rangle =0.
\end{equation}
As the set $\left\{  \varphi_{1},\varphi_{2},\psi_{1},\psi_{2}\right\}  $ we
can span the same space by $\left\{  \varphi_{1},\varphi_{2},\psi_{2}\right\}
$ and obtain
\begin{equation}
\left\vert
\begin{array}
[c]{cc}%
\left\langle \left.  \varphi_{1}\right\vert \varphi_{1}\right\rangle  &
\left\langle \left.  \varphi_{1}\right\vert \psi_{2}\right\rangle \\
\left\langle \left.  \varphi_{2}\right\vert \varphi_{1}\right\rangle  &
\left\langle \left.  \varphi_{2}\right\vert \psi_{2}\right\rangle
\end{array}
\right\vert =\left\langle \left.  \varphi_{1}\right\vert \varphi
_{1}\right\rangle \left\langle \left.  \varphi_{2}\right\vert \psi
_{2}\right\rangle -\left\langle \left.  \varphi_{1}\right\vert \psi
_{2}\right\rangle \left\langle \left.  \varphi_{2}\right\vert \varphi
_{1}\right\rangle =0.
\end{equation}
As%
\begin{equation}
\hat{\lambda}_{2}=\left\vert \varphi_{1}\varphi_{2}\right\rangle \left\langle
\varphi_{1}\varphi_{2}\right\vert +\left\vert \psi_{1}\psi_{2}\right\rangle
\left\langle \psi_{1}\psi_{2}\right\vert =\left\vert \varphi_{1}\varphi
_{2}\right\rangle \left\langle \varphi_{1}\varphi_{2}\right\vert +\left\vert
\psi_{1}\varphi_{2}\right\rangle \left\langle \psi_{1}\varphi_{2}\right\vert
\end{equation}%
\begin{equation}
\hat{\lambda}_{1}=\left\vert \varphi_{2}\right\rangle \left\langle \varphi
_{2}\right\vert
\end{equation}
satisfies
\begin{equation}
\hat{\lambda}_{1}\wedge\hat{\lambda}_{2}=0.
\end{equation}

\end{enumerate}

Hence
\begin{align}
&  \hat{\lambda}_{2}^{2}=\\
&  \left(  \sum_{\substack{2\leq j<k\leq r\\2\leq l\leq r}}\lambda
_{1jkl}\left\vert \psi_{1}\psi_{l}\right\rangle \left\langle \psi_{j}\psi
_{k}\right\vert +\sum_{\substack{2\leq j<k\leq r\\2\leq l\leq r}%
}\lambda_{jk1l}\left\vert \psi_{j}\psi_{k}\right\rangle \left\langle \psi
_{1}\psi_{l}\right\vert +\sum_{2\leq j,k\leq r}\lambda_{1j1k}\left\vert
\psi_{1}\psi_{j}\right\rangle \left\langle \psi_{1}\psi_{k}\right\vert \right)
\\
&  \left(  \sum_{\substack{2\leq j<k\leq r\\2\leq l\leq r}}\lambda
_{1jkl}\left\vert \psi_{1}\psi_{l}\right\rangle \left\langle \psi_{j}\psi
_{k}\right\vert +\sum_{\substack{2\leq j<k\leq r\\2\leq l\leq r}%
}\lambda_{jk1l}\left\vert \psi_{j}\psi_{k}\right\rangle \left\langle \psi
_{1}\psi_{l}\right\vert +\sum_{2\leq j,k\leq r}\lambda_{1j1k}\left\vert
\psi_{1}\psi_{j}\right\rangle \left\langle \psi_{1}\psi_{k}\right\vert
\right)  \times\\
&  =\sum_{\substack{2\leq j<k\leq r\\2\leq m<n\leq r\\2\leq l\leq r}%
}\lambda_{jk1l}\lambda_{1lmn}\left\vert \psi_{j}\psi_{k}\right\rangle
\left\langle \psi_{m}\psi_{n}\right\vert +\sum_{\substack{2\leq j<k\leq
r\\2\leq l,m\leq r}}\lambda_{jk1l}\lambda_{1l1m}\left\vert \psi_{j}\psi
_{k}\right\rangle \left\langle \psi_{1}\psi_{m}\right\vert \\
&  +\sum_{\substack{2\leq j,l\leq r\\2\leq m<n\leq r}}\lambda_{1j1l}%
\lambda_{1lmn}\left\vert \psi_{1}\psi_{j}\right\rangle \left\langle \psi
_{m}\psi_{n}\right\vert +\sum_{2\leq j,k\leq r}\lambda_{1j1m}^{2}\left\vert
\psi_{1}\psi_{j}\right\rangle \left\langle \psi_{1}\psi_{m}\right\vert
+\sum_{\substack{2\leq j<k\leq r\\2\leq l,m\leq r}}\lambda_{1ljk}%
\lambda_{jk1m}\left\vert \psi_{1}\psi_{l}\right\rangle \left\langle \psi
_{1}\psi_{m}\right\vert \\
&  =-\sum_{\substack{2\leq j<k\leq r\\2\leq l\leq r}}\lambda_{1jkl}\left\vert
\psi_{1}\psi_{l}\right\rangle \left\langle \psi_{j}\psi_{k}\right\vert
-\sum_{\substack{2\leq j<k\leq r\\2\leq l\leq r}}\lambda_{jk1l}\left\vert
\psi_{j}\psi_{k}\right\rangle \left\langle \psi_{1}\psi_{l}\right\vert
-\sum_{2\leq j,k\leq r}\lambda_{1j1m}\left\vert \psi_{1}\psi_{j}\right\rangle
\left\langle \psi_{1}\psi_{m}\right\vert .
\end{align}


\begin{enumerate}
\item Let
\begin{equation}
\hat{\lambda}_{1}=\left\vert \psi_{1}\right\rangle \left\langle \psi
_{1}\right\vert +\left\vert \psi_{2}\right\rangle \left\langle \psi
_{2}\right\vert
\end{equation}
then
\begin{align}
\hat{\lambda}_{2}  &  =\sum_{3\leq i<k\leq r}\lambda_{12ik}\left\vert \psi
_{1}\psi_{2}\right\rangle \left\langle \psi_{i}\psi_{k}\right\vert
+\sum_{3\leq i,k\leq r}\lambda_{1i2k}\left\vert \psi_{1}\psi_{i}\right\rangle
\left\langle \psi_{2}\psi_{k}\right\vert +\sum_{3\leq i,k\leq r}\lambda
_{2i1k}\left\vert \psi_{2}\psi_{i}\right\rangle \left\langle \psi_{1}\psi
_{k}\right\vert +\sum_{3\leq i,k\leq r}\lambda_{ik12}\left\vert \psi_{i}%
\psi_{k}\right\rangle \left\langle \psi_{1}\psi_{2}\right\vert \nonumber\\
&  +\sum_{3\leq i<r}\lambda_{121i}\left\vert \psi_{1}\psi_{2}\right\rangle
\left\langle \psi_{1}\psi_{i}\right\vert +\sum_{3\leq i<r}\lambda
_{1i12}\left\vert \psi_{1}\psi_{i}\right\rangle \left\langle \psi_{1}\psi
_{2}\right\vert +\sum_{3\leq i<r}\lambda_{122i}\left\vert \psi_{1}\psi
_{2}\right\rangle \left\langle \psi_{2}\psi_{i}\right\vert +\sum_{3\leq
i<r}\lambda_{2i12}\left\vert \psi_{2}\psi_{i}\right\rangle \left\langle
\psi_{1}\psi_{2}\right\vert \nonumber\\
&  +\lambda_{1212}\left\vert \psi_{1}\psi_{2}\right\rangle \left\langle
\psi_{1}\psi_{2}\right\vert .
\end{align}%
\begin{align}
\left(  \hat{\lambda}_{2}\right)  ^{2}  &  =\sum_{3\leq i<k\leq r}%
\lambda_{12ik}\left\vert \psi_{1}\psi_{2}\right\rangle \left\langle \psi
_{i}\psi_{k}\right\vert +\sum_{3\leq i,k\leq r}\lambda_{1i2k}\left\vert
\psi_{1}\psi_{i}\right\rangle \left\langle \psi_{2}\psi_{k}\right\vert
+\sum_{3\leq i,k\leq r}\lambda_{2i1k}\left\vert \psi_{2}\psi_{i}\right\rangle
\left\langle \psi_{1}\psi_{k}\right\vert +\sum_{3\leq i,k\leq r}\lambda
_{ik12}\left\vert \psi_{i}\psi_{k}\right\rangle \left\langle \psi_{1}\psi
_{2}\right\vert \nonumber\\
&  +\sum_{3\leq i<r}\lambda_{121i}\left\vert \psi_{1}\psi_{2}\right\rangle
\left\langle \psi_{1}\psi_{i}\right\vert +\sum_{3\leq i<r}\lambda
_{1i12}\left\vert \psi_{1}\psi_{i}\right\rangle \left\langle \psi_{1}\psi
_{2}\right\vert +\sum_{3\leq i<r}\lambda_{122i}\left\vert \psi_{1}\psi
_{2}\right\rangle \left\langle \psi_{2}\psi_{i}\right\vert +\sum_{3\leq
i<r}\lambda_{2i12}\left\vert \psi_{2}\psi_{i}\right\rangle \left\langle
\psi_{1}\psi_{2}\right\vert \nonumber\\
&  +\lambda_{1212}\left\vert \psi_{1}\psi_{2}\right\rangle \left\langle
\psi_{1}\psi_{2}\right\vert .
\end{align}
The condition $\hat{\lambda}_{2}\wedge\hat{\lambda}_{2}=0$ gives
\begin{align}
\hat{\lambda}_{2}\wedge\hat{\lambda}_{2}  &  =\sum_{3\leq i,k\leq r}%
\lambda_{12ik}\left\vert \psi_{1}\psi_{2}\right\rangle \left\langle \psi
_{i}\psi_{k}\right\vert \wedge\sum_{3\leq i,k\leq r}\lambda_{jl12}\left\vert
\psi_{j}\psi_{l}\right\rangle \left\langle \psi_{1}\psi_{2}\right\vert
+\sum_{3\leq i,k\leq r}\lambda_{1i2k}\left\vert \psi_{1}\psi_{i}\right\rangle
\left\langle \psi_{2}\psi_{k}\right\vert \wedge\sum_{3\leq j,l\leq r}%
\lambda_{2j1l}\left\vert \psi_{2}\psi_{j}\right\rangle \left\langle \psi
_{1}\psi_{l}\right\vert \nonumber\\
&  =\sum_{\substack{3\leq i<k\leq r\\3\leq j<l\leq r}}\lambda_{12ik}%
\lambda_{jl12}\left\vert \psi_{1}\wedge\psi_{2}\wedge\psi_{j}\wedge\psi
_{l}\right\rangle \left\langle \psi_{1}\wedge\psi_{2}\wedge\psi_{i}\wedge
\psi_{k}\right\vert +\sum_{3\leq i,j,k,l\leq r}\lambda_{1i2k}\lambda
_{2j1l}\left\vert \psi_{1}\wedge\psi_{2}\wedge\psi_{i}\wedge\psi
_{j}\right\rangle \left\langle \psi_{1}\wedge\psi_{2}\wedge\psi_{k}\wedge
\psi_{l}\right\vert =0\\
&
\end{align}%
\begin{equation}
P_{2}\wedge P_{2}+\hat{\lambda}_{2}^{2}-\hat{\lambda}_{2}+\left(  P_{r}\wedge
P_{2}\right)  \hat{\lambda}_{2}+\hat{\lambda}_{2}\left(  P_{r}\wedge
P_{2}\right)  =0
\end{equation}%
\begin{align}
P_{2}\wedge P_{2}  &  =2\left\vert \psi_{1}\psi_{2}\right\rangle \left\langle
\psi_{1}\psi_{2}\right\vert \\
P_{r}\wedge P_{2}  &  =\sum_{3\leq i,k\leq r}\left\{  \left\vert \psi_{1}%
\psi_{i}\right\rangle \left\langle \psi_{1}\psi_{k}\right\vert +\left\vert
\psi_{2}\psi_{i}\right\rangle \left\langle \psi_{2}\psi_{k}\right\vert
\right\}  +2\left\vert \psi_{1}\psi_{2}\right\rangle \left\langle \psi_{1}%
\psi_{2}\right\vert \\
\left(  P_{r}\wedge P_{2}\right)  \hat{\lambda}_{2}  &  =2\sum_{3\leq i<k\leq
r}\lambda_{12ik}\left\vert \psi_{1}\psi_{2}\right\rangle \left\langle \psi
_{i}\psi_{k}\right\vert +2\sum_{3\leq i<r}\lambda_{121i}\left\vert \psi
_{1}\psi_{2}\right\rangle \left\langle \psi_{1}\psi_{i}\right\vert
+\lambda_{1212}\left\vert \psi_{1}\psi_{2}\right\rangle \left\langle \psi
_{1}\psi_{2}\right\vert \\
&  +\sum_{\substack{3\leq i\leq r\\1\leq l<k\leq r}}\lambda_{1i2k}\left\vert
\psi_{1}\psi_{i}\right\rangle \left\langle \psi_{l}\psi_{k}\right\vert
+\sum_{\substack{3\leq i\leq r\\1\leq l<k\leq r}}\lambda_{1i2k}\left\vert
\psi_{2}\psi_{i}\right\rangle \left\langle \psi_{l}\psi_{k}\right\vert
\end{align}


\item Let
\begin{equation}
\hat{\lambda}_{1}=\left\vert \psi_{1}\right\rangle \left\langle \psi
_{1}\right\vert +\left\vert \psi_{2}\right\rangle \left\langle \psi
_{2}\right\vert +\left\vert \psi_{3}\right\rangle \left\langle \psi
_{3}\right\vert
\end{equation}
then
\begin{align}
\hat{\lambda}_{2}  &  =\sum_{4\leq i\leq r}\left\{  \lambda_{123i}\left\vert
\psi_{1}\psi_{2}\right\rangle \left\langle \psi_{3}\psi_{i}\right\vert
+\lambda_{132i}\left\vert \psi_{1}\psi_{3}\right\rangle \left\langle \psi
_{2}\psi_{i}\right\vert +\lambda_{231i}\left\vert \psi_{2}\psi_{3}%
\right\rangle \left\langle \psi_{1}\psi_{i}\right\vert +\right. \nonumber\\
&  \left.  \lambda_{3i12}\left\vert \psi_{3}\psi_{i}\right\rangle \left\langle
\psi_{1}\psi_{2}\right\vert +\lambda_{2i13}\left\vert \psi_{2}\psi
_{i}\right\rangle \left\langle \psi_{1}\psi_{3}\right\vert +\lambda
_{1i23}\left\vert \psi_{1}\psi_{i}\right\rangle \left\langle \psi_{2}\psi
_{3}\right\vert \right\} \nonumber\\
&  +\lambda_{1213}\left\vert \psi_{1}\psi_{2}\right\rangle \left\langle
\psi_{1}\psi_{3}\right\vert +\lambda_{1312}\left\vert \psi_{1}\psi
_{3}\right\rangle \left\langle \psi_{1}\psi_{2}\right\vert +\lambda
_{1223}\left\vert \psi_{1}\psi_{2}\right\rangle \left\langle \psi_{2}\psi
_{3}\right\vert \nonumber\\
&  +\lambda_{2312}\left\vert \psi_{2}\psi_{3}\right\rangle \left\langle
\psi_{1}\psi_{2}\right\vert +\lambda_{2313}\left\vert \psi_{2}\psi
_{3}\right\rangle \left\langle \psi_{1}\psi_{3}\right\vert +\lambda
_{1323}\left\vert \psi_{1}\psi_{3}\right\rangle \left\langle \psi_{2}\psi
_{3}\right\vert .
\end{align}
In particular note that in this case $\hat{\lambda}_{2}$ \emph{has no diagonal
elements}. The condition $\hat{\lambda}_{2}\wedge\hat{\lambda}_{2}=0$ gives
\begin{align}
&  \hat{\lambda}_{2}\wedge\hat{\lambda}_{2}=\nonumber\\
&  \sum_{1\leq i\leq r}\left\{  \left\vert \lambda_{123i}\right\vert
^{2}\left\vert \psi_{1}\psi_{2}\psi_{3}\psi_{i}\right\rangle \left\langle
\psi_{3}\psi_{i}\psi_{1}\psi_{2}\right\vert +\left\vert \lambda_{132i}%
\right\vert ^{2}\left\vert \psi_{1}\psi_{3}\psi_{2}\psi_{i}\right\rangle
\left\langle \psi_{2}\psi_{i}\psi_{1}\psi_{3}\right\vert +\left\vert
\lambda_{231i}\right\vert ^{2}\left\vert \psi_{2}\psi_{3}\psi_{1}\psi
_{i}\right\rangle \left\langle \psi_{1}\psi_{i}\psi_{2}\psi_{3}\right\vert
\right\}  =0.
\end{align}
Thus
\begin{equation}
\lambda_{1234}=\lambda_{1324}=\lambda_{1423}=\lambda_{3412}=\lambda
_{2413}=\lambda_{2314}=0
\end{equation}
giving
\begin{align}
\hat{\lambda}_{2}  &  =\lambda_{1213}\left\vert \psi_{1}\psi_{2}\right\rangle
\left\langle \psi_{1}\psi_{3}\right\vert +\lambda_{1312}\left\vert \psi
_{1}\psi_{3}\right\rangle \left\langle \psi_{1}\psi_{2}\right\vert
+\lambda_{1223}\left\vert \psi_{1}\psi_{2}\right\rangle \left\langle \psi
_{2}\psi_{3}\right\vert \nonumber\\
&  +\lambda_{2312}\left\vert \psi_{2}\psi_{3}\right\rangle \left\langle
\psi_{1}\psi_{2}\right\vert +\lambda_{2313}\left\vert \psi_{2}\psi
_{3}\right\rangle \left\langle \psi_{1}\psi_{3}\right\vert +\lambda
_{1323}\left\vert \psi_{1}\psi_{3}\right\rangle \left\langle \psi_{2}\psi
_{3}\right\vert .
\end{align}


\item Let
\begin{equation}
\hat{\lambda}_{1}=\left\vert \psi_{1}\right\rangle \left\langle \psi
_{1}\right\vert +\left\vert \psi_{2}\right\rangle \left\langle \psi
_{2}\right\vert +\left\vert \psi_{3}\right\rangle \left\langle \psi
_{3}\right\vert +\left\vert \psi_{4}\right\rangle \left\langle \psi
_{4}\right\vert
\end{equation}
then
\begin{align}
\hat{\lambda}_{2}=  &  \lambda_{1234}\left\vert \psi_{1}\psi_{2}\right\rangle
\left\langle \psi_{3}\psi_{4}\right\vert +\lambda_{3412}\left\vert \psi
_{3}\psi_{4}\right\rangle \left\langle \psi_{1}\psi_{2}\right\vert
+\lambda_{1324}\left\vert \psi_{1}\psi_{3}\right\rangle \left\langle \psi
_{2}\psi_{4}\right\vert +\lambda_{2413}\left\vert \psi_{2}\psi_{4}%
\right\rangle \left\langle \psi_{1}\psi_{3}\right\vert \nonumber\\
&  +\lambda_{1423}\left\vert \psi_{1}\psi_{4}\right\rangle \left\langle
\psi_{2}\psi_{3}\right\vert +\lambda_{2314}\left\vert \psi_{2}\psi
_{3}\right\rangle \left\langle \psi_{1}\psi_{4}\right\vert .
\end{align}
Again there are no diagonal elements.
\begin{align}
\hat{\lambda}_{2}\wedge\hat{\lambda}_{2}  &  =\lambda_{1234}\lambda
_{3412}\left\vert \psi_{1}\psi_{2}\right\rangle \left\langle \psi_{3}\psi
_{4}\right\vert \wedge\left\vert \psi_{3}\psi_{4}\right\rangle \left\langle
\psi_{1}\psi_{2}\right\vert +\lambda_{1324}\lambda_{2413}\left\vert \psi
_{1}\psi_{3}\right\rangle \left\langle \psi_{2}\psi_{4}\right\vert
\wedge\left\vert \psi_{2}\psi_{4}\right\rangle \left\langle \psi_{1}\psi
_{3}\right\vert \nonumber\\
&  +\lambda_{1423}\lambda_{2314}\left\vert \psi_{1}\psi_{4}\right\rangle
\left\langle \psi_{2}\psi_{3}\right\vert \wedge\left\vert \psi_{2}\psi
_{3}\right\rangle \left\langle \psi_{1}\psi_{4}\right\vert \nonumber\\
&  =\left(  \lambda_{1234}\lambda_{3412}+\lambda_{1324}\lambda_{2413}%
+\lambda_{1423}\lambda_{2314}\right)  \left\vert \psi_{1}\psi_{2}\psi_{3}%
\psi_{4}\right\rangle \left\langle \psi_{1}\psi_{2}\psi_{3}\psi_{4}\right\vert
\nonumber\\
&  =\left(  \left\vert \lambda_{1234}\right\vert ^{2}+\left\vert
\lambda_{1324}\right\vert ^{2}+\left\vert \lambda_{1423}\right\vert
^{2}\right)  \left\vert \psi_{1}\psi_{2}\psi_{3}\psi_{4}\right\rangle
\left\langle \psi_{1}\psi_{2}\psi_{3}\psi_{4}\right\vert =0.
\end{align}
Thus
\begin{equation}
\lambda_{1234}=\lambda_{1324}=\lambda_{1423}=\lambda_{3412}=\lambda
_{2413}=\lambda_{2314}=0
\end{equation}
and
\begin{equation}
\hat{\lambda}_{2}=0,
\end{equation}
so there are no solutions if $\hat{\lambda}_{1}$ is rank 4.

\item Let
\begin{equation}
\hat{\lambda}_{1}=\left\vert \psi_{1}\right\rangle \left\langle \psi
_{1}\right\vert +\left\vert \psi_{2}\right\rangle \left\langle \psi
_{2}\right\vert +\left\vert \psi_{3}\right\rangle \left\langle \psi
_{3}\right\vert +\left\vert \psi_{4}\right\rangle \left\langle \psi
_{4}\right\vert +\left\vert \psi_{5}\right\rangle \left\langle \psi
_{5}\right\vert
\end{equation}
then there are no solutions i.e no $\hat{\lambda}_{2}$ such that $\hat
{\lambda}_{1}\wedge\hat{\lambda}_{2}=0.$
\end{enumerate}

\bigskip%

\begin{gather}
\hat{\lambda}_{1}\wedge\hat{\lambda}_{2}=P_{\nu}\wedge\hat{\lambda}_{2}%
=\Gamma_{2}^{3}\left(  \hat{\lambda}_{2}\right)  -P_{r-\nu}\wedge\hat{\lambda
}_{2}=0\\
\Downarrow\\
\Gamma_{2}^{3}\left(  \hat{\lambda}_{2}\right)  =P_{r}\wedge\hat{\lambda}%
_{2}=P_{r-\nu}\wedge\hat{\lambda}_{2}%
\end{gather}


\bigskip%

\begin{equation}
\hat{\lambda}_{1}\wedge\hat{\lambda}_{1}+\hat{\lambda}_{2}^{2}-\hat{\lambda
}_{2}+\Gamma_{1}^{2}\left(  \hat{\lambda}_{1}\right)  \hat{\lambda}_{2}%
+\hat{\lambda}_{2}\Gamma_{1}^{2}\left(  \hat{\lambda}_{1}\right)  =0
\end{equation}%
\begin{equation}
P_{\nu}\wedge P_{\nu}+\hat{\lambda}_{2}^{2}-\hat{\lambda}_{2}+P_{\nu}%
\wedge\left(  P_{\nu}+P_{r-\nu}\right)  \hat{\lambda}_{2}+\hat{\lambda}%
_{2}P_{\nu}\wedge\left(  P_{\nu}+P_{r-\nu}\right)  =0
\end{equation}%
\begin{equation}
P_{\nu}\wedge P_{\nu}+\hat{\lambda}_{2}^{2}-\hat{\lambda}_{2}+\left(  P_{\nu
}\wedge P_{\nu}+P_{\nu}\wedge P_{r-\nu}\right)  \hat{\lambda}_{2}+\hat
{\lambda}_{2}\left(  P_{\nu}\wedge P_{\nu}+P_{\nu}\wedge P_{r-\nu}\right)  =0
\end{equation}%
\begin{equation}
P_{\nu}\wedge P_{\nu}\wedge\hat{\lambda}_{2}+\hat{\lambda}_{2}^{2}\wedge
\hat{\lambda}_{2}-\hat{\lambda}_{2}\wedge\hat{\lambda}_{2}+\left(  \left(
P_{\nu}\wedge P_{\nu}+P_{\nu}\wedge P_{r-\nu}\right)  \hat{\lambda}%
_{2}\right)  \wedge\hat{\lambda}_{2}+\left(  \hat{\lambda}_{2}\left(  P_{\nu
}\wedge P_{\nu}+P_{\nu}\wedge P_{r-\nu}\right)  \right)  \wedge\hat{\lambda
}_{2}=0
\end{equation}%
\begin{equation}
\hat{\lambda}_{2}^{2}\wedge\hat{\lambda}_{2}+\left(  \left(  P_{\nu}\wedge
P_{\nu}+P_{\nu}\wedge P_{r-\nu}\right)  \hat{\lambda}_{2}\right)  \wedge
\hat{\lambda}_{2}+\left(  \hat{\lambda}_{2}\left(  P_{\nu}\wedge P_{\nu
}+P_{\nu}\wedge P_{r-\nu}\right)  \right)  \wedge\hat{\lambda}_{2}=0
\end{equation}%
\begin{equation}
\left(  \hat{\lambda}_{2}^{2}+\left(  \left(  P_{\nu}\wedge P_{\nu}+P_{\nu
}\wedge P_{r-\nu}\right)  \hat{\lambda}_{2}\right)  +\left(  \hat{\lambda}%
_{2}\left(  P_{\nu}\wedge P_{\nu}+P_{\nu}\wedge P_{r-\nu}\right)  \right)
\right)  \wedge\hat{\lambda}_{2}=0
\end{equation}%
\begin{equation}
\left(  \hat{\lambda}_{2}^{2}+\left(  P_{r}\wedge P_{r}-P_{r-\nu}\wedge
P_{r}\right)  \hat{\lambda}_{2}+\hat{\lambda}_{2}\left(  P_{r}\wedge
P_{r}-P_{r-\nu}\wedge P_{r}\right)  \right)  \wedge\hat{\lambda}_{2}=0
\end{equation}%
\begin{equation}
\left(  \hat{\lambda}_{2}^{2}+\hat{\lambda}_{2}-\left(  P_{r-\nu}\wedge
P_{r}\right)  \hat{\lambda}_{2}+\hat{\lambda}_{2}-\left(  P_{r-\nu}\wedge
P_{r}\right)  \hat{\lambda}_{2}\right)  \wedge\hat{\lambda}_{2}=0
\end{equation}%
\begin{equation}
\left(  \overset{}{\hat{\lambda}_{2}^{2}-\left(  P_{r-\nu}\wedge P_{r}\right)
\hat{\lambda}_{2}-\left(  P_{r-\nu}\wedge P_{r}\right)  \hat{\lambda}_{2}%
}\right)  \wedge\hat{\lambda}_{2}=0
\end{equation}
Given a projector $\hat{\lambda}_{1}$ what solutions for $\hat{\lambda}_{2}$
are possible for $\hat{\lambda}_{1}\wedge\hat{\lambda}_{2}=0$, first we
determine the subspace $\left\{  \left\vert \psi_{k}\psi_{l}\right\rangle
\left\langle \psi_{m}\psi_{n}\right\vert \right\}  ,$ which is the kernal of
the linear map $\sum_{1\leq j\leq v}\left\vert \psi_{j}\right\rangle
\left\langle \psi_{j}\right\vert \wedge$ i.e.
\begin{equation}
\sum_{1\leq j\leq v}\left\vert \psi_{j}\right\rangle \left\langle \psi
_{j}\right\vert \wedge\left\vert \psi_{k}\psi_{l}\right\rangle \left\langle
\psi_{m}\psi_{n}\right\vert =0.
\end{equation}
This gives for $\nu=1$
\begin{equation}
\left\vert \psi_{1}\right\rangle \left\langle \psi_{1}\right\vert
\wedge\left\vert \psi_{k}\psi_{l}\right\rangle \left\langle \psi_{m}\psi
_{n}\right\vert =0\Rightarrow k=1\text{ or }l=1\text{ or }m=1\text{ or }n=1
\end{equation}
i.e. $\left\vert \psi_{k}\psi_{l}\right\rangle \left\langle \psi_{m}\psi
_{n}\right\vert $ must be a linear combination of
\begin{equation}
\left\{  \left\vert \psi_{1}\psi_{l}\right\rangle \left\langle \psi_{m}%
\psi_{n}\right\vert ,\left\vert \psi_{k}\psi_{l}\right\rangle \left\langle
\psi_{1}\psi_{n}\right\vert ,\left\vert \psi_{1}\psi_{l}\right\rangle
\left\langle \psi_{1}\psi_{n}\right\vert ;l,m,n\right\}  .
\end{equation}
For $\nu=2$ the operators $\left\{  \left\vert \psi_{k}\psi_{l}\right\rangle
\left\langle \psi_{m}\psi_{n}\right\vert \right\}  $ must be "orthogonal" to
both $\left\vert \psi_{1}\right\rangle \left\langle \psi_{1}\right\vert $ and
$\left\vert \psi_{2}\right\rangle \left\langle \psi_{2}\right\vert $ and must
be linear combinations of
\begin{equation}
\left\vert \psi_{1}\psi_{l}\right\rangle \left\langle \psi_{2}\psi
_{n}\right\vert ,\left\vert \psi_{k}\psi_{l}\right\rangle \left\langle
\psi_{1}\psi_{2}\right\vert ,\left\vert \psi_{2}\psi_{l}\right\rangle
\left\langle \psi_{1}\psi_{n}\right\vert ,\left\vert \psi_{1}\psi
_{2}\right\rangle \left\langle \psi_{m}\psi_{n}\right\vert .
\end{equation}
For $\nu=3$ the operators $\left\{  \left\vert \psi_{k}\psi_{l}\right\rangle
\left\langle \psi_{m}\psi_{n}\right\vert \right\}  $ must be "orthogonal" to
$\left\vert \psi_{1}\right\rangle \left\langle \psi_{1}\right\vert $,
$\left\vert \psi_{2}\right\rangle \left\langle \psi_{2}\right\vert $ and
$\left\vert \psi_{3}\right\rangle \left\langle \psi_{3}\right\vert $ thus they
must be linear combinations of
\begin{equation}
\left\vert \psi_{1}\psi_{3}\right\rangle \left\langle \psi_{2}\psi
_{n}\right\vert ,\left\vert \psi_{1}\psi_{l}\right\rangle \left\langle
\psi_{2}\psi_{3}\right\vert ,\left\vert \psi_{k}\psi_{3}\right\rangle
\left\langle \psi_{1}\psi_{2}\right\vert ,\left\vert \psi_{2}\psi
_{3}\right\rangle \left\langle \psi_{1}\psi_{n}\right\vert ,\left\vert
\psi_{2}\psi_{l}\right\rangle \left\langle \psi_{1}\psi_{3}\right\vert
,\left\vert \psi_{1}\psi_{2}\right\rangle \left\langle \psi_{3}\psi
_{n}\right\vert .
\end{equation}
For $\nu=4$ the operators $\left\{  \left\vert \psi_{k}\psi_{l}\right\rangle
\left\langle \psi_{m}\psi_{n}\right\vert \right\}  $ must be "orthogonal" to
$\left\vert \psi_{1}\right\rangle \left\langle \psi_{1}\right\vert $,
$\left\vert \psi_{2}\right\rangle \left\langle \psi_{2}\right\vert ,\left\vert
\psi_{3}\right\rangle \left\langle \psi_{3}\right\vert $ and $\left\vert
\psi_{4}\right\rangle \left\langle \psi_{4}\right\vert $ and thus must be
linear combinations of
\begin{equation}
\left\vert \psi_{1}\psi_{3}\right\rangle \left\langle \psi_{2}\psi
_{4}\right\vert ,\left\vert \psi_{1}\psi_{4}\right\rangle \left\langle
\psi_{2}\psi_{3}\right\vert ,\left\vert \psi_{3}\psi_{4}\right\rangle
\left\langle \psi_{1}\psi_{2}\right\vert ,\left\vert \psi_{2}\psi
_{3}\right\rangle \left\langle \psi_{1}\psi_{4}\right\vert ,\left\vert
\psi_{2}\psi_{4}\right\rangle \left\langle \psi_{1}\psi_{3}\right\vert
,\left\vert \psi_{1}\psi_{2}\right\rangle \left\langle \psi_{3}\psi
_{4}\right\vert .
\end{equation}
For $\nu\geq5$ the kernal of the linear map $\sum_{1\leq j\leq5}\left\vert
\psi_{j}\right\rangle \left\langle \psi_{j}\right\vert \wedge$ is empty.

$\lceil$The operator $\hat{\lambda}_{2}$ is self adjoint thus one needs to
consider basis elements
\begin{equation}
\left\{  \frac{1}{\sqrt{2}}\left(  \left\vert \psi_{k}\psi_{l}\right\rangle
\left\langle \psi_{m}\psi_{n}\right\vert +\left\vert \psi_{m}\psi
_{n}\right\rangle \left\langle \psi_{k}\psi_{l}\right\vert \right)  ,\frac
{i}{\sqrt{2}}\left(  \left\vert \psi_{k}\psi_{l}\right\rangle \left\langle
\psi_{m}\psi_{n}\right\vert -\left\vert \psi_{m}\psi_{n}\right\rangle
\left\langle \psi_{k}\psi_{l}\right\vert \right)  \right\}
\end{equation}
which in order to be orthogonal to $\left\vert \psi_{1}\right\rangle
\left\langle \psi_{1}\right\vert $ and $\left\vert \psi_{2}\right\rangle
\left\langle \psi_{2}\right\vert $ must be of the form
\begin{align}
&  \frac{1}{\sqrt{2}}\left(  \left\vert \psi_{1}\psi_{l}\right\rangle
\left\langle \psi_{2}\psi_{n}\right\vert +\left\vert \psi_{2}\psi
_{l}\right\rangle \left\langle \psi_{1}\psi_{n}\right\vert \right)  ,\frac
{1}{\sqrt{2}}\left(  \left\vert \psi_{k}\psi_{l}\right\rangle \left\langle
\psi_{1}\psi_{2}\right\vert +\left\vert \psi_{1}\psi_{2}\right\rangle
\left\langle \psi_{k}\psi_{l}\right\vert \right) \nonumber\\
&  \frac{i}{\sqrt{2}}\left(  \left\vert \psi_{1}\psi_{l}\right\rangle
\left\langle \psi_{2}\psi_{n}\right\vert -\left\vert \psi_{2}\psi
_{l}\right\rangle \left\langle \psi_{1}\psi_{n}\right\vert \right)  ,\frac
{i}{\sqrt{2}}\left(  \left\vert \psi_{k}\psi_{l}\right\rangle \left\langle
\psi_{1}\psi_{2}\right\vert -\left\vert \psi_{1}\psi_{2}\right\rangle
\left\langle \psi_{k}\psi_{l}\right\vert \right)  ,
\end{align}
while in order it be orthogonal to $\left\vert \psi_{1}\right\rangle
\left\langle \psi_{1}\right\vert $, $\left\vert \psi_{2}\right\rangle
\left\langle \psi_{2}\right\vert $ and $\left\vert \psi_{3}\right\rangle
\left\langle \psi_{3}\right\vert $ they must be of the form
\begin{align}
&  \frac{1}{\sqrt{2}}\left(  \left\vert \psi_{1}\psi_{3}\right\rangle
\left\langle \psi_{2}\psi_{n}\right\vert +\left\vert \psi_{2}\psi
_{n}\right\rangle \left\langle \psi_{1}\psi_{3}\right\vert \right)  ,\frac
{1}{\sqrt{2}}\left(  \left\vert \psi_{3}\psi_{l}\right\rangle \left\langle
\psi_{1}\psi_{2}\right\vert +\left\vert \psi_{1}\psi_{2}\right\rangle
\left\langle \psi_{3}\psi_{l}\right\vert \right) \nonumber\\
&  \frac{1}{\sqrt{2}}\left(  \left\vert \psi_{2}\psi_{3}\right\rangle
\left\langle \psi_{1}\psi_{n}\right\vert +\left\vert \psi_{1}\psi
_{n}\right\rangle \left\langle \psi_{2}\psi_{3}\right\vert \right)
,\nonumber\\
&  \frac{i}{\sqrt{2}}\left(  \left\vert \psi_{1}\psi_{3}\right\rangle
\left\langle \psi_{2}\psi_{n}\right\vert -\left\vert \psi_{2}\psi
_{n}\right\rangle \left\langle \psi_{1}\psi_{3}\right\vert \right)  ,\frac
{i}{\sqrt{2}}\left(  \left\vert \psi_{3}\psi_{l}\right\rangle \left\langle
\psi_{1}\psi_{2}\right\vert -\left\vert \psi_{1}\psi_{2}\right\rangle
\left\langle \psi_{3}\psi_{l}\right\vert \right) \nonumber\\
&  \frac{i}{\sqrt{2}}\left(  \left\vert \psi_{2}\psi_{3}\right\rangle
\left\langle \psi_{1}\psi_{n}\right\vert -\left\vert \psi_{1}\psi
_{n}\right\rangle \left\langle \psi_{2}\psi_{3}\right\vert \right)  ,
\end{align}
while in order be orthogonal to $\left\vert \psi_{1}\right\rangle \left\langle
\psi_{1}\right\vert $, $\left\vert \psi_{2}\right\rangle \left\langle \psi
_{2}\right\vert ,\left\vert \psi_{3}\right\rangle \left\langle \psi
_{3}\right\vert $ and $\left\vert \psi_{4}\right\rangle \left\langle \psi
_{4}\right\vert $ it must be of the form%
\begin{align}
&  \frac{1}{\sqrt{2}}\left(  \left\vert \psi_{1}\psi_{3}\right\rangle
\left\langle \psi_{2}\psi_{4}\right\vert +\left\vert \psi_{2}\psi
_{4}\right\rangle \left\langle \psi_{1}\psi_{3}\right\vert \right)  ,\frac
{1}{\sqrt{2}}\left(  \left\vert \psi_{1}\psi_{2}\right\rangle \left\langle
\psi_{3}\psi_{4}\right\vert +\left\vert \psi_{3}\psi_{4}\right\rangle
\left\langle \psi_{1}\psi_{2}\right\vert \right) \nonumber\\
&  \frac{1}{\sqrt{2}}\left(  \left\vert \psi_{1}\psi_{4}\right\rangle
\left\langle \psi_{2}\psi_{3}\right\vert +\left\vert \psi_{2}\psi
_{3}\right\rangle \left\langle \psi_{1}\psi_{4}\right\vert \right)
,\nonumber\\
&  \left\vert \psi_{1}\psi_{2}\right\rangle \left\langle \psi_{1}\psi
_{2}\right\vert ,\left\vert \psi_{1}\psi_{3}\right\rangle \left\langle
\psi_{1}\psi_{3}\right\vert ,\left\vert \psi_{1}\psi_{4}\right\rangle
\left\langle \psi_{1}\psi_{4}\right\vert ,\left\vert \psi_{2}\psi
_{3}\right\rangle \left\langle \psi_{2}\psi_{3}\right\vert ,\left\vert
\psi_{2}\psi_{4}\right\rangle \left\langle \psi_{2}\psi_{4}\right\vert
,\left\vert \psi_{3}\psi_{4}\right\rangle \left\langle \psi_{3}\psi
_{4}\right\vert ,\nonumber\\
&  \frac{i}{\sqrt{2}}\left(  \left\vert \psi_{1}\psi_{3}\right\rangle
\left\langle \psi_{2}\psi_{4}\right\vert -\left\vert \psi_{2}\psi
_{4}\right\rangle \left\langle \psi_{1}\psi_{3}\right\vert \right)  ,\frac
{i}{\sqrt{2}}\left(  \left\vert \psi_{3}\psi_{4}\right\rangle \left\langle
\psi_{1}\psi_{2}\right\vert -\left\vert \psi_{1}\psi_{2}\right\rangle
\left\langle \psi_{3}\psi_{4}\right\vert \right) \nonumber\\
&  \frac{i}{\sqrt{2}}\left(  \left\vert \psi_{2}\psi_{3}\right\rangle
\left\langle \psi_{1}\psi_{4}\right\vert -\left\vert \psi_{1}\psi
_{4}\right\rangle \left\langle \psi_{2}\psi_{3}\right\vert \right)  .
\end{align}
$\rfloor$

The condition
\begin{equation}
\hat{\lambda}_{2}\wedge\hat{\lambda}_{2}=0
\end{equation}
gives
\begin{subequations}
\begin{align}
\sum_{1\leq l\leq\binom{r}{2}}\beta_{l}\left\vert g_{l}\right\rangle
\left\langle g_{l}\right\vert \wedge\sum_{1\leq m\leq\binom{r}{2}}\beta
_{m}\left\vert g_{m}\right\rangle \left\langle g_{m}\right\vert  &
=2\sum_{1\leq l<m\leq\binom{r}{2}}\beta_{l}\beta_{m}\left\vert g_{l}\wedge
g_{m}\right\rangle \left\langle g_{l}\wedge g_{m}\right\vert +\sum_{1\leq
l\leq\binom{r}{2}}\beta_{l}^{2}\left\vert g_{l}^{2}\right\rangle \left\langle
g_{l}^{2}\right\vert =0\nonumber\\
&  \Updownarrow\nonumber\\
\left\vert g_{l}\wedge g_{m}\right\rangle \left\langle g_{l}\wedge
g_{m}\right\vert  &  =0;1\leq l<m\leq r,\beta_{l}\neq0\text{ and }\beta
_{m}\neq0\\
&  \Updownarrow\nonumber\\
\left\vert g_{l}\wedge g_{m}\right\rangle  &  =0;1\leq l<m\leq r,\beta_{l}%
\neq0\text{ and }\beta_{m}\neq0
\end{align}%
\end{subequations}
\begin{equation}
\left\vert g_{l}\wedge g_{m}\right\rangle =\frac{1}{2}\sum_{\substack{1\leq
j,k\leq r\\1\leq p,q\leq r}}g_{ljk}g_{mpq}\left\vert \varphi_{j}\wedge
\varphi_{k}\wedge\varphi_{p}\wedge\varphi_{q}\right\rangle =K\sum_{1\leq
j<k<p<q\leq r}\left(  g_{ljk}g_{mpq}\right)  \left\vert \varphi_{j}%
\wedge\varphi_{k}\wedge\varphi_{p}\wedge\varphi_{q}\right\rangle
\end{equation}%
\begin{equation}
Tr\left\{  \hat{\lambda}_{2}\wedge\hat{\lambda}_{2}\right\}  =\sum_{1\leq
P<Q\leq\binom{r}{2}}\alpha_{P}\alpha_{Q}=0
\end{equation}
gives
\begin{align}
&  \sum_{\substack{2<i<r\\2\leq j,k\leq r}}\left\{  \lambda_{1ijk}\left\vert
\psi_{1}\psi_{i}\right\rangle \left\langle \psi_{j}\psi_{k}\right\vert
+\lambda_{jk1i}\left\vert \psi_{j}\psi_{k}\right\rangle \left\langle \psi
_{1}\psi_{i}\right\vert \right\}  \wedge\sum_{\substack{2<i<r\\2\leq j,k\leq
r}}\left\{  \lambda_{1ijk}\left\vert \psi_{1}\psi_{i}\right\rangle
\left\langle \psi_{j}\psi_{k}\right\vert +\lambda_{jk1i}\left\vert \psi
_{j}\psi_{k}\right\rangle \left\langle \psi_{1}\psi_{i}\right\vert \right\}
\nonumber\\
&  +\sum_{\substack{2<i<r\\2\leq j,k\leq r}}\left\{  \lambda_{1ijk}\left\vert
\psi_{1}\psi_{i}\right\rangle \left\langle \psi_{j}\psi_{k}\right\vert
+\lambda_{jk1i}\left\vert \psi_{j}\psi_{k}\right\rangle \left\langle \psi
_{1}\psi_{i}\right\vert \right\}  \wedge\sum_{2\leq i,j\leq r}\lambda
_{1i1j}\left\vert \psi_{1}\psi_{i}\right\rangle \left\langle \psi_{1}\psi
_{j}\right\vert \nonumber\\
&  +\sum_{2\leq i,j\leq r}\lambda_{1i1j}\left\vert \psi_{1}\psi_{i}%
\right\rangle \left\langle \psi_{1}\psi_{j}\right\vert \wedge\sum
_{\substack{2<i<r\\2\leq j,k\leq r}}\left\{  \lambda_{1ijk}\left\vert \psi
_{1}\psi_{i}\right\rangle \left\langle \psi_{j}\psi_{k}\right\vert
+\lambda_{jk1i}\left\vert \psi_{j}\psi_{k}\right\rangle \left\langle \psi
_{1}\psi_{i}\right\vert \right\} \nonumber\\
&  +\sum_{2\leq i,j\leq r}\lambda_{1i1j}\left\vert \psi_{1}\psi_{i}%
\right\rangle \left\langle \psi_{1}\psi_{j}\right\vert \wedge\sum_{2\leq
i,j\leq r}\lambda_{1i1j}\left\vert \psi_{1}\psi_{i}\right\rangle \left\langle
\psi_{1}\psi_{j}\right\vert \left.  =\right.  0
\end{align}
leading to%
\begin{align}
&  \sum_{\substack{2<i<r\\2\leq j,k\leq r}}\left\{  \lambda_{1ijk}\left\vert
\psi_{1}\psi_{i}\right\rangle \left\langle \psi_{j}\psi_{k}\right\vert
+\lambda_{jk1i}\left\vert \psi_{j}\psi_{k}\right\rangle \left\langle \psi
_{1}\psi_{i}\right\vert \right\}  \wedge\sum_{\substack{2<i<r\\2\leq j,k\leq
r}}\left\{  \lambda_{1ijk}\left\vert \psi_{1}\psi_{i}\right\rangle
\left\langle \psi_{j}\psi_{k}\right\vert +\lambda_{jk1i}\left\vert \psi
_{j}\psi_{k}\right\rangle \left\langle \psi_{1}\psi_{i}\right\vert \right\}
\nonumber\\
&  =\sum_{\substack{2\leq i_{1}<\leq,2\leq i_{2}\leq r\\2\leq j_{1},k_{1}\leq
r,2\leq j_{2},k_{2}\leq r}}\left\{  \lambda_{1i_{1}j_{1}k_{1}}\lambda
_{j_{2}k_{2}1i_{2}}\left\vert \psi_{1}\psi_{i_{1}}\right\rangle \wedge
\left\vert \psi_{j_{2}}\psi_{k_{2}}\right\rangle \left\langle \psi_{j_{1}}%
\psi_{k_{1}}\right\vert \wedge\left\langle \psi_{1}\psi_{i_{2}}\right\vert
\right. \nonumber\\
&  \left.  +\lambda_{1i_{2}j_{2}k_{2}}\lambda_{j_{1}k_{1}1i_{1}}\left\vert
\psi_{1}\psi_{i_{2}}\right\rangle \wedge\left\vert \psi_{j_{1}}\psi_{k_{1}%
}\right\rangle \left\langle \psi_{j_{2}}\psi_{k_{2}}\right\vert \wedge
\left\langle \psi_{1}\psi_{i_{1}}\right\vert \right\} \nonumber\\
&  =\sum_{\substack{2\leq i_{1}<\leq,2\leq i_{2}\leq r\\2\leq j_{1},k_{1}\leq
r,2\leq j_{2},k_{2}\leq r}}\left\{  \lambda_{1i_{1}j_{1}k_{1}}\lambda
_{j_{2}k_{2}1i_{2}}\left\vert \left(  \psi_{1}\psi_{i_{1}}\right)
\wedge\left(  \psi_{j_{2}}\psi_{k_{2}}\right)  \right\rangle \left\langle
\left(  \psi_{j_{1}}\psi_{k_{1}}\right)  \wedge\left(  \psi_{1}\psi_{i_{2}%
}\right)  \right\vert \right. \nonumber\\
&  \left.  +\lambda_{1i_{2}j_{2}k_{2}}\lambda_{j_{1}k_{1}1i_{1}}\left\vert
\left(  \psi_{1}\psi_{i_{2}}\right)  \wedge\left(  \psi_{j_{1}}\psi_{k_{1}%
}\right)  \right\rangle \left\langle \left(  \psi_{j_{2}}\psi_{k_{2}}\right)
\wedge\left(  \psi_{1}\psi_{i_{1}}\right)  \right\vert \right\} \\
&  =2\sum_{\substack{2\leq i_{1}<\leq,2\leq i_{2}\leq r\\2\leq j_{1},k_{1}\leq
r,2\leq j_{2},k_{2}\leq r}}\left\{  \lambda_{1i_{1}j_{1}k_{1}}\lambda
_{j_{2}k_{2}1i_{2}}\left\vert \left(  \psi_{1}\psi_{i_{1}}\right)
\wedge\left(  \psi_{j_{2}}\psi_{k_{2}}\right)  \right\rangle \left\langle
\left(  \psi_{j_{1}}\psi_{k_{1}}\right)  \wedge\left(  \psi_{1}\psi_{i_{2}%
}\right)  \right\vert \right\}  \left.  =\right.  0\\
&  \Rightarrow\sum_{\substack{2\leq i_{1}<\leq,2\leq i_{2}\leq r\\2\leq
j_{1},k_{1}\leq r,2\leq j_{2},k_{2}\leq r}}\left\{  \lambda_{1i_{1}j_{1}k_{1}%
}\lambda_{j_{2}k_{2}1i_{2}}\left\vert \psi_{1}\psi_{i_{1}}\psi_{j_{2}}%
\psi_{k_{2}}\right\rangle \left\langle \psi_{1}\psi_{i_{2}}\psi_{j_{1}}%
\psi_{k_{1}}\right\vert \right\}  \left.  =\right.  0
\end{align}



\end{document}